
\documentclass[final]{siamart0216}
%\documentclass[review]{siamart0216}
\usepackage{amsfonts}
\usepackage{todonotes}
\usepackage{amsfonts,amsmath,amssymb,stackengine}
\usepackage{mathrsfs,mathtools,stmaryrd,wasysym}
\usepackage{enumerate}
\usepackage{xspace,mydef} 
\usepackage{esint}
\usepackage{graphicx,psfrag}
\usepackage{hyperref}
\usepackage{pgf,tikz,pgfplots}
\usepackage{pstricks-add}
\usepackage{enumitem}
\usepackage{yfonts}
\usepackage{rotating}

\usetikzlibrary{arrows}
\DeclareMathAlphabet{\mathpzc}{OT1}{pzc}{m}{it}

%% ============== TO BE COMMENTED ===========
% \usepackage[notcite,notref]{showkeys} % To see crossreferences.
% \newcommand{\alert}[1]{\noindent{\textcolor{red}{#1}}}
% \newcommand{\EO}[1]{\noindent{\textcolor{blue}{#1}}}
%  \newcommand{\GC}[1]{\noindent{\textcolor{magenta}{#1}}}%% ===========================================

\DeclareMathOperator{\supp}{supp}
\newsiamremark{remark}{Remark}

\newcommand{\norm}[1]{{\left\vert\kern-0.25ex\left\vert\kern-0.25ex\left\vert #1 
\right\vert\kern-0.25ex\right\vert\kern-0.25ex\right\vert}}


\newcommand{\TheTitle}{
Bilinear optimal control for the Stokes--Brinkman equations: a priori and a posteriori error analyses}
\newcommand{\ShortTitle}{The bilinear optimal control of the Stokes--Brinkman equations}
\newcommand{\TheAuthors}{A. Allendes, G. Campa\~na and E.~Ot\'arola}

\headers{\ShortTitle}{\TheAuthors}

\title{{\TheTitle}\thanks{AA is partially supported by ANID through FONDECYT project 1210729. GC is partially supported by ANID-Subdirecci\'on de Capital Humano/Doctorado Nacional/2020--2120092. 
EO is partially supported by ANID through FONDECYT project 1220156.}}

\author{Alejandro Allendes\thanks{Departamento de Matem\'atica, Universidad T\'ecnica Federico Santa Mar\'ia, Valpara\'iso, Chile.
(\email{alejandro.allendes@usm.cl}, \url{http://aallendes.mat.utfsm.cl/}).}
\and 
Gilberto Campa\~na\thanks{Departamento de Ciencias, Universidad T\'ecnica Federico Santa Mar\'ia, Valpara\'iso, Chile.
(\email{gilberto.campana@usm.cl}.}
\and
Enrique Ot\'arola\thanks{Departamento de Matem\'atica, Universidad T\'ecnica Federico Santa Mar\'ia, Valpara\'iso, Chile.
(\email{enrique.otarola@usm.cl}, \url{http://eotarola.mat.utfsm.cl/}).}   
}

\ifpdf
\hypersetup{
  pdftitle={\TheTitle},
  pdfauthor={\TheAuthors}
}
\fi

\date{Draft version of \today.}

\begin{document}

\maketitle


\begin{abstract}
We analyze a bilinear optimal control problem for the Stokes--Brinkman equations: the control variable enters the state equations as a coefficient. In two- and three-dimensional Lipschitz domains, we perform a complete continuous analysis that includes the existence of solutions and first- and second-order optimality conditions. We also develop two finite element methods that differ fundamentally in whether the admissible control set is discretized or not. For each of the proposed methods, we perform a convergence analysis and derive a priori error estimates; the latter under the assumption that the domain is convex. Finally, assuming that the domain is Lipschitz, we develop an a posteriori error estimator for each discretization scheme and obtain a global reliability bound.
\end{abstract}

\begin{keywords}
bilinear optimal control, Stokes--Brinkman equations, first- and second-order optimality conditions, finite elements, convergence, a priori error bounds, a posteriori error bounds.
\end{keywords}

% REQUIRED
\begin{AMS}
35Q35,  % PDEs in connection with fluid mechanics
49K20,  % Optimality conditions for problems involving partial differential equations
49M25,  % Discrete approximations in optimal control
65K10,  % Numerical optimization and variational techniques
65N15,  % Error bounds for boundary value problems involving PDEs
65N30,  % Finite element, Rayleigh-Ritz and Galerkin methods for boundary value problems involving PDEs
65N50,  % Mesh generation, refinement, and adaptive methods for boundary value problems involving PDEs
76D07.  % Stokes and related (Oseen, etc.) flows
\end{AMS}

%%%%%%%%%%%%%%%%%%%%%%%%%%%%%%%%%%%%%%%%%%%%%%%%%%%%%%%%%%%%%%%%
%%%%%%%%%%%%%%%%%%%%%%%%%%%%%%%%%%%%%%%%%%%%%%%%%%%%%%%%%%%%%%%%
%%%%%%%%%%%%%%%%%%%%%%%%%%%%%%%%%%%%%%%%%%%%%%%%%%%%%%%%%%%%%%%%
%%%%%%%%%%%%%%%%%%%%%%%%%%%%%%%%%%%%%%%%%%%%%%%%%%%%%%%%%%%%%%%%
%%%%%%%%%%%%%%%%%%%%%%%%%%%%%%%%%%%%%%%%%%%%%%%%%%%%%%%%%%%%%%%%
%%%%%%%%%%%%%%%%%%%%%%%%%%%%%%%%%%%%%%%%%%%%%%%%%%%%%%%%%%%%%%%%
%%%%%%%%%%%%%%%%%%%%%%%%%%%%%%%%%%%%%%%%%%%%%%%%%%%%%%%%%%%%%%%%
%%%%%%%%%%%%%%%%%%%%%%%%%%%%%%%%%%%%%%%%%%%%%%%%%%%%%%%%%%%%%%%%

\section{Introduction}
\label{sec:intro}
In this paper we deal with the \emph{analysis} and \emph{discretization} of a bilinear optimal control problem for the Stokes--Brinkman equations: the control variable enters the state equations as a coefficient and not as a source term. To be precise, we let $d\in\{2,3\}$ and $\Omega \subset \mathbb{R}^d$ be an open and bounded domain with Lipschitz boundary $\partial\Omega$. Given a regularization parameter $\alpha>0$ and a desired velocity field $\mathbf{y}_\Omega\in \mathbf{L}^2(\Omega)$, we introduce the cost functional
\begin{equation}\label{eq:functional}
J(\mathbf{y},u) :=
\dfrac{1}{2}\|\mathbf{y}-\mathbf{y}_\Omega\|_{\mathbf{L}^2(\Omega)}^2
+
\dfrac{\alpha}{2}\|u\|^2_{L^2(\Omega)}.
\end{equation}
Let $\mathbf{f} \in \mathbf{H}^{-1}(\Omega)$ be the body force acting on the fluid. The bilinear optimal control problem reads as follows: Find $\min J(\mathbf{y},u)$ subject to the \emph{Stokes--Brinkman} equations
\begin{align}\label{eq:state}
-\Delta \mathbf{y} + u\mathbf{y}+\nabla \mathsf{p} =  \mathbf{f}
\text{ in } \Omega,
\quad
\text{div}~\mathbf{y}  =  0
\text{ in } \Omega,
\quad
\mathbf{y}  =  0
\text{ on } \partial\Omega,
\end{align}
where $\mathbf{y}$ is the velocity field, $\mathsf{p}$ is the pressure of the fluid, and $u$ corresponds to the control variable, which is chosen so that
\begin{equation}\label{eq:set_adm}
u\in \mathbb{U}_{ad},
\quad
\mathbb{U}_{ad}:=
\{
v \in L^2(\Omega): 0<\mathsf{a}\leq v\leq \mathsf{b} \text{ a.e. in } \Omega
\}.
\end{equation}
The control bounds $\mathsf{a}$ and $\mathsf{b}$ both belong to $\mathbb{R}$ and are chosen so that $\mathsf{a} < \mathsf{b}$.

When the state equations are replaced by the problem $-\Delta u + q u = f$ in $\Omega$ and $u = 0$ on $\partial \Omega$, there are several works that provide error estimates for finite element discretizations of \eqref{eq:functional}--\eqref{eq:set_adm}. To the best of our knowledge, \cite{MR2536007} is the first paper that provides an analysis. In this paper, the authors investigate discretization methods and derive error bounds \cite[Corollaries 5.6 and 5.10]{MR2536007}. These results were later extended to mixed and stabilized methods in \cite{MR3103238} and \cite{MR3693332}, respectively. Regarding the analysis of \emph{a posteriori} error estimates, we refer the reader to \cite{MR2680928} and \cite{MR4450052}. We would also like to refer to the paper \cite{MR4081184}, in which the author studies an optimal control problem with a bilinear term in the boundary condition of the state equation. Finally, we mention the recent advances in the semilinear scenario presented in \cite{MR4679130} and \cite{casas2024error}.

As far as we know, this paper is the \emph{first} to study optimality conditions and finite element schemes for the control problem \eqref{eq:state}--\eqref{eq:set_adm}. The analysis is not trivial and involves a number of difficulties. To overcome them, we had to provide several results. In the following, we list what we consider to be the main contributions of our work:

$\bullet$ \emph{Existence results and optimality conditions:} We prove that the optimal control problem admits at least one optimal solution and derive first- and necessary and sufficient second-order optimality conditions with a minimal gap.

$\bullet$ \emph{Finite element schemes:} We propose two different schemes: a semidiscrete method in which the control set is not discretized, and a fully discrete method in which such a set is discretized with piecewise constant functions. For each scheme, we prove convergence results and a priori error bounds: $\mathcal{O}(h)$ for the fully discrete scheme and $\mathcal{O}(h^2)$ for the semidiscrete one. We also study a posteriori error bounds.

%%%%%%%%%%%%%%%%%%%%%%%%%%%%%%%%%%%%%%%%%%%%%%%%%%%%%%%%%%%%%%%%
%%%%%%%%%%%%%%%%%%%%%%%%%%%%%%%%%%%%%%%%%%%%%%%%%%%%%%%%%%%%%%%%
%%%%%%%%%%%%%%%%%%%%%%%%%%%%%%%%%%%%%%%%%%%%%%%%%%%%%%%%%%%%%%%%
%%%%%%%%%%%%%%%%%%%%%%%%%%%%%%%%%%%%%%%%%%%%%%%%%%%%%%%%%%%%%%%%
%%%%%%%%%%%%%%%%%%%%%%%%%%%%%%%%%%%%%%%%%%%%%%%%%%%%%%%%%%%%%%%%
%%%%%%%%%%%%%%%%%%%%%%%%%%%%%%%%%%%%%%%%%%%%%%%%%%%%%%%%%%%%%%%%
%%%%%%%%%%%%%%%%%%%%%%%%%%%%%%%%%%%%%%%%%%%%%%%%%%%%%%%%%%%%%%%%
%%%%%%%%%%%%%%%%%%%%%%%%%%%%%%%%%%%%%%%%%%%%%%%%%%%%%%%%%%%%%%%%

\section{Notation and preliminary remarks}
\label{sec:notation}
Let us establish the notation and introduce the framework within which we will work on our manuscript.

%%%%%%%%%%%%%%%%%%%%%%%%%%%%%%%%%%%%%%%%%%%%%%%%%%%%%%%%%%%%%%%%
%%%%%%%%%%%%%%%%%%%%%%%%%%%%%%%%%%%%%%%%%%%%%%%%%%%%%%%%%%%%%%%%
%%%%%%%%%%%%%%%%%%%%%%%%%%%%%%%%%%%%%%%%%%%%%%%%%%%%%%%%%%%%%%%%
%%%%%%%%%%%%%%%%%%%%%%%%%%%%%%%%%%%%%%%%%%%%%%%%%%%%%%%%%%%%%%%%

\subsection{Notation}
Throughout the text, we use the classical notation for Lebesgue and Sobolev spaces and denote the vector-valued counterparts of these spaces with bold letters. In particular, we set
$
\mathbf{V}(\Omega):=\{\mathbf{v}\in\mathbf{H}_0^1(\Omega):\text{div }\mathbf{v}=0\}.
$

Let $\mathcal{X}$ and $\mathcal{Y}$ be Banach function spaces. We write $\mathcal{X} \hookrightarrow \mathcal{Y}$ to indicate that $\mathcal{X}$ is continuously embedded in $\mathcal{Y}$. We denote by $\mathcal{W}'$ and $\|\cdot\|_{\mathcal{W}}$ the dual space and the norm of $\mathcal{W}$, respectively. Given $\mathtt{p}\in(1, \infty)$, we denote by $\mathtt{p}\prime$ its H\"older conjugate, which is such that $1/\mathtt{p} + 1/\mathtt{p}\prime = 1$. The relation $a\lesssim b$ states that $a\leq Cb$ with a positive constant $C$ that does not depend on $a$, $b$, or the discretization parameters. The value of $C$ might change at each occurrence.

\section{The Stokes--Brinkman problem}
We present some of the main results related to the well-posedness and regularity of the Stokes--Brinkman equations \eqref{eq:state}. Let $\boldsymbol{\mathfrak{f}}$ be in $\mathbf{H}^{-1}(\Omega)$, and let $\mathfrak{u}$ be an arbitrary function in  $\mathbb{U}_{ad}$. Under these assumptions, we introduce the following weak formulation: Find $(\boldsymbol{\mathsf{y}},\mathfrak{p})\in \mathbf{H}_0^1(\Omega)\times L_0^2(\Omega)$ such that
\begin{equation}
\label{eq:brinkman_problem}
(\nabla \boldsymbol{\mathsf{y}},\nabla \mathbf{v})_{\mathbf{L}^2(\Omega)}
+
( \mathfrak{u}\boldsymbol{\mathsf{y}}, \mathbf{v})_{\mathbf{L}^2(\Omega)}
-
(\mathfrak{p},\text{div }\mathbf{v})_{L^2(\Omega)}
=
\langle \boldsymbol{\mathfrak{f}},\mathbf{v} \rangle,
\quad
(\mathsf{q},\text{div }\boldsymbol{\mathsf{y}})_{L^2(\Omega)}
=
0
\end{equation}
for all $(\mathbf{v},\mathsf{q}) \in \mathbf{H}_0^1(\Omega) \times L_0^2(\Omega)$. As a first ingredient to analyze \eqref{eq:brinkman_problem}, we recall that the divergence operator is surjective from $\mathbf{H}_0^1(\Omega)$ to $L_0^2(\Omega)$. This implies the existence of a constant $\beta>0$ such that \cite[Chapter I, Section 5.1]{MR851383}, \cite[Corollary B.71]{Guermond-Ern}
\begin{equation}\label{eq:infsup}
\sup_{\mathbf{v}\in \mathbf{H}_0^1(\Omega)}\dfrac{(\mathsf{q},\text{div }\mathbf{v})_{L^2(\Omega)}}{\|\nabla \mathbf{v}\|_{\mathbf{L}^2(\Omega)}}\geq \beta \|\mathsf{q}\|_{L^2(\Omega)}\quad\forall \mathsf{q}\in L_0^2(\Omega).
\end{equation}
As a second ingredient, we note that $(\mathbf{w},\mathbf{v}) \mapsto (\nabla \mathbf{w},\nabla \mathbf{v})_{\mathbf{L}^2(\Omega)}
+
( \mathfrak{u}\mathbf{w}, \mathbf{v})_{\mathbf{L}^2(\Omega)}$ defines a bilinear, continuous, and coercive form in $\mathbf{H}_0^1(\Omega) \times \mathbf{H}_0^1(\Omega)$. We can therefore apply the inf-sup theory for saddle point problems to obtain the well-posedness of \eqref{eq:brinkman_problem} and
\begin{align}\label{eq:well_posedness_brinkman}
	\|\nabla \boldsymbol{\mathsf{y}}\|_{\mathbf{L}^2(\Omega)}+\|\mathfrak{p}\|_{L^2(\Omega)}\lesssim \| \boldsymbol{\mathfrak{f}}\|_{\mathbf{H}^{-1}(\Omega)}.
\end{align}

The next result shows that it is possible to obtain an $\mathbf{L}^{\infty}(\Omega)$-regularity result for the velocity component if we assume better regularity properties for the datum $\boldsymbol{\mathfrak{f}}$.

\begin{theorem}[$\mathbf{L}^{\infty}(\Omega)$-regularity]\label{theorem_regularity_Linfty}
If $\Omega$ is Lipschitz, $\boldsymbol{\mathfrak{f}}\in\mathbf{L}^2(\Omega)$, and $\mathfrak{u}\in \mathbb{U}_{ad}$, then there exists $\kappa$ such that the solution $(\boldsymbol{\mathsf{y}},\mathfrak{p})$ of \eqref{eq:brinkman_problem} belongs to $\mathbf{H}_0^{1}(\Omega) \cap \mathbf{W}^{1,\kappa}(\Omega) \times L_0^2(\Omega) \cap L^{\kappa}(\Omega)$, where $\kappa > 4$ if $d=2$ and $\kappa >3$ if $d=3$. Furthermore, we have
\begin{equation}\label{eq:estimate_extra}
\| \nabla \boldsymbol{\mathsf{y}}\|_{\mathbf{L}^{\kappa}(\Omega)} +\|\mathfrak{p}\|_{L^{\kappa}(\Omega)}
\lesssim \| \boldsymbol{\mathfrak{f}}\|_{\mathbf{L}^2(\Omega)}.
\end{equation}
In particular, the velocity field $\boldsymbol{\mathsf{y}}$ belongs to $\mathbf{C}(\bar \Omega)$.
\end{theorem}
\begin{proof}
We begin the proof by writing the momentum equation of \eqref{eq:brinkman_problem} as $-\Delta \boldsymbol{\mathsf{y}} +\nabla \mathfrak{p} = \boldsymbol{\mathfrak{f}} - \mathfrak{u} \boldsymbol{\mathsf{y}}$ in $\Omega$. Note that $\boldsymbol{\mathfrak{f}} - \mathfrak{u} \boldsymbol{\mathsf{y}} \in \mathbf{L}^2(\Omega)$. Based on the basic Sobolev embedding in \cite[Theorem 4.12, Part I, Case C]{MR2424078}, we deduce the existence of $\kappa > 4$ when $d=2$ and $\kappa >3$ when $d=3$ such that $\boldsymbol{\mathfrak{f}} - \mathfrak{u} \boldsymbol{\mathsf{y}} \in \mathbf{W}^{-1,\kappa}(\Omega)$. The desired regularity for $(\boldsymbol{\mathsf{y}},\mathfrak{p})$ thus results from a direct application of \cite[(1.52)]{MR2987056}. Note that $\kappa$ could be further restricted. Finally, the fact that $\boldsymbol{\mathsf{y}} \in \mathbf{C}(\bar \Omega)$ follows from \cite[Theorem 4.12, Part II]{MR2424078}.
\end{proof}

We also present the following regularity result for the pair $(\boldsymbol{\mathsf{y}},\mathfrak{p})$ that solves \eqref{eq:brinkman_problem}.

\begin{theorem}[regularity]
\label{eq:theorem_regularity}
Let $\Omega\subset \mathbb{R}^d$ be a convex polytope, $\boldsymbol{\mathfrak{f}}\in\mathbf{L}^2(\Omega)$, and $\mathfrak{u}\in \mathbb{U}_{ad}$. If $(\boldsymbol{\mathsf{y}},\mathfrak{p})$ solves \eqref{eq:brinkman_problem}, then $\boldsymbol{\mathsf{y}}\in\mathbf{H}^2(\Omega)\cap\mathbf{H}_0^1(\Omega)$, $\mathfrak{p}\in H^1(\Omega)\cap L_0^2(\Omega)$, and
\begin{align}\label{eq:regularity}
\|\boldsymbol{\mathsf{y}}\|_{\mathbf{H}^2(\Omega)}+\|\mathfrak{p}\|_{H^1(\Omega)}\lesssim \|\boldsymbol{\mathfrak{f}}\|_{\mathbf{L}^2(\Omega)},
\end{align}
with a hidden constant that is independent of $(\boldsymbol{\mathsf{y}},\mathfrak{p})$ and $\mathfrak{u}$.
\end{theorem}
\begin{proof}
	Since $\boldsymbol{\mathfrak{f}}-\mathfrak{u}\mathbf{y}\in \mathbf{L}^2(\Omega)$, the proof follows immediately from an application of \cite{MR0977489,MR1301452} for $d=3$ and \cite[Theorem 2]{MR0404849} for $d=2$.
\end{proof}
%%%%%%%%%%%%%%%%%%%%%%%%%%%%%%%%%%%%%%%%%%%%%%%%%%%%%%%%%%%%%%%%
%%%%%%%%%%%%%%%%%%%%%%%%%%%%%%%%%%%%%%%%%%%%%%%%%%%%%%%%%%%%%%%%
%%%%%%%%%%%%%%%%%%%%%%%%%%%%%%%%%%%%%%%%%%%%%%%%%%%%%%%%%%%%%%%%

\section{The optimal control problem}
In this section, we introduce a weak formulation for
the optimal control problem \eqref{eq:min_functional}--\eqref{eq:brinkamn_problem_state} and analyze an existence result and first- and second-order optimality conditions.

\subsection{Weak formulation}\label{section:weak_formulation}
Given $\mathbf{f}\in\mathbf{H}^{-1}(\Omega)$, find
\begin{equation}
\label{eq:min_functional}
\min \{J(\mathbf{y},u):(\mathbf{y},u)\in\mathbf{H}_0^1(\Omega)\times \mathbb{U}_{ad}\}
\end{equation}
subject to the \emph{state equations} 
\begin{equation}\label{eq:brinkamn_problem_state}
(\nabla \mathbf{y},\nabla \mathbf{v})_{\mathbf{L}^2(\Omega)}+( u\mathbf{y}, \mathbf{v})_{\mathbf{L}^2(\Omega)}-(\mathsf{p},\text{div }\mathbf{v})_{L^2(\Omega)}
=
\langle \mathbf{f},\mathbf{v} \rangle,
\quad
(\mathsf{q},\text{div }\mathbf{y})_{L^2(\Omega)}
=
0
\end{equation}
for all $\mathbf{v}\in \mathbf{H}_0^1(\Omega)$ and $\mathsf{q}\in L_0^2(\Omega)$, respectively. Recall that $J$ is the cost functional defined in \eqref{eq:functional}. We note that due to de Rham's Theorem, the state equations \eqref{eq:brinkamn_problem_state} are equivalent to the following \emph{constrained formulation} \cite[Section 4.1.3]{Guermond-Ern}:
\begin{equation}\label{eq:brinkamn_problem_reduced}
\mathbf{y}\in \mathbf{V}(\Omega):
\quad
(\nabla \mathbf{y},\nabla \mathbf{v})_{\mathbf{L}^2(\Omega)}
+
( u\mathbf{y}, \mathbf{v})_{\mathbf{L}^2(\Omega)} 
=  
\langle \mathbf{f},\mathbf{v} \rangle
\quad 
\forall\mathbf{v}\in \mathbf{V}(\Omega).
\end{equation}


\subsection{Existence of solutions}\label{section:existence}
The existence of an optimal solution is as follows.
 
\begin{theorem}[existence of an optimal solution]\label{theorem_existence_optimalcontrol} The optimal control problem \eqref{eq:min_functional}--\eqref{eq:brinkamn_problem_state} admits at least one global solution $(\bar{\mathbf{y}},\bar{\mathsf{p}},\bar{u})\in \mathbf{H}_0^1(\Omega)\times L_0^2(\Omega)\times \mathbb{U}_{ad}$.
\end{theorem}
\begin{proof}
The proof follows from the direct method of calculus of variations. Let $\{(\mathbf{y}_{k},\mathsf{p}_k,u_{k})\}_{k\in\mathbb{N}}$ be a minimizing sequence, i.e., for $k\in \mathbb{N}$, $(\mathbf{y}_{k},\mathsf{p}_k) \in \mathbf{H}_0^1(\Omega)\times L_0^2(\Omega)$ solves  \eqref{eq:brinkamn_problem_state} where $u$ is replaced by $u_{k}$ and $\{ (\mathbf{y}_{k}, \mathsf{p}_k,u_k) \}_{k \in \mathbb{N}}$ is such that $J(\mathbf{y}_{k},u_{k}) \rightarrow \mathfrak{i}:=\inf\{ J(\mathbf{y}, u):(\mathbf{y}, \mathsf{p}, u)\in \mathbf{H}_0^1(\Omega)\times L_0^2(\Omega) \times \mathbb{U}_{ad}\}$ as $k\uparrow \infty$. Since $\mathbb{U}_{ad}$ is weakly sequentially compact in $L^2(\Omega)$, there exists a nonrelabeled subsequence $\{u_{k}\}_{k\in\mathbb{N}} \subset \mathbb{U}_{ad}$ such that $u_{k}\rightharpoonup \bar{u}$ in $L^2(\Omega)$ as $ k \uparrow \infty$ and $\bar{u} \in \mathbb{U}_{ad}$. On the other hand, since $u_k\in\mathbb{U}_{ad}$ for every $k\in\mathbb{N}$, the bound \eqref{eq:well_posedness_brinkman} shows that $\{(\mathbf{y}_k,\mathsf{p}_k)\}_{k\in\mathbb{N}}$ is uniformly bounded in $\mathbf{H}_0^1(\Omega)\times L_0^2(\Omega)$. We can therefore conclude the existence of a nonrelabeled subsequence $\{(\mathbf{y}_k,\mathsf{p}_k)\}_{k\in\mathbb{N}}$ such that $(\mathbf{y}_k,\mathsf{p}_k)\rightharpoonup (\bar{\mathbf{y}},\bar{\mathsf{p}})$ in $\mathbf{H}_0^1(\Omega)\times L_0^2(\Omega)$ as $k\uparrow\infty$; $(\bar{\mathbf{y}},\bar{\mathsf{p}})$ is the natural candidate for an optimal state. In the following, we prove that $(\bar{\mathbf{y}},\bar{\mathsf{p}})$ solves problem \eqref{eq:brinkamn_problem_state}, where $u$ is replaced by $\bar{u}$, and that $(\bar{\mathbf{y}},\bar{\mathsf{p}},\bar{u})$ is optimal.

With the weak convergence $(\mathbf{y}_k,\mathsf{p}_k)\rightharpoonup (\bar{\mathbf{y}},\bar{\mathsf{p}})$ in $\mathbf{H}_0^1(\Omega)\times L_0^2(\Omega)$ as $k\uparrow\infty$ at hand, we immediately obtain that for every $\mathbf{v}\in\mathbf{H}_0^1(\Omega)$ and $\mathsf{q}\in L_0^2(\Omega)$,
\begin{align*}
|(\nabla (\bar{\mathbf{y}}-\mathbf{y}_k),\nabla \mathbf{v})_{\mathbf{L}^2(\Omega)}|\to 0,~ |(\bar{\mathsf{p}}-\mathsf{p}_k,\text{div }\mathbf{v})_{L^2(\Omega)}|\to 0,~ |(\mathsf{q},\text{div}(\bar{\mathbf{y}}-\mathbf{y}_k))_{L^2(\Omega)}|\to 0
\end{align*}
as $k\uparrow \infty$. To analyze the convergence of the bilinear term, we use $u_k \rightharpoonup\bar{u}$ in $L^2 (\Omega)$ as $k\uparrow\infty$, the fact that $\bar{\mathbf{y}}\cdot\mathbf{v}\in L^2(\Omega)$, $\mathbf{y}_k\rightharpoonup \bar{\mathbf{y}}$ in $\mathbf{H}_0^1(\Omega)$ as $k\uparrow\infty$, and the compact embedding $\mathbf{H}_0^1(\Omega) \hookrightarrow\mathbf{L}^4(\Omega)$ \cite[Theorem 6.3, Part I]{MR2424078} to conclude that
\begin{multline*}
|(\bar{u}\bar{\mathbf{y}},\mathbf{v})_{\mathbf{L}^2(\Omega)} - (u_k\mathbf{y}_k,\mathbf{v})_{\mathbf{L}^2(\Omega)}|
\leq 
|\left((\bar{u} - u_k)\bar{\mathbf{y}},\mathbf{v}\right)_{\mathbf{L}^2(\Omega)}|
\\
+ 
\| u_k \|_{L^2(\Omega)}
\| \bar{\mathbf{y}} - \mathbf{y}_k \|_{\mathbf{L}^4(\Omega)}
\| \mathbf{v} \|_{\mathbf{L}^4(\Omega)}
\to 0, 
\qquad k\uparrow\infty,
\end{multline*}
where we have also used that $\| u_k \|_{L^2(\Omega)} \leq \mathcal{M}$ for every $k \in \mathbb{N}$. We have thus proved that $(\bar{\mathbf{y}},\bar{\mathsf{p}})$ solves problem \eqref{eq:brinkamn_problem_state}, where $u$ is replaced by $\bar{u}$. Finally, the optimality of $(\bar{\mathbf{y}},\bar{\mathsf{p}},\bar{u})$ follows from the strong convergence $\mathbf{y}_k\to\bar{\mathbf{y}}$ in $\mathbf{L}^2(\Omega)$ as $k\uparrow\infty$ and the fact that the square of $\|\cdot\|_{L^2(\Omega)}$ es weakly lower semicontinuous in $L^2(\Omega)$.
\end{proof}

\subsection{First-order optimality conditions}\label{section:firstconditions}
In this section, we obtain first-order optimality conditions for the optimal control problem \eqref{eq:min_functional}--\eqref{eq:brinkamn_problem_state}.
In the absence of convexity, we proceed as usual in the framework of \emph{local solutions} in $L^2(\Omega)$ \cite[page 207]{MR2583281}: A control $\bar{u}\in\mathbb{U}_{ad}$ is said to be \emph{locally optimal} in $L^2(\Omega)$ for \eqref{eq:min_functional}--\eqref{eq:brinkamn_problem_state} if there exists $\delta>0$ such that $J(\bar{\mathbf{y}},\bar{u})\leq J(\mathbf{y},u)$ for all $u\in \mathbb{U}_{ad}$ such that $\|u-\bar{u}\|_{L^2(\Omega)}\leq \delta$. Here, $\bar{\mathbf{y}}$ and $\mathbf{y}$ denote the velocity fields associated with $\bar{u}$ and $u$, respectively.

We begin our analysis by introducing $\mathcal{U}:=\{u\in L^\infty(\Omega):\exists \mathsf{c}>0\text{ such that }u(x)>\mathsf{c}>0\text{ for a.e. } x\in\Omega\}$ \cite[page 783]{MR2536007}. We note that $\mathbb{U}_{ad}\subset\mathcal{U}$. We also introduce the control-to-state map $\mathcal{S}:\mathcal{U}\to \mathbf{H}_0^1(\Omega)\times L_0^2(\Omega)$, which assigns the unique velocity-pressure pair $(\mathbf{y},\mathsf{p}) \in \mathbf{H}_0^1(\Omega) \times L_0^2(\Omega)$ to a control $u \in \mathcal{U}$ using the equations \eqref{eq:brinkamn_problem_state}.

The following result establishes differentiability properties for the operator $\mathcal{S}$.

\begin{theorem}[differentiability properties of $\mathcal{S}$]\label{diff_S}
The control to state map $\mathcal{S}:\mathcal{U}\to \mathbf{H}_0^1(\Omega)\times L_0^2(\Omega)$ is of class $\mathcal{C}^2$ with respect to the $L^{\infty}(\Omega)$-topology. Moreover, if $v\in L^{\infty}(\Omega)$, then $(\boldsymbol{\varphi},\zeta)=\mathcal{S}'(u)v\in \mathbf{H}_0^1(\Omega)\times L_0^2(\Omega)$ is the unique solution to
\begin{equation}\label{eq:problem_S_first}
\begin{array}{rl}
(\nabla \boldsymbol{\varphi},\nabla \mathbf{v})_{\mathbf{L}^2(\Omega)}+( u\boldsymbol{\varphi}, \mathbf{v})_{\mathbf{L}^2(\Omega)}-(\zeta,\mathrm{div~} \mathbf{v})_{L^2(\Omega)}
=
&
\hspace{-2mm}
-(v\mathbf{y},\mathbf{v})_{\mathbf{L}^2(\Omega)},
\\
(\mathsf{s},\mathrm{div~} \boldsymbol{\varphi})_{L^2(\Omega)}
=
&
\hspace{-2mm}
0
\end{array}
\end{equation}
for all $(\mathbf{v},\mathsf{s}) \in \mathbf{H}_0^1(\Omega) \times L_0^2(\Omega)$. Here, $(\mathbf{y},\mathsf{p})=\mathcal{S}u$. In addition, if $v_1,v_2 \in L^{\infty}(\Omega)$, then $(\boldsymbol{\psi},\xi)=\mathcal{S}''(u)(v_1,v_2)\in \mathbf{H}_0^1(\Omega)\times L_0^2(\Omega)$ is the unique solution to
\begin{multline}\label{eq:problem_S_second}
(\nabla \boldsymbol{\psi},\nabla \mathbf{v})_{\mathbf{L}^2(\Omega)}+( u\boldsymbol{\psi},\mathbf{v})_{\mathbf{L}^2(\Omega)}-(\xi,\mathrm{div~} \mathbf{v})_{L^2(\Omega)}
\\
=
-(v_2\boldsymbol{\varphi}_{v_1},\mathbf{v})_{\mathbf{L}^2(\Omega)}-(v_1\boldsymbol{\varphi}_{v_2},\mathbf{v})_{\mathbf{L}^2(\Omega)},
\qquad
(\mathsf{s},\mathrm{div~} \boldsymbol{\psi})_{L^2(\Omega)}  = 0
\end{multline}
for all $(\mathbf{v},\mathsf{s}) \in \mathbf{H}_0^1(\Omega) \times L_0^2(\Omega)$. Here, $(\boldsymbol{\varphi}_{v_i},\zeta_{v_i})=\mathcal{S}'(u)v_i$ and $i\in\{1,2\}$.
\end{theorem}
\begin{proof}
The first-order Fr\'echet differentiability of $\mathcal{S}$ follows from a simple modification of the arguments developed in the proof of \cite[Theorem 4.17]{MR2583281}. These arguments also show that $(\boldsymbol{\varphi},\zeta) = \mathcal{S}'(u)v$ corresponds to the unique solution to \eqref{eq:problem_S_first}. The second-order differentiability follows by applying similar arguments. The fact that $(\boldsymbol{\psi},\xi)$ solves \eqref{eq:problem_S_second} follows from the arguments developed in \cite[Theorem 4.24]{MR2583281}. Finally, we note that the well-posedness of problems \eqref{eq:problem_S_first} and \eqref{eq:problem_S_second} is trivial.
\end{proof}

\begin{remark}[non-differentiability]
The operator $\mathcal{S}$ is not Fr\'echet differentiable with respect to the $L^2(\Omega)$-topology; see \cite[Remark 2.8]{MR2536007}.
\end{remark}

To introduce the so-called \emph{reduced cost functional}, we define the map
\begin{align}\label{eq:mapG}
\mathcal{G}: \mathcal{U}\to \mathbf{H}_0^1(\Omega) :
\quad u\to \mathbf{y},
\end{align}
where $\mathbf{y}$ corresponds to the velocity component of the pair $(\mathbf{y},\mathsf{p})=\mathcal{S}u$. With these ingredients, we introduce the reduced cost functional $j:\mathcal{U}\to\mathbb{R}_0^+$ as $j(u) := J(\mathcal{G}u,u)$.

To obtain first-order optimality conditions, we present the following basic result: If $\bar{u} \in \mathbb{U}_{ad}$ is locally optimal for problem \eqref{eq:min_functional}--\eqref{eq:brinkamn_problem_state}, then $\bar{u}$ satisfies \cite[Lemma 4.18]{MR2583281}
\begin{equation}\label{eq:ineq_variational}
j'(\bar{u})(u-\bar{u})\geq 0\quad \forall u\in \mathbb{U}_{ad},
\end{equation}
where $j'(\bar{u})$ denotes the Gate\^aux derivative of $j$ at $\bar{u}$. To explore the inequality \eqref{eq:ineq_variational}, we introduce $(\mathbf{z},\mathsf{r})\in\mathbf{H}_0^1(\Omega)\times L_0^2(\Omega)$ as the solution to the \emph{adjoint equations}
\begin{equation}\label{eq:brinkamn_problem_adjoint}
\begin{array}{rl}
(\nabla \mathbf{v}, \nabla \mathbf{z})_{\mathbf{L}^2(\Omega)}
+
( u\mathbf{z}, \mathbf{v})_{\mathbf{L}^2(\Omega)}
+
(\mathsf{r},\text{div }\mathbf{v})_{L^2(\Omega)} = &\hspace{-2mm} (\mathbf{y}-\mathbf{y}_{\Omega},\mathbf{v})_{\mathbf{L}^2(\Omega)},
\\
(\mathsf{s},\text{div }\mathbf{z})_{L^2(\Omega)}  = & \hspace{-2mm} 0
\end{array}
\end{equation}
for all $(\mathbf{v},\mathsf{s}) \in \mathbf{H}_0^1(\Omega) \times L_0^2(\Omega)$. Here, $(\mathbf{y},\mathsf{p})=\mathcal{S}u$ is the solution to \eqref{eq:brinkamn_problem_state}. We note that the adjoints equations \eqref{eq:brinkamn_problem_adjoint} are well-posed.

We are now in a position to establish necessary first-order optimality conditions as in the scalar case described in \cite[identity (2.6)]{MR2536007}.

\begin{theorem}[first-order optimality conditions]
Every locally optimal control $\bar{u}\in \mathbb{U}_{ad}$ for the optimal control problem \eqref{eq:min_functional}--\eqref{eq:brinkamn_problem_state} satisfies the variational inequality
\begin{eqnarray}\label{eq:variational_inequality}
(\alpha \bar{u}-\bar{\mathbf{y}}\cdot \bar{\mathbf{z}},u-\bar{u})_{L^2(\Omega)}\geq 0\quad \forall u\in \mathbb{U}_{ad}.
\end{eqnarray}
Here, $(\bar{\mathbf{z}},\bar{\mathsf{r}})$ denotes the solution to the adjoint equations \eqref{eq:brinkamn_problem_adjoint}, where $u$ and $\mathbf{y}$ are replaced by $\bar{u}$ and $\bar{\mathbf{y}}=\mathcal{G}\bar{u}$, respectively.
\end{theorem}

\begin{proof}
As a first step, we rewrite the variational inequality \eqref{eq:ineq_variational} as follows:
\begin{eqnarray}\label{eq:variat_ineq_01}
(\mathcal{G}'(\bar{u})(u-\bar{u}),\bar{\mathbf{y}}-\mathbf{y}_{\Omega})_{\mathbf{L}^2(\Omega)}
+
(\alpha \bar{u},u-\bar{u})_{L^2(\Omega)}\geq 0
\quad
\forall u\in \mathbb{U}_{ad},
\end{eqnarray}
where $\bar{\mathbf{y}}=\mathcal{G}\bar{u}$. Define $(\boldsymbol{\varphi},\zeta):=\mathcal{S}'(\bar{u})(u-\bar{u})$ and note that $(\boldsymbol{\varphi},\zeta)$ solves problem \eqref{eq:problem_S_first} with $u$ and $v$ replaced by $\bar{u}$ and $u-\bar{u}$, respectively. We now set $\mathbf{v}=\bar{\mathbf{z}}$ in the first equation of this system to obtain
$
(\nabla \boldsymbol{\varphi},\nabla \bar{\mathbf{z}})_{\mathbf{L}^2(\Omega)}
+
( \bar{u}\boldsymbol{\varphi}, \bar{\mathbf{z}})_{\mathbf{L}^2(\Omega)}
-
(\zeta,\mathrm{div~} \bar{\mathbf{z}})_{L^2(\Omega)}
=  
-((u-\bar{u})\bar{\mathbf{y}},\bar{\mathbf{z}})_{\mathbf{L}^2(\Omega)}.
$
Similarly, we set $\mathbf{v} = \boldsymbol{\varphi}$ as a test function in the the first equation of \eqref{eq:brinkamn_problem_adjoint} to obtain
$
(\nabla \boldsymbol{\varphi},\nabla \bar{\mathbf{z}})_{\mathbf{L}^2(\Omega)}
+
( \bar{u}\bar{\mathbf{z}}, \boldsymbol{\varphi})_{\mathbf{L}^2(\Omega)}
+
(\mathsf{r},\text{div }\boldsymbol{\varphi})_{L^2(\Omega)} = (\bar{\mathbf{y}}-\mathbf{y}_{\Omega},\boldsymbol{\varphi})_{\mathbf{L}^2(\Omega)}.
$
Since $(\zeta,\mathrm{div~} \bar{\mathbf{z}})_{L^2(\Omega)} = 0$ and $(\mathsf{r},\text{div }\boldsymbol{\varphi})_{L^2(\Omega)} = 0$ for $\zeta \in L_0^2(\Omega)$ and $\mathsf{r} \in L_0^2(\Omega)$, respectively, we can conclude that $(\bar{\mathbf{y}}-\mathbf{y}_{\Omega},\boldsymbol{\varphi})_{\mathbf{L}^2(\Omega)} = -((u-\bar{u})\bar{\mathbf{y}},\bar{\mathbf{z}})_{\mathbf{L}^2(\Omega)}$. The desired variational inequality \eqref{eq:variational_inequality} thus follows from \eqref{eq:variat_ineq_01}. Note that $\boldsymbol{\varphi} = \mathcal{G}'(\bar u)(u - \bar u)$.
\end{proof}

The following projection formula will be important for the derivation of regularity estimates: If $\bar{u} \in \mathbb{U}_{ad}$ is a locally optimal control for  \eqref{eq:min_functional}--\eqref{eq:brinkamn_problem_state}, then \cite[section 4.6.1]{MR2583281}
\begin{equation}\label{eq:representation}
\bar{u}(x):=\Pi_{[\mathsf{a},\mathsf{b}]}(\alpha^{-1}\bar{\mathbf{y}}\cdot\bar{\mathbf{z}})\text{ a.e. } x \in \Omega,
\end{equation}
where $\Pi_{[\mathsf{a},\mathsf{{b}]}}:L^1(\Omega)\to\mathbb{U}_{ad}$ is defined by $\Pi_{[\mathsf{a},\mathsf{b}]}(v):=\min\{\mathsf{b},\max\{v,\mathsf{a}\}\}$ a.e.~in $\Omega$.

The Sobolev regularity of an optimal control variable is as follows.

\begin{theorem}[regularity of an optimal control]
\label{thm:regularity_optimal_control}
Let $\bar{u}$ be a locally optimal control for problem \eqref{eq:min_functional}--\eqref{eq:brinkamn_problem_state}. If $\Omega$ is Lipschitz and $\mathbf{f},\mathbf{y}_{\Omega}\in\mathbf{L}^2(\Omega)$, then $\bar{u}\in H^1(\Omega)$ and
\begin{align}\label{eq:control_regul_estimate}
\|\bar{u}\|_{H^1(\Omega)}
\lesssim
\|\mathbf{f}\|_{\mathbf{L}^2(\Omega)}
\left[
\|\mathbf{f}\|_{\mathbf{L}^2(\Omega)}+\|\mathbf{y}_{\Omega}\|_{\mathbf{L}^2(\Omega)}
\right].
\end{align}
\end{theorem} 
\begin{proof}
Let $(\bar{\mathbf{y}},\bar{\mathsf{p}})=\mathcal{S}\bar{u}$, and let $(\mathbf{z},\mathsf{r})\in\mathbf{H}_0^1(\Omega)\times L_0^2(\Omega)$ be the solution to  \eqref{eq:brinkamn_problem_adjoint} with $u$ and $\mathbf{y}$ replaced by $\bar{u}$ and $\bar{\mathbf{y}}$, respectively. Since $\mathbf{f},\mathbf{y}_{\Omega}\in\mathbf{L}^2(\Omega)$, we can apply the result in Theorem \ref{theorem_regularity_Linfty} to immediately obtain that $\bar{\mathbf{y}},\bar{\mathbf{z}}\in\mathbf{H}_0^1(\Omega)\cap\mathbf{L}^{\infty}(\Omega)$. A simple calculation thus shows that $\nabla (\bar{\mathbf{y}}\cdot\bar{\mathbf{z}}) \in \mathbf{L}^2(\Omega)$ Consequently, the desired regularity property and the bound \eqref{eq:control_regul_estimate} for $\bar{u}$ follow from the projection formula \eqref{eq:representation}.
\end{proof}

\subsection{Second-order optimality conditions}\label{section:secondconditions}
In this section, we derive necessary and sufficient second-order optimality conditions. To this end, we begin our analysis with the following instrumental result.

\begin{theorem}[$j$ is of class $C^2$ and $j''$ is Lipschitz]
\label{eq:theorem_properties_j}
The reduced cost $j: \mathbb{U}_{ad} \subset \mathcal{U}\to\mathbb{R}_0^+$ is of class $C^2$. Moreover, for each $u\in\mathbb{U}_{ad}$ and $v\in L^{\infty}(\Omega)$, it holds that
\begin{align}
\label{eq:identity_j}
j''(u)v^2=\alpha\|v\|_{L^2(\Omega)}^2-2(v\boldsymbol{\varphi},\mathbf{z})_{\mathbf{L}^2(\Omega)}+\|\boldsymbol{\varphi}\|_{\mathbf{L}^2(\Omega)}^2.
\end{align}
Here, $(\mathbf{z},\mathsf{r})\in\mathbf{H}_0^1(\Omega)\times L_0^2(\Omega)$ denotes the solution to \eqref{eq:brinkamn_problem_adjoint} and $(\boldsymbol{\varphi},\zeta)=\mathcal{S}'(u)v$. For every $u_1,u_2\in\mathbb{U}_{ad}$ and $v\in L^{\infty}(\Omega)$, we have the following bound:
\begin{eqnarray}
\label{eq:lemma_j}
|j''(u_1)v^2-j''(u_2)v^2|\leq \mathfrak{D}\|u_1-u_2\|_{L^2(\Omega)}\|v\|_{L^2(\Omega)}^2,
\end{eqnarray}
where $\mathfrak{D}>0$ is a constant that depends on $\|\mathbf{f}\|_{\mathbf{H}^{-1}(\Omega)}$ and $\|\mathbf{y}_\Omega\|_{\mathbf{L}^2(\Omega)}$.
\end{theorem}

\begin{proof}
The fact that $j$ belongs to the class $C^2$ is a consequence of the differentiability properties of the control-to-state map $\mathcal{S}$ given in Theorem \ref{diff_S}. Let us now prove \eqref{eq:identity_j} and \eqref{eq:lemma_j} by means of three steps.

\emph{Step 1.} \emph{The identity \eqref{eq:identity_j}:} For every $u\in\mathbb{U}_{ad}$ and $v \in L^{\infty}(\Omega)$, we write
\begin{align}\label{eq:second_derivative_aux}
j''(u)v^2 =\alpha\|v\|_{L^2(\Omega)}^2+(\mathbf{y}-\mathbf{y}_{\Omega},\boldsymbol{\psi})_{\mathbf{L}^2(\Omega)}+\|\boldsymbol{\varphi}\|_{\mathbf{L}^2(\Omega)}^2.
\end{align}
Here, $(\boldsymbol{\varphi},\zeta) = \mathcal{S}'(u)v$, i.e., $(\boldsymbol{\varphi},\zeta)$ corresponds to the solution to \eqref{eq:problem_S_first} and
$(\boldsymbol{\psi},\xi)$ is defined as the solution to \eqref{eq:problem_S_second} with $v_1 = v_2 = v$ and $\boldsymbol{\varphi}_{v} = \boldsymbol{\varphi}$. Note that $\boldsymbol{\varphi}=\mathcal{G}'(u)v$ and $\boldsymbol{\psi}=\mathcal{G}''(u)v^2$. We now set $(\mathbf{v},\mathsf{s})=(\boldsymbol{\psi},0)$ in \eqref{eq:brinkamn_problem_adjoint} and $(\mathbf{v},\mathsf{s}) = (\mathbf{z},0)$ in \eqref{eq:problem_S_second} to obtain the identity $(\mathbf{y}-\mathbf{y}_{\Omega},\boldsymbol{\psi})_{\mathbf{L}^2(\Omega)}=-2(v\boldsymbol{\varphi},\mathbf{z})_{\mathbf{L}^2(\Omega)}$. Replacing this relation into \eqref{eq:second_derivative_aux} gives the desired identity \eqref{eq:identity_j}.


\emph{Step 2:} \emph{The bound \eqref{eq:lemma_j}:} Let $u_1,u_2\in \mathbb{U}_{ad}$ and $v\in L^{\infty}(\Omega)$. Define $(\boldsymbol{\varphi}_i,\zeta_i):=\mathcal{S}'(u_i)v$ with $i \in \{ 1,2 \}$. We note that $(\boldsymbol{\varphi}_i,\zeta_i)$ solves \eqref{eq:problem_S_first} where $u$ and $\mathbf{y}$ are replaced by $u_i$ and $\mathbf{y}_i$, respectively. A simple application of the identity \eqref{eq:identity_j} results in
\begin{multline}
\label{eq:aux_identity_propj}
j''(u_1)v^2-j''(u_2)v^2=2(v(\boldsymbol{\varphi}_2-\boldsymbol{\varphi}_1),\mathbf{z}_2)_{\mathbf{L}^2(\Omega)}+2(v\boldsymbol{\varphi}_1,\mathbf{z}_2-\mathbf{z}_1)_{\mathbf{L}^2(\Omega)}\\
(\boldsymbol{\varphi}_1-\boldsymbol{\varphi}_2,\boldsymbol{\varphi}_1+\boldsymbol{\varphi}_2)_{\mathbf{L}^2(\Omega)}=:\text{I}+\text{II}+\text{III},
\end{multline}
where $(\mathbf{z}_i,\mathsf{r}_i) \in\mathbf{H}_0^1(\Omega) \times L_0^2(\Omega)$ corresponds to the solution of the adjoint problem \eqref{eq:brinkamn_problem_adjoint}, where $\mathbf{y}$ and $u$ are replaced by $\mathbf{y}_i$ and $u_i$, respectively. In the following, we analyze each of the terms \text{I}, \text{II}, and \text{III} involved in the identity \eqref{eq:aux_identity_propj}.

We first bound \text{I}. Given the embedding $\mathbf{H}_0^1(\Omega) \hookrightarrow \mathbf{L}^{r}(\Omega)$ with $r<\infty$ when $d=2$ and $r\leq6$ when $d=3$, and the fact that $\mathbf{f} \in \mathbf{H}^{-1}(\Omega)$ and $\mathbf{y}_{\Omega}\in\mathbf{L}^2(\Omega)$, we obtain
\begin{align}\label{eq:estimates_extra}
\|\mathbf{y}_i\|_{\mathbf{L}^{r}(\Omega)}\lesssim \|\mathbf{f}\|_{\mathbf{H}^{-1}(\Omega)},
\quad 
\|\mathbf{z}_i\|_{\mathbf{L}^{r}(\Omega)}\lesssim \|\mathbf{f}\|_{\mathbf{H}^{-1}(\Omega)}+\|\mathbf{y}_{\Omega}\|_{\mathbf{L}^{2}(\Omega)},
\quad
i\in\{1,2\}.
\end{align}
We thus bound $\text{I}$ as follows: $\text{I} \leq 2 \|v\|_{L^2(\Omega)}\|\boldsymbol{\varphi}_2-\boldsymbol{\varphi}_1\|_{\mathbf{L}^4(\Omega)}\|\mathbf{z}_2\|_{\mathbf{L}^4(\Omega)}$ $\lesssim$ $\|v\|_{L^2(\Omega)}\|\nabla(\boldsymbol{\varphi}_2-\boldsymbol{\varphi}_1)\|_{\mathbf{L}^2(\Omega)}$, where the hidden constant depends on $\|\mathbf{f}\|_{\mathbf{H}^{-1}(\Omega)}$ and $\|\mathbf{y}_\Omega\|_{\mathbf{L}^2(\Omega)}$. To control $\|\nabla(\boldsymbol{\varphi}_2-\boldsymbol{\varphi}_1)\|_{\mathbf{L}^2(\Omega)}$, we invoke the problem that solves $(\boldsymbol{\varphi}_2-\boldsymbol{\varphi}_1,\zeta_2-\zeta_1)$:
\begin{multline*}
(\nabla (\boldsymbol{\varphi}_2-\boldsymbol{\varphi}_1),\nabla \mathbf{v})_{\mathbf{L}^2(\Omega)}+( u_2(\boldsymbol{\varphi}_2-\boldsymbol{\varphi}_1), \mathbf{v})_{\mathbf{L}^2(\Omega)}-(\zeta_2-\zeta_1,\text{div }\mathbf{v})_{L^2(\Omega)} \\= (\boldsymbol{\varphi}_1(u_1-u_2)-v(\mathbf{y}_2-\mathbf{y}_1),\mathbf{v})_{\mathbf{L}^2(\Omega)},\qquad
(\mathsf{s},\text{div }(\boldsymbol{\varphi}_2-\boldsymbol{\varphi}_1))_{L^2(\Omega)}  =  0,
\end{multline*}
for all $(\mathbf{v},\mathsf{s}) \in \mathbf{H}_0^1(\Omega) \times L_0^2(\Omega)$. If we set $\mathbf{v}=\boldsymbol{\varphi}_2-\boldsymbol{\varphi}_1$ and $\mathsf{s}=0$, we obtain
\begin{multline}\label{eq:estimate_aux1_theorem_second}
\|\nabla (\boldsymbol{\varphi}_2-\boldsymbol{\varphi}_1)\|_{\mathbf{L}^2(\Omega)}\lesssim \|\nabla\boldsymbol{\varphi}_1\|_{\mathbf{L}^{2}(\Omega)}\|u_1-u_2\|_{L^2(\Omega)}\\+\|v\|_{L^2(\Omega)}\|\nabla(\mathbf{y}_2-\mathbf{y}_1)\|_{\mathbf{L}^{2}(\Omega)}.
\end{multline}
A bound for $\|\nabla\boldsymbol{\varphi}_1\|_{\mathbf{L}^{2}(\Omega)}$ can be derived using stability estimates for the problems that $(\boldsymbol{\varphi}_1,\zeta_1)$ and $(\mathbf{y}_1,\mathsf{p}_1)$ solve. In fact,
\begin{align}\label{eq:estimate_aux2_theorem_second}
\|\nabla\boldsymbol{\varphi}_1\|_{\mathbf{L}^{2}(\Omega)}
\lesssim 
\|v \mathbf{y}_1\|_{\mathbf{H}^{-1}(\Omega)}
\lesssim \|v\|_{L^2(\Omega)} \|\mathbf{y}_1\|_{\mathbf{L}^{4}(\Omega)}
\lesssim \|v\|_{L^2(\Omega)}\|\mathbf{f}\|_{\mathbf{H}^{-1}(\Omega)}.
\end{align}
On the other hand, we note that $(\mathbf{y}_2-\mathbf{y}_1,\mathsf{p}_2-\mathsf{p}_1)\in\mathbf{H}_0^1(\Omega)\times L_0^2(\Omega)$ solves
\begin{multline*}
(\nabla (\mathbf{y}_2-\mathbf{y}_1),\nabla \mathbf{v})_{\mathbf{L}^2(\Omega)}+( u_2(\mathbf{y}_2-\mathbf{y}_1), \mathbf{v})_{\mathbf{L}^2(\Omega)}-(\mathsf{p}_2-\mathsf{p}_1,\text{div }\mathbf{v})_{L^2(\Omega)} \\
= (\mathbf{y}_1(u_1-u_2),\mathbf{v})_{\mathbf{L}^2(\Omega)},
\quad
(\mathsf{s},\text{div }(\mathbf{y}_2-\mathbf{y}_1))_{L^2(\Omega)}  =  0
\quad
\forall (\mathbf{v},\mathsf{s}) \in \mathbf{H}_0^1(\Omega) \times L_0^2(\Omega).
\end{multline*}
If we set $(\mathbf{v},\mathsf{s}) = (\mathbf{y}_2-\mathbf{y}_1,0)$ and use basic arguments, we obtain the bound
\begin{multline}\label{eq:estimate_aux3_theorem_second}
\|\nabla(\mathbf{y}_2-\mathbf{y}_1)\|_{\mathbf{L}^{2}(\Omega)}
\lesssim 
\|\nabla\mathbf{y}_1\|_{\mathbf{L}^{2}(\Omega)} \|u_1-u_2\|_{L^2(\Omega)}
\lesssim 
\|\mathbf{f}\|_{\mathbf{H}^{-1}(\Omega)}
\|u_1-u_2\|_{L^2(\Omega)}.
\end{multline} 
A collection of the estimates \eqref{eq:estimate_aux1_theorem_second}, \eqref{eq:estimate_aux2_theorem_second}, and \eqref{eq:estimate_aux3_theorem_second} leads to the conclusion
\begin{align}\label{eq:I_estimates}
\text{I}\lesssim\|v\|^2_{L^2(\Omega)}\|u_1-u_2\|_{L^2(\Omega)},
\end{align}
with a hidden constant that depends on $\|\mathbf{f}\|_{\mathbf{H}^{-1}(\Omega)}$ and $\|\mathbf{y}_\Omega\|_{\mathbf{L}^2(\Omega)}$. 

In the following, we control $\text{II}$ in \eqref{eq:aux_identity_propj}. To do so, we use H\"older's inequality and \eqref{eq:estimate_aux2_theorem_second} to obtain  $\text{II}\lesssim \|v\|_{L^2(\Omega)}\|\boldsymbol{\varphi}_1\|_{\mathbf{L}^4(\Omega)}\|\mathbf{z}_2-\mathbf{z}_1\|_{\mathbf{L}^{4}(\Omega)}\lesssim \|v\|_{L^2(\Omega)}^2 \|\nabla(\mathbf{z}_2-\mathbf{z}_1)\|_{\mathbf{L}^{2}(\Omega)}$. To control $\|\nabla(\mathbf{z}_2-\mathbf{z}_1)\|_{\mathbf{L}^{2}(\Omega)}$, we first note that  $(\mathbf{z}_2-\mathbf{z}_1,\mathsf{r}_2-\mathsf{r}_1)$ solves
\begin{multline*}
(\nabla \mathbf{v},\nabla (\mathbf{z}_2-\mathbf{z}_1))_{\mathbf{L}^2(\Omega)}+( u_2(\mathbf{z}_2-\mathbf{z}_1), \mathbf{v})_{\mathbf{L}^2(\Omega)}+(\mathsf{r}_2-\mathsf{r}_1,\text{div }\mathbf{v})_{L^2(\Omega)} \\= ((\mathbf{y}_2-\mathbf{y}_1)+\mathbf{z}_1(u_1-u_2),\mathbf{v})_{\mathbf{L}^2(\Omega)},\qquad
(\mathsf{s},\text{div }(\mathbf{z}_2-\mathbf{z}_1))_{L^2(\Omega)}  =  0,
\end{multline*}
for all $(\mathbf{v},\mathsf{s}) \in \mathbf{H}_0^1(\Omega) \times L_0^2(\Omega)$. If we set $\mathbf{v}=\mathbf{z}_2-\mathbf{z}_1$ and $\mathsf{s}=0$, we immediately obtain
\begin{align}
\|\nabla(\mathbf{z}_2-\mathbf{z}_1)\|_{\mathbf{L}^{2}(\Omega)}\lesssim \|\nabla(\mathbf{y}_2-\mathbf{y}_1)\|_{\mathbf{L}^{2}(\Omega)}+\|\nabla\mathbf{z}_1\|_{\mathbf{L}^{2}(\Omega)}\|u_2-u_1\|_{L^2(\Omega)}.
\end{align}
The bound $\|\nabla(\mathbf{y}_2-\mathbf{y}_1)\|_{\mathbf{L}^{2}(\Omega)}\lesssim\|u_2-u_1\|_{L^2(\Omega)}$, which follows from \eqref{eq:estimate_aux3_theorem_second}, and the stability estimate $\|\nabla\mathbf{z}_1\|_{\mathbf{L}^{2}(\Omega)} \lesssim \|\mathbf{f}\|_{\mathbf{H}^{-1}(\Omega)} + \|\mathbf{y}_{\Omega}\|_{\mathbf{L}^{2}(\Omega)}$ allow us to obtain
\begin{align}\label{eq:II_estimates}
\text{II}\lesssim \|v\|_{L^2(\Omega)}^2 \|\nabla(\mathbf{z}_2-\mathbf{z}_1)\|_{\mathbf{L}^{2}(\Omega)}\lesssim \|v\|_{L^2(\Omega)}^2 \|u_2-u_1\|_{L^2(\Omega)}.
\end{align}

Finally, we control $\text{III}$. To accomplish this task, we use basic inequalities, \eqref{eq:estimate_aux2_theorem_second}, and an adjustment of \eqref{eq:estimate_aux2_theorem_second} for $\boldsymbol{\varphi}_2$, and $\|\boldsymbol{\varphi}_2-\boldsymbol{\varphi}_1\|_{\mathbf{L}^2(\Omega)}\lesssim\|v\|_{L^2(\Omega)}\|u_1-u_2\|_{L^2(\Omega)}$:
\begin{align}\label{eq:III_estimates}
\text{III}\lesssim \|v\|_{L^2(\Omega)}\|\boldsymbol{\varphi}_2-\boldsymbol{\varphi}_1\|_{\mathbf{L}^2(\Omega)}\lesssim\|v\|_{L^2(\Omega)}^2 \|u_2-u_1\|_{L^2(\Omega)}.
\end{align}

\emph{Step 3:} In view of \eqref{eq:aux_identity_propj}, we infer the desired estimate \eqref{eq:lemma_j} by collecting the bounds derived in \eqref{eq:I_estimates}, \eqref{eq:II_estimates}, \eqref{eq:III_estimates}. This concludes the proof.
\end{proof}

\subsubsection{Necessary second-order optimality conditions}

We follow the arguments developed in the proof of \cite[Theorem 23]{MR3586845} and present necessary second-order optimality conditions. To this end, we introduce some concepts: Let $(\bar{u},\bar{\mathbf{y}},\bar{\mathbf{z}})\in\mathbb{U}_{ad}\times \mathbf{H}_0^1(\Omega)\times \mathbf{H}_0^1(\Omega)$ satisfy the first-order optimality conditions \eqref{eq:brinkamn_problem_state}, \eqref{eq:brinkamn_problem_adjoint}, and \eqref{eq:variational_inequality}. We define $\bar{\mathfrak{d}}:=\alpha\bar{u}-\bar{\mathbf{y}}\cdot\bar{\mathbf{z}}$ and the cone of critical directions
\begin{align}\label{eq:cone}
C_{\bar{u}}:=\{v\in L^{2}(\Omega) \text{ satisfying } \eqref{eq:cone_condition} \text{ and } v(x)=0 \text{ if } \bar{\mathfrak{d}}(x)\neq 0\},
\end{align}
where condition \eqref{eq:cone_condition} reads as follows:
\begin{align}\label{eq:cone_condition}
v(x)\geq 0 \text{ a.e. } x\in\Omega	\text{ if } \bar{u}(x)=\mathsf{a},\quad v(x)\leq 0 \text{ a.e. } x\in\Omega	\text{ if } \bar{u}(x)= \mathsf{b}.
\end{align}

We are now in a position to present necessary second-order optimality conditions.

\begin{theorem}[necessary second-order optimality conditions]\label{theorem_neccesary_conditions} If $\bar{u}\in \mathbb{U}_{ad}$ denotes a locally optimal control for \eqref{eq:min_functional}--\eqref{eq:brinkamn_problem_state}, then $j''(\bar{u})v^2\geq 0$ for all $v\in C_{\bar{u}}$.
\end{theorem}
\begin{proof}
The proof follows from an adaptation of the arguments from \cite[Theorem 23]{MR3586845} and an application of Theorem \ref{eq:theorem_properties_j}. In a first step, we define for $w\in C_{\bar{u}}$, for a.e.~$x\in\Omega$, and for $k\in\mathbb{N}$, the function 
\begin{eqnarray*}
w_k(x):=\left\{
\begin{array}{cl}
0 & \text{if } x: \mathsf{a}<\bar{u}(x)<\mathsf{a}+k^{-1},\quad \mathsf{b}-k^{-1}<\bar{u}(x)<\mathsf{b},\\
\Pi_{[-k,k]}(w(x)) & \text{otherwise}.
\end{array}
\right.
\end{eqnarray*}
Starting from the fact that $w\in C_{\bar{u}}$, we obtain that $w_k\in C_{\bar{u}}\cap L^{\infty}(\Omega)$ for $k \in \mathbb{N}$. We also note that $w_k(x)\to w(x)$ as $k\uparrow\infty$ for a.e.~$x\in\Omega$ and that $|w_k(x)|\leq |w(x)|$ for a.e.~$x\in\Omega$. As a consequence of Lebesgue's dominated convergence theorem, we have that $w_k\to w$ in $L^2(\Omega)$ as $k\uparrow\infty$. We now define $\rho^{\star}=\min\{k^{-2},(b-a)k^{-1}\}$ and note that for $\rho \in (0, \rho^{\star}]$ we have $\bar{u}+\rho w_k \in \mathbb{U}_{ad}$. Thus, since $\bar{u}+\rho w_k$ is admissible, we use the fact that $\bar{u}$ is a local minimizer to obtain the inequality $j(\bar{u})\leq j(\bar{u}+\rho w_k)$ for $\rho$ sufficiently small. We now apply Taylor's theorem for $j$ at $\bar{u}$, and the relation $j'(\bar{u})w_k=0$, which follows from the fact that $w_k\in C_{\bar{u}}$, to obtain 
\begin{align*}
0\leq j(\bar{u}+\rho w_k)-j(\bar{u})=\rho j'(\bar{u})w_k+\frac{\rho^2}{2}j''(\bar{u}+\rho \theta_k w_k)w_k^2 =\frac{\rho^2}{2}j''(\bar{u}+\rho\theta_k w_k)w_k^2,
\end{align*}
where $\theta_k\in(0,1)$. Based on the estimate \eqref{eq:lemma_j}, we obtain the following result as $\rho\downarrow 0$:
\begin{align*}
|j''(\bar{u}+\rho\theta_k w_k)w_k^2-j''(\bar{u})w_k^2|\lesssim\rho \|w_k\|_{L^2(\Omega)}^3\to 0,
\qquad \rho\downarrow 0.
\end{align*}
This implies that $j''(\bar{u})w_k^2\geq 0$ for every $k\in\mathbb{N}$. We now use the convergence $w_k\to w$ in $L^2(\Omega)$ as $k\uparrow\infty$ and the identity \eqref{eq:identity_j} to obtain that $j''(\bar{u})w^2 \geq 0$ for all $w \in C_{\bar{u}}$.
\end{proof}

\subsubsection{Sufficient second-order optimality conditions}

We now establish sufficient second-order optimality conditions \cite[Theorem 2.3]{MR2902693}, \cite[Theorem 23]{MR3586845}.

\begin{theorem}[sufficient second-order optimality conditions]\label{theorem_second_order_optimal}
Let $(\bar{u},\bar{\mathbf{y}},\bar{\mathbf{z}}) \in \mathbb{U}_{ad} \times \mathbf{H}_0^1(\Omega) \times \mathbf{H}_0^1(\Omega)$ satisfy the first-order optimality conditions \eqref{eq:brinkamn_problem_state}, \eqref{eq:brinkamn_problem_adjoint}, and \eqref{eq:variational_inequality}. If $j''(\bar{u})v^2> 0$ for all $v\in C_{\bar{u}}\setminus\{0\}$, then there exist $\delta>0$ and $\sigma>0$ such that
\begin{eqnarray}\label{eq:second_order_condition}
j(u)\geq j(\bar{u})+\frac{\sigma}{2}\|u-\bar{u}\|_{L^2(\Omega)}^2\quad\forall u\in\mathbb{U}_{ad}\cap B_{\delta}(\bar{u}),
\end{eqnarray}
where $B_{\delta}(\bar{u})$ denotes the closed ball in $L^2(\Omega)$ of radius $\delta$ and center $\bar{u}$.
\end{theorem}

\begin{proof}
We proceed by contradiction and assume that for every $k\in\mathbb{N}$ we can find an element $u_k\in\mathbb{U}_{ad}$ such that
\begin{align}\label{eq:estimate_contradiction}
\|\bar{u}-u_k\|_{L^2(\Omega)}<k^{-1},
\quad
j(u_k)<j(\bar{u})+(2k)^{-1}\|\bar{u}-u_k\|_{L^2(\Omega)}^2.
\end{align}
For $k\in\mathbb{N}$, we define $\rho_k := \|u_k - \bar{u}\|_{L^2(\Omega)}$ and $w_k:=(u_k-\bar{u})/\rho_k$. It is clear that, for every $k\in\mathbb{N}$, $\|w_k\|_{L^2(\Omega)}=1$. We can therefore deduce the existence of a nonrelabeled subsequence $\{w_k\}_{k\geq 1}\subset L^2(\Omega)$ such that $w_k \rightharpoonup w$ in $L^2(\Omega)$ as $k \uparrow\infty$. In the following, we prove that the limit point $w$ belongs to the cone $C_{\bar{u}}$ and that $w\equiv 0$.

\emph{Step 1:} $w\in C_{\bar{u}}$. We first note that the set of elements satisfying \eqref{eq:cone_condition} is closed and convex in $L^2(\Omega)$. Consequently, this set is weakly sequentially closed in $L^2(\Omega)$, and the weak limit $w$ satisfies \eqref{eq:cone_condition}. To prove that $w\in C_{\bar{u}}$, we thus need to verify the remaining condition in \eqref{eq:cone}. To do so, we apply the mean value theorem and the right-hand side estimate in \eqref{eq:estimate_contradiction} to obtain
\begin{align}\label{eq:estimate_jprime}
j'(\tilde{u}_k)w_k=\rho_k^{-1}(j(u_k)-j(\bar{u}))< (2k)^{-1}\rho_k\to 0,\quad k\uparrow \infty,
\end{align}
where $\tilde{u}_k=\bar{u}+\theta_k (u_k-\bar{u})$ and $\theta_k\in (0,1)$. We now define $(\tilde{\mathbf{y}}_k,\tilde{\mathsf{p}}_k):=\mathcal{S}\tilde{u}_k$ and $(\tilde{\mathbf{z}}_k,\tilde{\mathsf{r}}_k)\in\mathbf{H}_0^1(\Omega)\times L_0^2(\Omega)$ as the solution to problem \eqref{eq:brinkamn_problem_adjoint}, where $\mathbf{y}$ and $u$ are replaced by $\tilde{\mathbf{y}}_k$ and $\tilde{u}_k$, respectively. A basic stability bound for the problem that $(\bar{\mathbf{y}}-\tilde{\mathbf{y}}_k,\bar{\mathsf{p}}-\tilde{\mathsf{p}}_k)$ solves combined with the fact that $\tilde{u}_k \to \bar{u}$ in $L^2(\Omega)$ as $k\uparrow \infty$ results in $\tilde{\mathbf{y}}_k \to \bar{\mathbf{y}}$ in $\mathbf{H}_0^1(\Omega)$ as $k\uparrow \infty$. Using similar arguments as above, we can conclude that $\tilde{\mathbf{z}}_k\to \bar{\mathbf{z}}$ in $\mathbf{H}_0^1(\Omega)$ as $k\uparrow \infty$. With these convergence results at hand, we obtain  $\alpha\tilde{u}_k-\tilde{\mathbf{y}}_k\cdot\tilde{\mathbf{z}}_k=:\tilde{\mathfrak{d}}_k\to\bar{\mathfrak{d}} = \alpha\bar{u}-\bar{\mathbf{y}}\cdot\bar{\mathbf{z}}$ in $L^2(\Omega)$ as $k\uparrow \infty$. In fact, we have
$
\|\tilde{\mathfrak{d}}_k-\bar{\mathfrak{d}}\|_{L^2(\Omega)}
\lesssim
\alpha \|\tilde{u}_k-\bar{u}\|_{L^2(\Omega)}
+
\|\nabla\tilde{\mathbf{y}}_k\|_{\mathbf{L}^2(\Omega)}\|\nabla( \tilde{\mathbf{z}}_k-\bar{\mathbf{z}})\|_{\mathbf{L}^2(\Omega)}
+ 
\|\nabla\bar{\mathbf{z}}\|_{\mathbf{L}^2(\Omega)}\|\nabla( \tilde{\mathbf{y}}_k-\bar{\mathbf{y}})\|_{\mathbf{L}^2(\Omega)} \rightarrow 0$ as $k \uparrow \infty$, where we have used $\|\nabla\bar{\mathbf{z}}\|_{\mathbf{L}^2(\Omega)}\lesssim \|\mathbf{f}\|_{\mathbf{H}^{-1}(\Omega)}+\|\mathbf{y}_\Omega\|_{\mathbf{L}^2(\Omega)}$ and $\|\nabla\tilde{\mathbf{y}}_k\|_{\mathbf{L}^2(\Omega)}\lesssim \|\mathbf{f}\|_{\mathbf{H}^{-1}(\Omega)}$, which is uniform in $k$. We now use $w_k \rightharpoonup w$ in $L^2(\Omega)$ as $k \uparrow\infty$ and \eqref{eq:estimate_jprime} to arrive at $j'(\bar{u})w = ( \bar{\mathfrak{d}}, w )_{L^2(\Omega)} = \lim_{k\uparrow \infty} ( \tilde{\mathfrak{d}}_k,w_k)_{L^2(\Omega)} =\lim_{k\uparrow \infty}j'(\tilde{u}_k)w_k\leq 0$. On the other hand, the first-order condition \eqref{eq:variational_inequality} shows that $\int_{\Omega}\bar{\mathfrak{d}}(x)w(x)\mathrm{d}x \geq 0$. As a result, we deduce that $\int_{\Omega}|\bar{\mathfrak{d}}(x)w(x)|\mathrm{d}x=\int_{\Omega}\bar{\mathfrak{d}}(x)w(x)\mathrm{d}x=0$. We can thus conclude that $\bar{\mathfrak{d}} \neq 0$ implies that $w=0$ a.e.~in $\Omega$. Therefore, $w\in C_{\bar{u}}$.

\emph{Step 2:} $w\equiv 0$. Using Taylor's theorem, $j'(\bar{u})(u_k-\bar{u})\geq 0$, and the right-hand side estimate in \eqref{eq:estimate_contradiction}, we obtain $j''(\hat{u}_k)w_k^2 < k^{-1}\to 0$ as $k \uparrow \infty$, where $\hat{u}_k:= \bar{u}+\hat{\theta}_k (u_k-\bar{u})$ and $\theta_k \in (0,1).$ We now prove the following result: $j''(\bar{u})w^2\leq \liminf_{k\uparrow \infty} j''(\hat{u}_k)w_k^2$. To do this, we refer to \eqref{eq:identity_j} and write
\begin{equation*}
j''(\hat{u}_k)w_k^2=\alpha\|w_k\|_{L^2(\Omega)}^2-2(w_k\hat{\boldsymbol{\varphi}}_k,\hat{\mathbf{z}}_k)_{\mathbf{L}^2(\Omega)}+\|\hat{\boldsymbol{\varphi}}_k\|_{\mathbf{L}^2(\Omega)}^2.
\end{equation*}
Here, $(\hat{\mathbf{y}}_k,\hat{\mathsf{p}}_k):=\mathcal{S}\hat{u}_k$, $(\hat{\mathbf{z}}_k,\hat{\mathsf{r}}_k)$ denotes the solution to problem \eqref{eq:brinkamn_problem_adjoint}, where $u$ and $\mathbf{y}$ are replaced by $\hat{u}_k$ and $\hat{\mathbf{y}}_k$, respectively, and $(\hat{\boldsymbol{\varphi}}_k,\hat{\zeta}_k)$ solves \eqref{eq:problem_S_first}, where $u$, $v$ and $\mathbf{y}$ are replaced by $\hat{u}_k$, $w_k$ and $\hat{\mathbf{y}}_k$, respectively. We now note that
\begin{multline*}
\|\nabla(\bar{\mathbf{y}}-\hat{\mathbf{y}}_k)\|_{\mathbf{L}^2(\Omega)}+\|\nabla(\bar{\mathbf{z}}-\hat{\mathbf{z}}_k)\|_{\mathbf{L}^2(\Omega)}
\\
\lesssim
(\|\mathbf{f}\|_{\mathbf{H}^{-1}(\Omega)}
+\|\mathbf{y}_{\Omega}\|_{\mathbf{L}^2(\Omega)})
\|\bar{u}-\hat{u}_k\|_{L^2(\Omega)}\to 0,
\quad k\uparrow\infty.
\end{multline*}

We now examine the convergence properties of the sequence $\{ ( \hat{\boldsymbol{\varphi}}_k,\hat{\zeta}_k) \}_{k \geq 1}$. For this purpose, we first write the problem that solves $(\boldsymbol{\varphi}-\hat{\boldsymbol{\varphi}}_k, \zeta-\hat{\zeta}_k)$:
\begin{multline*}
(\nabla (\boldsymbol{\varphi}-\hat{\boldsymbol{\varphi}}_k),\nabla \mathbf{v})_{\mathbf{L}^2(\Omega)}+( \bar{u}(\boldsymbol{\varphi}-\hat{\boldsymbol{\varphi}}_k), \mathbf{v})_{\mathbf{L}^2(\Omega)}-(\zeta-\hat{\zeta}_k,\text{div }\mathbf{v})_{L^2(\Omega)} \\=- ((w-w_k)\bar{\mathbf{y}},\mathbf{v})_{\mathbf{L}^2(\Omega)}+(w_k(\hat{\mathbf{y}}_k-\bar{\mathbf{y}}),\mathbf{v})_{\mathbf{L}^2(\Omega)}+((\hat{u}_k-\bar{u})\hat{\boldsymbol{\varphi}}_k,\mathbf{v})_{\mathbf{L}^2(\Omega)}
=:\text{I}_k+\text{II}_k+\text{III}_k,
\end{multline*}
and $(\mathsf{s},\text{div }(\boldsymbol{\varphi} - \hat{\boldsymbol{\varphi}}_k))_{L^2(\Omega)}  =  0$ for all $\mathbf{v}\in \mathbf{H}_0^1(\Omega)$ and $\mathsf{s}\in L_0^2(\Omega)$, respectively. Since $\bar{\mathbf{y}}\cdot\mathbf{v}\in L^{2}(\Omega)$, the weak convergence $w_k \rightharpoonup w$ in $L^2(\Omega)$ as $k \uparrow\infty$ allows us to obtain that $|\text{I}_k|\to 0$. We control the term $\text{II}_k$ using the Sobolev embedding $\mathbf{H}_0^1(\Omega) \hookrightarrow \mathbf{L}^{4}(\Omega)$:
\begin{equation*}
|\text{II}_k|
\lesssim
\|w_k\|_{L^2(\Omega)}\|\nabla(\bar{\mathbf{y}}-\hat{\mathbf{y}}_k)\|_{\mathbf{L}^{2}(\Omega)}\|\nabla\mathbf{v}\|_{\mathbf{L}^{2}(\Omega)}
\to 0,
\quad k\uparrow\infty.
\end{equation*}
The control of the term $\text{III}_k$ follows similar arguments. In fact, we have
\begin{align*}
|\text{III}_k|\leq
\| \hat{u}_k-\bar{u}\|_{L^{2}(\Omega)}
\|\nabla\hat{\boldsymbol{\varphi}}_k\|_{\mathbf{L}^{2}(\Omega)} \|\nabla\mathbf{v}\|_{\mathbf{L}^{2}(\Omega)}
\lesssim
\| \hat{u}_k-\bar{u}\|_{L^{2}(\Omega)} \|\nabla\mathbf{v}\|_{\mathbf{L}^{2}(\Omega)} \rightarrow 0
\end{align*}
as $k \uparrow \infty$, where we used $\hat{u}_k\to\bar{u}$ in $L^2(\Omega)$ as $k\uparrow\infty$ and $\|\nabla\hat{\boldsymbol{\varphi}}_k\|_{\mathbf{L}^2(\Omega)}\lesssim\|\mathbf{f}\|_{\mathbf{H}^{-1}(\Omega)}$, which is uniform in $k$. We can therefore conclude that $\hat{\boldsymbol{\varphi}}_k\rightharpoonup\boldsymbol{\varphi}$ in $\mathbf{H}_0^1(\Omega)$ as $k \uparrow \infty$ and that $\hat{\boldsymbol{\varphi}}_k \rightarrow \boldsymbol{\varphi}$ in $\mathbf{L}^r(\Omega)$ for every $r < \infty$ when $d = 2$ and for every $r < 6$ when $d = 3$. In particular, $\|\hat{\boldsymbol{\varphi}}_k\|_{\mathbf{L}^2(\Omega)} \to \|\boldsymbol{\varphi}\|_{\mathbf{L}^2(\Omega)}$ as $k\uparrow\infty$.

The convergence of $\{ (w_k\hat{\boldsymbol{\varphi}}_k,\hat{\mathbf{z}}_k)_{\mathbf{L}^2(\Omega)} \}_{k \geq 1}$ to $(w\boldsymbol{\varphi},\bar{\mathbf{z}})_{\mathbf{L}^2(\Omega)}$ in $\mathbb{R}$ is as follows:
\begin{multline*}
\left| (w\boldsymbol{\varphi},\bar{\mathbf{z}})_{\mathbf{L}^2(\Omega)}-(w_k\hat{\boldsymbol{\varphi}}_k,\hat{\mathbf{z}}_k)_{\mathbf{L}^2(\Omega)} \right|
\leq |((w-w_k)\boldsymbol{\varphi},\bar{\mathbf{z}})_{\mathbf{L}^2(\Omega)}|
\\
+\left| (w_k(\boldsymbol{\varphi}-\hat{\boldsymbol{\varphi}}_k),\bar{\mathbf{z}})_{\mathbf{L}^2(\Omega)} \right|
+
\left| (w_k\hat{\boldsymbol{\varphi}}_k,\bar{\mathbf{z}}-\hat{\mathbf{z}}_k)_{\mathbf{L}^2(\Omega)} \right|
\to 0,
\quad k\uparrow\infty,
\end{multline*}
where we used the weak convergence $w_k \rightharpoonup w$ in $L^2(\Omega)$ as $k\uparrow\infty$, the properties $\boldsymbol{\varphi}\cdot \bar{\mathbf{z}}\in L^2(\Omega)$ and $\|w_k\|_{L^2(\Omega)}=1$, the convergence result $\|\boldsymbol{\varphi}-\hat{\boldsymbol{\varphi}}_k \|_{\mathbf{L}^4(\Omega)}\to 0$ as $k\uparrow\infty$, the bound $\|\nabla \bar{\mathbf{z}}\|_{\mathbf{L}^2(\Omega)}\lesssim \|\mathbf{f}\|_{\mathbf{H}^{-1}(\Omega)}+\|\mathbf{y}_{\Omega}\|_{\mathbf{L}^2(\Omega)}$, and $\|\nabla(\bar{\mathbf{z}}-\hat{\mathbf{z}}_k)\|_{\mathbf{L}^{2}(\Omega)}\to 0$ as $k\uparrow\infty$.

A collection of the two previous convergence results, namely $(w_k\hat{\boldsymbol{\varphi}}_k,\hat{\mathbf{z}}_k)_{\mathbf{L}^2(\Omega)} \to (w\boldsymbol{\varphi},\bar{\mathbf{z}})_{\mathbf{L}^2(\Omega)}$ in $\mathbb{R}$ and $\|\hat{\boldsymbol{\varphi}}_k\|_{\mathbf{L}^2(\Omega)} \to \|\boldsymbol{\varphi}\|_{\mathbf{L}^2(\Omega)}$ in $\mathbb{R}$ as $k\uparrow\infty$, in conjunction with the fact that the square of $\| \cdot \|_{L^2(\Omega)}$ is weakly lower semicontinuous in $L^2(\Omega)$, show that $j''(\bar{u})w^2 \leq \liminf_{k\uparrow \infty} j''(\hat{u}_k)w_k^2$. Since $\liminf_{k\uparrow \infty} j''(\hat{u}_k)w_k^2 \leq 0$, we conclude that $j''(\bar{u})w^2 \leq 0$. Finally, we use that $w \in C_{\bar u}$ and that $\bar u$ satisfies the second-order condition $j''(\bar u) v^2 > 0$ for all $v \in C_{\bar u} \setminus \{ 0 \}$ to conclude that $w \equiv 0$.

\emph{Step 3:} \emph{The contradiction}. Since $w \equiv 0$, it is immediate that $\hat{\boldsymbol{\varphi}}_k \rightharpoonup 0$ in $\mathbf{H}_0^1(\Omega)$ as $k\uparrow\infty$. We therefore use the limit $\liminf_{k\uparrow \infty} j''(\hat{u}_k)w_k^2 \leq 0$ and the identities $\alpha = \alpha \|w_k\|_{L^2(\Omega)}^2 = j''(\hat{u}_k)w_k^2 + 2(w_k\hat{\boldsymbol{\varphi}}_k,\hat{\mathbf{z}}_k)_{\mathbf{L}^2(\Omega)}-\|\hat{\boldsymbol{\varphi}}_k\|_{\mathbf{L}^{2}(\Omega)}^2$ to conclude that $\alpha\leq 0$. This is a contradiction and concludes the proof.
\end{proof}

We conclude with an equivalent result that is important for deriving error bounds for the numerical methods that we will propose in our paper. To present such a result, we introduce the set $C_{\bar{u}}^{\tau}:=\{w\in L^2(\Omega) \text{ satisfying } \eqref{eq:cone_condition} \text{ and } w(x)=0 \text{ if } |\bar{\mathfrak{d}}(x)|>\tau\}$.

\begin{theorem}[equivalent optimality conditions]
Let $(\bar{u}, \bar{\mathbf{y}},\bar{\mathbf{z}}) \in \mathbb{U}_{ad} \times \mathbf{H}_0^1(\Omega) \times \mathbf{H}_0^1(\Omega)$ satisfy \eqref{eq:brinkamn_problem_state}, \eqref{eq:brinkamn_problem_adjoint}, and \eqref{eq:variational_inequality}. Then, the following statements are equivalent:
\begin{align}\label{equivalence_condition}
j''(\bar{u})v^2 >0\quad \forall v \in  C_{\bar{u}}\setminus\{0\} \iff \exists \mu,\tau>0: j''(\bar{u})v^2  \geq \mu\|v\|_{L^2{(\Omega)}}^2 \quad \forall v\in C_{\bar{u}}^{\tau}.
\end{align}
\end{theorem}
\begin{proof}
	The result follows from the arguments developed in the proofs of \cite[Theorem 25]{MR3586845} and Theorem \ref{theorem_second_order_optimal}. For the sake of brevity, we omit the details.
\end{proof}
%%%%%%%%%%%%%%%%%%%%%%%%%%%%%%%%%%%%%%%%%%%%%%%%%%%%%%%%%%%%%%%%
%%%%%%%%%%%%%%%%%%%%%%%%%%%%%%%%%%%%%%%%%%%%%%%%%%%%%%%%%%%%%%%%
%%%%%%%%%%%%%%%%%%%%%%%%%%%%%%%%%%%%%%%%%%%%%%%%%%%%%%%%%%%%%%%%

\section{Finite element approximation of the state equations}\label{fem}
In this section, we present a standard finite element method for approximating the solution of the Stokes--Brinkman equations \eqref{eq:brinkman_problem}. Before describing the numerical method, we introduce some terminology and basic ingredients \cite{MR2373954,CiarletBook,Guermond-Ern}.

We denote by $\mathscr{T} = \{ K\}$ a conforming partition or mesh of $\bar{\Omega}$ into closed simplices $K$ with size $h_K = \text{diam}(K)$. Define $h:=\max \{ h_K: K \in \mathscr{T} \}$. We denote by $\mathbb{T}$ a collection of conforming meshes $\mathscr{T}$, which are refinements of an initial mesh $\mathscr{T}_0$. We assume that the collection $\mathbb{T}$ satisfies the so-called shape regularity condition \cite{MR2373954,CiarletBook,Guermond-Ern}. We define $\mathscr{S}$ as the set of interelement boundaries $\gamma$ of $\T$. For $K \in \T$, we define $\mathscr{S}_K$ as the subset of $\mathscr{S}$ that contains the sides of $K$. We denote by $\mathcal{N}_{\gamma}$, for $\gamma \in \mathscr{S}$, the subset of $\T$ that contains the two elements that have $\gamma$ as a side. Finally, we define the following \emph{stars} or \emph{patches} associated with an element $K \in \T$:
\begin{equation}
\label{eq:patch}
\mathcal{N}_K=  \{K' \in \T: \mathscr{S}_K \cap \mathscr{S}_{K'} \neq \emptyset \},
\qquad 
\mathcal{N}_K^*= \{K' \in \T: K \cap {K'} \neq \emptyset \}.
\end{equation}
In an abuse of notation, we use $\mathcal{N}_K$ and $\mathcal{N}_K^*$ in the following to denote either the sets themselves or the union of their elements.

To perform an a posteriori error analysis, we introduce the following notation. If $K^{+}, K^{-} \in \T$ are such that $K^{+} \neq K^{-}$ and $\partial K^{+} \cap \partial K^{-} = \gamma$, we define the \emph{jump} of the discrete tensor-valued function $\mathbf{w}_{h}$ on $\gamma$ by $\llbracket \mathbf{w}_{h} \cdot \mathbf{n} \rrbracket:= \mathbf{w}_{h} \cdot \mathbf{n}^{+} |^{}_{K^{+}} + \mathbf{w}_{h} \cdot \mathbf{n}^{-}  |^{}_{K^{-}}$, where $\mathbf{n}^{+}$ and $\mathbf{n}^{-}$ denote the unit normals on $\gamma$ pointing towards $K^{+}$ and $K^{-}$, respectively.
% Let $\mathbf{w}_{h}$ be a discrete tensor-valued function and let $\gamma \in \mathscr{S}$ be an internal side. We define the \emph{jump} or \emph{interelement residual} of $\mathbf{w}_{h}$ on $\gamma$ by
% $\llbracket \mathbf{w}_{h} \cdot \mathbf{n} \rrbracket:= \mathbf{w}_{h} \cdot \mathbf{n}^{+} |^{}_{K^{+}} + \mathbf{w}_{h} \cdot \mathbf{n}^{-}  |^{}_{K^{-}},$
% where $\mathbf{n}^{+}$ and $\mathbf{n}^{-}$ denote the unit normals on $\gamma$ pointing towards $K^{+}$ and $K^{-}$, respectively. Here, $K^{+}$, $K^{-} \in \T$ are such that $K^{+} \neq K^{-}$ and $\partial K^{+} \cap \partial K^{-} = \gamma$.


\subsection{Finite element spaces}
Let $\mathscr{T}$ be a mesh in $\mathbb{T}$. To approximate the velocity field $\boldsymbol{\mathsf{y}}$ and the pressure $\mathfrak{p}$ of the fluid, we consider a pair $(\mathbf{X}_h,M_h)$ of finite element spaces satisfying a uniform discrete inf-sup condition: There exists a constant $\tilde{\beta}>0$ independent of $h$ such that
\begin{equation}\label{eq:discrete_infsup}
\inf _{\mathsf{q}_h \in M_h} \sup _{\mathbf{v}_h \in \mathbf{X}_h} \frac{ \int_{\Omega} \mathsf{q}_h\operatorname{div} \mathbf{v}_h \mathrm{d}x }{\left\|\nabla\mathbf{v}_h\right\|_{\mathbf{L}^2(\Omega)}\left\|\mathsf{q}_h\right\|_{L^2(\Omega)}} \geqslant \tilde{\beta}>0.
\end{equation}
We will look in particular at the following pairs, which are significant in the literature:

\begin{enumerate}
\item[(a)] The \emph{mini element} \cite[Section 4.2.4]{Guermond-Ern}: In this case,
\begin{eqnarray}\label{ME:vel_space}
\mathbf{X}_h&=&\{\mathbf{v}_{h}\in\mathbf{C}(\bar{\Omega}): \mathbf{v}_{h}|_{K}\in[\mathbb{W}(K)]^{d} \,\,\forall K\in\T \} \cap \mathbf{H}_0^1(\Omega),
\\ 
\label{ME:press_space}
M_h&=&\{\mathsf{q}_{h}\in L_0^2(\Omega)\cap C(\bar{\Omega}): \mathsf{q}_{h}|_{K}\in\mathbb{P}_1(K) \,\,\forall K\in\T \},
\end{eqnarray}
where $\mathbb{W}(K):=\mathbb{P}_1(K)\oplus \mathbb{B}(K)$, and $\mathbb{B}(K)$ denotes the space spanned by a local bubble function.

\item[(b)] The lowest order \emph{Taylor--Hood element} \cite[Section 4.2.5]{Guermond-Ern}: In this scenario,
\begin{eqnarray}\label{TH:vel_space}
\mathbf{X}_h
&=&
\{\mathbf{v}_{h}\in\mathbf{C}(\bar{\Omega}): \mathbf{v}_{h}|_{K}\in[\mathbb{P}_2(K)]^{d} \,\,\forall K\in\T \} \cap \mathbf{H}_0^1(\Omega),\\ \label{TH:press_space}
M_h
&=&
\{\mathsf{q}_{h}\in L_0^2(\Omega)\cap C(\bar{\Omega}): \mathsf{q}_{h}|_{K}\in\mathbb{P}_1(K) \,\,\forall K\in\T \}.
\end{eqnarray}
\end{enumerate}

A proof of the inf-sup condition \eqref{eq:discrete_infsup} for the mini element can be found in \cite[Lemma 4.20]{Guermond-Ern}. Provided that the mesh $\T_h$ contains at least three triangles in two dimensions and that each tetrahedron has at least one internal vertex in three dimensions, a proof of \eqref{eq:discrete_infsup} for the Taylor--Hood element can be found in \cite[Theorem 8.8.1]{MR3097958} and \cite[Theorem 8.8.2]{MR3097958}, respectively.

\subsection{The finite element method}
We introduce the following approximation of problem \eqref{eq:brinkman_problem}: Given $\boldsymbol{\mathfrak{f}} \in \mathbf{H}^{-1}(\Omega)$ and $\mathfrak{u} \in \mathbb{U}_{ad}$, find $(\boldsymbol{\mathsf{y}}_{h},\mathfrak{p}_{h})\in \mathbf{X}_h\times M_h$ such that
\begin{equation}\label{eq:modeldiscrete}
\begin{array}{rcl}
(\nabla \boldsymbol{\mathsf{y}}_{h},\nabla \mathbf{v}_{h})_{\mathbf{L}^2(\Omega)}\!+\!( \mathfrak{u}\boldsymbol{\mathsf{y}}_{h}, \mathbf{v}_{h})_{\mathbf{L}^2(\Omega)}\!-\!(\mathfrak{p}_{h},\text{div }\mathbf{v}_{h})_{L^2(\Omega)}& = & \langle \boldsymbol{\mathfrak{f}},\mathbf{v}_{h} \rangle,\\
(\mathsf{q}_{h},\text{div }\boldsymbol{\mathsf{y}}_{h})_{L^2(\Omega)} & = & 0
\end{array}
\end{equation}
for all $\mathbf{v}_{h}\in \mathbf{X}_h$ and $\mathsf{q}_{h}\in M_h$, respectively. The existence and uniqueness of $(\boldsymbol{\mathsf{y}}_{h},\mathfrak{p}_{h})$ $\in \mathbf{X}_h\times M_h$ is standard. In addition, we have $\|\nabla \boldsymbol{\mathsf{y}}_{h}\|_{\mathbf{L}^2(\Omega)} + \|\mathfrak{p}_{h}\|_{L^2(\Omega)}\lesssim \| \boldsymbol{\mathfrak{f}}\|_{\mathbf{H}^{-1}(\Omega)}$.

In the following, we present a priori error bounds for the approximation \eqref{eq:modeldiscrete}.

\begin{theorem}[error estimates]\label{eq:theorem_error_estimates}
Let $\Omega\subset \mathbb{R}^d$ be a convex polytope, let $\boldsymbol{\mathfrak{f}}\in\mathbf{L}^2(\Omega)$, and let $\mathfrak{u}\in \mathbb{U}_{ad}$. Let $(\boldsymbol{\mathsf{y}},\mathfrak{p})\in\mathbf{H}^2(\Omega)\cap\mathbf{H}_0^1(\Omega) \times H^1(\Omega)\cap L_0^2(\Omega)$ be the solution to the problem \eqref{eq:brinkman_problem}, and let $(\boldsymbol{\mathsf{y}}_{h},\mathfrak{p}_{h})\in\mathbf{X}_h\times  M_h$ be its finite element approximation obtained as the solution to \eqref{eq:modeldiscrete}. Then, we have the following error estimate:
\begin{align}\label{eq:error_estimate_Fem}
\|\nabla(\boldsymbol{\mathsf{y}}-\boldsymbol{\mathsf{y}}_{h})\|_{\mathbf{L}^2(\Omega)}+\|\mathfrak{p}-\mathfrak{p}_{h}\|_{L^2(\Omega)}
\lesssim 
h
\|\boldsymbol{\mathfrak{f}}\|_{\mathbf{L}^2(\Omega)},
\end{align}
where the hidden constant is independent of $(\boldsymbol{\mathsf{y}},\mathfrak{p})$, $(\boldsymbol{\mathsf{y}}_{h},\mathfrak{p}_{h})$, and $h$.
\end{theorem}
\begin{proof}
A quasi-best approximation result follows from a direct application of \cite[Lemma 2.44]{Guermond-Ern}. With this result in hand, the error bound \eqref{eq:error_estimate_Fem} follows from the regularity results of Theorem \ref{eq:theorem_regularity} and standard interpolation error estimates. For the sake of brevity, we omit the details of the proof.
\end{proof}

\begin{remark}[Taylor--Hood approximation]
If the solution $(\mathbf{y},\mathsf{p})$ of problem \eqref{eq:brinkman_problem} belongs to $\mathbf{H}^3(\Omega) \cap\mathbf{H}_0^1(\Omega)\times H^2(\Omega)\cap L_0^2(\Omega)$, then we have the following error bound:
\begin{align*}
\|\nabla(\boldsymbol{\mathsf{y}}-\boldsymbol{\mathsf{y}}_{h})\|_{\mathbf{L}^2(\Omega)}+\|\mathfrak{p}-\mathfrak{p}_{h}\|_{L^2(\Omega)}
\lesssim 
h^2 \left( \|\boldsymbol{\mathsf{y}}\|_{\mathbf{H}^3(\Omega)} + \|\mathfrak{p}\|_{H^2(\Omega)} \right).
\end{align*}
\end{remark}

\subsection{A posteriori error estimates}

Following \cite{MR1445736,MR1259620}, we present an a posteriori error estimator for the approximation \eqref{eq:modeldiscrete} of problem \eqref{eq:brinkman_problem}: For $K\in \T$ and an internal side $\gamma\in\mathscr{S}$, we define the \emph{element residual} $\mathcal{R}_K$ and the \emph{interelement residual} $\mathcal{J}_{\gamma}$ as
\begin{align}\label{eq:residuals}
\mathcal{R}_K:= (\boldsymbol{\mathfrak{f}}+\Delta \boldsymbol{\mathsf{y}}_{h}\!-\mathfrak{u}\boldsymbol{\mathsf{y}}_{h}\!-\!\nabla\mathfrak{p}_{h})|_K,
\qquad 
\mathcal{J}_{\gamma}:= \llbracket(\nabla \boldsymbol{\mathsf{y}}_{h}-\mathfrak{p}_{h}\mathbf{I})\cdot \mathbf{n}\rrbracket,
\end{align}
respectively. Here, $(\boldsymbol{\mathsf{y}}_{h},\mathfrak{p}_{h})$ denotes the solution to the discrete problem \eqref{eq:modeldiscrete} and $\mathbf{I}$ denotes the identity matrix in $\mathbb{R}^{d\times d}$. We thus define the \emph{local error indicators} $\mathcal{E}_{K}$ and the corresponding a posteriori error estimator $\mathcal{E}_{h}$ as follows:
\begin{align}\label{eq:indicator_estimator_e}
\mathcal{E}_{K}^2:=h_K^2  \|\mathcal{R}_K\|_{\mathbf{L}^2(K)}^2
+h_K  \|\mathcal{J}_{\gamma} \|_{\mathbf{L}^2(\partial K\setminus\partial \Omega)}^2+\|\text{div }\boldsymbol{\mathsf{y}}_{h}\|_{L^2(K)}^2,
\quad 
\mathcal{E}_{h}^2:= \sum_{K\in\T}\mathcal{E}_{K}^2.
\end{align}

We present the following global reliability result.

\begin{theorem}[global reliability]
\label{eq:reliability} 
Let $\Omega \subset \mathbb{R}^d$ be a Lipschitz domain, let $\boldsymbol{\mathfrak{f}}\in\mathbf{L}^2(\Omega)$, and let $\mathfrak{u}\in\mathbb{U}_{ad}$. Let $(\boldsymbol{\mathsf{y}},\mathfrak{p})\in\mathbf{H}_0^1(\Omega)\times L_0^2(\Omega)$ be the solution to \eqref{eq:brinkman_problem}, and let $(\boldsymbol{\mathsf{y}}_{h},\mathfrak{p}_{h})\in \mathbf{X}_h\times M_h$ be its finite element approximation, which solves \eqref{eq:modeldiscrete}. Then, 
\begin{eqnarray}\label{eq:bound_realiability}
\|\nabla(\boldsymbol{\mathsf{y}}-\boldsymbol{\mathsf{y}}_{h})\|_{\mathbf{L}^2(\Omega)}^2+\|\mathfrak{p}-\mathfrak{p}_{h}\|_{L^2(\Omega)}^2\lesssim \mathcal{E}_{h}^2,
\end{eqnarray}
where the hidden constant is independent of $(\boldsymbol{\mathsf{y}},\mathfrak{p})$ and $(\boldsymbol{\mathsf{y}}_{h},\mathfrak{p}_{h})$, the size of the elements in the mesh $\T$, and $\# \T$.
\end{theorem}
\begin{proof}
The proof follows the same arguments as in \cite[Section 5]{MR1259620} and \cite[Section 3]{MR1445736}. For the sake of brevity, we omit the details.
\end{proof}
%LO QUE ESTA COMENTADO ES LA DEMOSTRACION
%\begin{proof} 
%We divide the proof into three steps.
%
%\emph{Step 1.} First, we derive a reliability analysis for the a posteriori error estimator presented in \eqref{eq:indicator_estimator_e}. We are inspired by \cite[section 3]{MR1445736}. To do this, we introduce the Ritz projection $(\mathbf{\Phi},\psi)$ of the residuals. We note that the pair $(\mathbf{\Phi},\psi)$ is defined as the solution to the following problem: Find $(\mathbf{\Phi},\psi)\in \mathbf{H}_0^1(\Omega)\times L_0^2(\Omega)$ such that
%\begin{equation}\label{eq:ritz}
%\begin{array}{rcl}
%(\nabla\mathbf{\Phi},\nabla \mathbf{v})_{\mathbf{L}^2(\Omega)}&=&(\nabla\mathbf{e}_{\boldsymbol{\mathsf{y}}},\nabla\mathbf{v})_{\mathbf{L}^2(\Omega)}+(\mathfrak{u}\mathbf{e}_{\boldsymbol{\mathsf{y}}},\mathbf{v})_{\mathbf{L}^2(\Omega)}-(e_{\mathfrak{p}},\text{div }\mathbf{v})_{L^2(\Omega)},\\
%(\psi,\mathsf{q})_{L^2(\Omega)}&=&(\mathsf{q},\text{div }\mathbf{e}_{\boldsymbol{\mathsf{y}}})_{L^2(\Omega)},
%\end{array}
%\end{equation}
%for all $(\mathbf{v},\mathsf{q})\in\mathbf{H}_0^1(\Omega)\times L_0^2(\Omega)$. Here, we have introduced the notation $(\mathbf{e}_{\boldsymbol{\mathsf{y}}},e_{\mathfrak{p}}):=(\boldsymbol{\mathsf{y}}-\boldsymbol{\mathsf{y}}_{h},\mathfrak{p}-\mathfrak{p}_{h})$. The existence and uniqueness of the pair $(\mathbf{\Phi},\psi)\in \mathbf{H}_0^1(\Omega)\times L_0^2(\Omega)$ follows from a simple application of the Lax--Milgram lemma. In addition, with similar arguments to those developed in \cite[Theorem 4]{MR1259620} yield the estimate
%\begin{align}\label{eq:ritz_estimate}
%\|\nabla\mathbf{e}_{\boldsymbol{\mathsf{y}}}\|_{\mathbf{L}^2(\Omega)}^2+\|e_{\mathfrak{p}}\|_{L^2(\Omega)}^2\lesssim \|\nabla\mathbf{\Phi}\|_{\mathbf{L}^2(\Omega)}^2+\|\psi\|_{L^2(\Omega)}^2.
%\end{align}
%
%\emph{Step 2.} Second, we dedicated to controlling the terms on the right-hand side in the inequality \eqref{eq:ritz_estimate}. Let $\mathbf{v}\in \mathbf{H}_0^1(\Omega)$ and set $\mathsf{q}=0$ in \eqref{eq:ritz}. This fact, combined with \eqref{eq:brinkman_problem}, yields the identity
%\begin{align*}
%(\nabla\mathbf{\Phi},\nabla \mathbf{v})_{\mathbf{L}^2(\Omega)}=(\boldsymbol{\mathfrak{f}},\mathbf{v})_{\mathbf{L}^2(\Omega)}-(\nabla \boldsymbol{\mathsf{y}}_{h},\nabla \mathbf{v})_{\mathbf{L}^2(\Omega)}\!-\!(\mathfrak{u}\boldsymbol{\mathsf{y}}_{h}, \mathbf{v})_{\mathbf{L}^2(\Omega)}\!+\!(\mathfrak{p}_{h},\text{div }\mathbf{v})_{L^2(\Omega)}.
%\end{align*}
%We use $\mathcal{C}_{h}:\mathbf{L}^1(\Omega)\to \mathbf{V}(\Omega)$ to set the Cl\'ement interpolation operator \cite{MR2373954,CiarletBook}. If we invoke the previous relation, the discrete problem \eqref{eq:modeldiscrete} with $\mathbf{v}_{h}=\mathcal{C}_{h}\mathbf{v}$, an elementwise integration by parts formula, approximation properties for $\mathcal{C}_{h}$, and the finite overlapping property of stars, we can arrive at $(\nabla\mathbf{\Phi},\nabla \mathbf{v})_{\mathbf{L}^2(\Omega)}\lesssim \mathcal{E}_{h}\|\nabla \mathbf{v}\|_{\mathbf{L}^2(\Omega)}$.
%Replacing the test function $\mathbf{v}=\mathbf{\Phi}$ in the above estimate, we immediately obtain the bound $\|\nabla\mathbf{\Phi}\|_{\mathbf{L}^2(\Omega)}\lesssim\mathcal{E}_{h}$. Now, let $\mathsf{q}\in L_0^2(\Omega)$ and set $\mathbf{v}=\mathbf{0}$ in \eqref{eq:ritz}. Applying the Cauchy--Schwarz inequality, we obtain
%$(\psi,\mathsf{q})_{L^2(\Omega)}^2\leq \sum_{K\in\T} \|\text{div }\boldsymbol{\mathsf{y}}_{h}\|_{L^2(K)}^2\|\mathsf{q}\|_{L^2(\Omega)}^2.
%$
%As a consequence, $\|\psi\|_{L^2(\Omega)}\lesssim \mathcal{E}_{h}$. Therefore, by collecting the above estimates, we can derive the estimate 
%\begin{align}\label{eq:ritz_estimate_02}
%\|\nabla\mathbf{\Phi}\|_{\mathbf{L}^2(\Omega)}^2+\|\psi\|_{L^2(\Omega)}^2\lesssim\mathcal{E}_{h}^2.
%\end{align}
%
%\emph{Step 3.} To conclude, we invoke \eqref{eq:ritz_estimate} and \eqref{eq:ritz_estimate_02} to arrive at the desired estimate.
%\end{proof}


\section{Finite element discretizations for the control problem}
We propose two discretization methods to approximate solutions of the control problem \eqref{eq:min_functional}--\eqref{eq:brinkamn_problem_state}: a semidiscrete method in which the control set is not discretized, and a fully discrete one in which the admissible control set is discretized with piecewise constant functions.

\subsection{The fully discrete scheme}
\label{fully_discrete_framework}
To present the fully discrete scheme, we introduce
$
\mathbb{U}_h:=\{u_{h}\in L^{\infty}(\Omega):u_{h}|_{K}\in\mathbb{P}_0(K)~\forall K\in\T\}
$
and $\mathbb{U}_{ad,h}:= \mathbb{U}_h\cap \mathbb{U}_{ad}$. The scheme is as follows: Find $\min J(\mathbf{y}_{h},u_{h})$ subject to the \emph{discrete state equations}
\begin{equation}\label{eq:brinkamn_problem_state_fd}
\begin{array}{rl}
(\nabla \mathbf{y}_{h},\nabla \mathbf{v}_{h})_{\mathbf{L}^2(\Omega)}
+
( u_{h}\mathbf{y}_{h}, \mathbf{v}_{h})_{\mathbf{L}^2(\Omega)}
-
(\mathsf{p}_{h},\text{div }\mathbf{v}_{h})_{L^2(\Omega)}
=& \langle \mathbf{f},\mathbf{v}_{h} \rangle,\\
(\mathsf{q}_{h},\text{div }\mathbf{y}_{h})_{L^2(\Omega)} =&  0,
\end{array}
\end{equation}
for all $(\mathbf{v}_{h},\mathsf{q}_{h}) \in \mathbf{X}_h \times M_h$, and the \emph{discrete control constraints} $u_{h}\in\mathbb{U}_{ad,h}$.

The scheme admits at least one solution. To present first-order optimality conditions, we introduce the discrete control-state map $\mathcal{G}_h: \mathbb{U}_{ad} \ni u \mapsto \mathbf{y}_h \in \mathbf{X}_h$, where $\mathbf{y}_h$ corresponds to the velocity component of the pair $(\mathbf{y}_h,\mathsf{p}_h)$ that solves problem \eqref{eq:modeldiscrete}. We also introduce $j_h(u):=J(\mathcal{G}_h u,u)$. The first-order optimality conditions are thus as follows: If $\bar{u}_{h}$ is locally optimal for the fully discrete scheme, then
\begin{eqnarray}\label{eq:variational_inequality_discrete}
j'_h(\bar{u}_h)(u_h-\bar{u}_h)=(\alpha \bar{u}_{h}-\bar{\mathbf{y}}_{h}\cdot \bar{\mathbf{z}}_{h},u_{h}-\bar{u}_{h})_{L^2(\Omega)}\geq 0\quad\forall u_{h}\in \mathbb{U}_{ad,h}.
\end{eqnarray}
Here, $\bar{\mathbf{y}}_h = \mathcal{G}_h \bar{\mathbf{u}}_{h}$ and the \emph{discrete adjoint state} $(\bar{\mathbf{z}}_{h},\bar{\mathsf{r}}_{h})$ corresponds to the solution of the \emph{discrete adjoint equations}: Find $(\bar{\mathbf{z}}_{h},\bar{\mathsf{r}}_{h})\in \mathbf{X}_h\times M_h$ such that
\begin{multline}
 \label{eq:brinkman_problem_adjoint_fully}
(\nabla \mathbf{v}_{h},\nabla \bar{\mathbf{z}}_{h})_{\mathbf{L}^2(\Omega)}
+
(\bar{u}_{h}\bar{\mathbf{z}}_{h}, \mathbf{v}_{h})_{\mathbf{L}^2(\Omega)}+
(\bar{\mathsf{r}}_{h},\text{div }\mathbf{v}_{h})_{L^2(\Omega)}
= (\bar{\mathbf{y}}_{h}
-
\mathbf{y}_{\Omega},\mathbf{v}_{h})_{\mathbf{L}^2(\Omega)},
\\
(\mathsf{s}_{h},\text{div }\bar{\mathbf{z}}_{h})_{L^2(\Omega)} =  0 \qquad \forall (\mathbf{v}_{h},\mathsf{s}_{h})\in \mathbf{X}_h\times M_h.
\end{multline}

\subsection{The semidiscrete scheme}
\label{semi_discrete_framework}
\label{sectionsemi-discrete}
In this scheme, only the state space $\mathbf{H}_0^1(\Omega) \times L_0^2(\Omega)$ is discretized; $\mathbb{U}_{ad}$ is not discretized. The semidiscrete scheme reads as follows: Find $\min J(\mathbf{y}_{h},\mathsf{u})$ subject to the \emph{discrete state equations}
\begin{multline}\label{eq:brinkamn_problem_state_sd}
(\nabla \mathbf{y}_{h},\nabla \mathbf{v}_{h})_{\mathbf{L}^2(\Omega)}
+
( \mathsf{u}\mathbf{y}_{h}, \mathbf{v}_{h})_{\mathbf{L}^2(\Omega)}
-
(\mathsf{p}_{h},\text{div }\mathbf{v}_{h})_{L^2(\Omega)} =  \langle \mathbf{f},\mathbf{v}_{h} \rangle,
\\
(\mathsf{q}_{h},\text{div }\mathbf{y}_{h})_{L^2(\Omega)} =  0 \qquad \forall (\mathbf{v}_{h},\mathsf{q}_{h}) \in \mathbf{X}_h \times M_h,
\end{multline}
and the \emph{control constraints} $\mathsf{u}\in\mathbb{U}_{ad}$. The existence of an optimal solution follows from standard arguments. Moreover, if $\bar{\mathsf{u}}\in\mathbb{U}_{ad}$ denotes a locally optimal control, then
\begin{eqnarray}\label{eq:variational_inequality_semi-discrete}
j'_h(\bar{\mathsf{u}})(u-\bar{\mathsf{u}})=(\alpha \bar{\mathsf{u}}-\bar{\mathbf{y}}_{h}\cdot \bar{\mathbf{z}}_{h},u-\bar{\mathsf{u}})_{L^2(\Omega)}\geq 0\quad\forall u\in \mathbb{U}_{ad}.
\end{eqnarray}
Here, $\bar{\mathbf{y}}_h = \mathcal{G}_h \bar{\mathsf{u}}$ and $(\bar{\mathbf{z}}_{h},\bar{\mathsf{r}}_{h})\in \mathbf{X}_h\times M_h $ is the solution to the following problem:
\begin{multline}\label{eq:brinkman_problem_adjoint_semi}
(\nabla \mathbf{v}_{h},\nabla \bar{\mathbf{z}}_{h})_{\mathbf{L}^2(\Omega)}
+
(\bar{\mathsf{u}}\bar{\mathbf{z}}_{h}, \mathbf{v}_{h})_{\mathbf{L}^2(\Omega)}
+
(\bar{\mathsf{r}}_{h},\text{div }\mathbf{v}_{h})_{L^2(\Omega)} =  (\bar{\mathbf{y}}_{h}-\mathbf{y}_{\Omega},\mathbf{v}_{h})_{\mathbf{L}^2(\Omega)},
\\
(\mathsf{s}_{h},\text{div }\bar{\mathbf{z}}_{h})_{L^2(\Omega)}  =  0 \qquad \forall (\mathbf{v}_{h},\mathsf{s}_{h})\in \mathbf{X}_h\times M_h.
\end{multline}

\subsection{Convergence of discretizations}

In this section, we analyze convergence properties for suitable discretizations of the state and adjoint equations and the fully and semidiscrete methods. To simplify the presentation of the material, we define
\begin{equation}
\label{eq:energy_norm}
 \| (\mathbf{v},\mathsf{q}) \|_{\Omega} := \|\nabla \mathbf{v} \|_{\mathbf{L}^2(\Omega)}+\|\mathsf{q}\|_{L^2(\Omega)},
 \quad
 (\mathbf{v},\mathsf{q}) \in \mathbf{H}_0^1(\Omega) \times L^2(\Omega).
\end{equation}

\subsubsection{Discretization of the state equations}
We have the following result.

\begin{proposition}[convergence and error bound]\label{theorem_convergence_01}
Let $\Omega \subset \mathbb{R}^d$ be a Lipschitz polytope, and let $\mathbf{f} \in \mathbf{H}^{-1}(\Omega)$. Let $(\mathbf{y},\mathsf{p})$ and $(\mathbf{y}_{h},\mathsf{p}_{h})$ be the unique solutions of problems \eqref{eq:brinkamn_problem_state} and \eqref{eq:brinkamn_problem_state_fd}, respectively. Then, we have the following result:
\begin{align}\label{eq:convergence_prop}
u_{h}\rightharpoonup u
\textrm{ in }
L^2(\Omega) 
\textrm{ as } h \rightarrow 0
\Rightarrow 
\mathbf{y}_{h} \to \mathbf{y}
\textrm{ in }
\mathbf{H}_0^1(\Omega),
\,
\mathsf{p}_h \rightharpoonup \mathsf{p}
\textrm{ in } L^2(\Omega) \textrm{ as } h \rightarrow 0.
\end{align}
If, in addition, $\Omega$ is convex and $\mathbf{f} \in \mathbf{L}^{2}(\Omega)$, then we have the error bound
\begin{align}\label{eq:errorbound}
\|\nabla(\mathbf{y}-\mathbf{y}_{h})\|_{\mathbf{L}^2(\Omega)}+\|\mathsf{p}-\mathsf{p}_{h}\|_{L^2(\Omega)}\lesssim h \|\mathbf{f}\|_{\mathbf{L}^2(\Omega)}+\|u-u_{h}\|_{L^2(\Omega)}.
\end{align}
\end{proposition}
\begin{proof}

We divide the proof into two steps.

\emph{Step 1:} \emph{The property \eqref{eq:convergence_prop}.} Let $(\mathbf{y}_{h}(u),\mathsf{p}_h(u))\in \mathbf{X}_h \times M_h$ be the solution to \eqref{eq:brinkamn_problem_state_fd}, where $u_{h}$ is replaced by $u$. A quasi-best approximation property, which follows from \cite[Lemma 2.44]{Guermond-Ern}, combined with a density argument, as the one developed in \cite[Corollary 1.109]{Guermond-Ern}, show that
$
\|(\mathbf{y}-\mathbf{y}_{h}(u),\mathsf{p}-\mathsf{p}_h(u))\|_{\Omega}
\rightarrow 0
$
as 
$h \rightarrow 0$. To control $\|(\mathbf{y}_{h}(u)-\mathbf{y}_{h},\mathsf{p}_h(u)-\mathsf{p}_h)\|_{\Omega}$, we note that $(\mathbf{y}_{h}(u)-\mathbf{y}_{h},\mathsf{p}_{h}(u)-\mathsf{p}_{h}) \in \mathbf{X}_h \times M_h$ solves
\begin{multline*}
(\nabla (\mathbf{y}_h(u)-\mathbf{y}_h),\nabla \mathbf{v}_h)_{\mathbf{L}^2(\Omega)}
+
( u_h(\mathbf{y}_h(u)-\mathbf{y}_h), \mathbf{v}_h)_{\mathbf{L}^2(\Omega)}
-
(\mathsf{p}_h(u)-\mathsf{p}_h,\text{div }\mathbf{v}_h)_{L^2(\Omega)} \\=((u_{h}-u)\mathbf{y}_h(u),\mathbf{v}_h)_{\mathbf{L}^2(\Omega)},\quad (\mathsf{q}_h,\text{div}(\mathbf{y}_h(u)-\mathbf{y}_h))_{L^2(\Omega)}  =  0
\end{multline*}
for all $(\mathbf{v}_{h},\mathsf{q}_{h}) \in \mathbf{X}_h  \times M_h$. If we set $\mathbf{v}_{h}=\mathbf{y}_h(u)-\mathbf{y}_h$ and $\mathsf{q}_{h}=0$, we obtain $\|\nabla(\mathbf{y}_{h}(u)-\mathbf{y}_{h})\|_{\mathbf{L}^2(\Omega)} \lesssim \| (u_{h}-u)\mathbf{y}_h(u) \|_{\mathbf{H}^{-1}(\Omega)}$. Given the convergence property $ \| \nabla (\mathbf{y}-\mathbf{y}_{h}(u))\|_{\mathbf{L}^2(\Omega)} \rightarrow 0 $ as $h \rightarrow 0$, we conclude that the sequence $\{ \mathbf{y}_h (\mathbf{u}) \cdot \mathbf{v} \}_{h>0}$ converges in $L^2(\Omega)$ as $h \rightarrow 0$, for a given $\mathbf{v} \in \mathbf{H}_0^1(\Omega)$. From this follows
\begin{align*}
\|\nabla(\mathbf{y}_{h}(u)-\mathbf{y}_{h})\|_{\mathbf{L}^2(\Omega)} \lesssim \| (u_{h}-u)\mathbf{y}_h(u) \|_{\mathbf{H}^{-1}(\Omega)} \rightarrow 0,
\quad h \rightarrow 0.
\end{align*}
Finally, the weak convergence  $\mathsf{p}_h(u)-\mathsf{p}_h \rightharpoonup 0$ in $L^2(\Omega)$ as $h \rightarrow 0$ follows from using the discrete problem that $(\mathbf{y}_{h}(u)-\mathbf{y}_{h},\mathsf{p}_{h}(u)-\mathsf{p}_{h}) \in \mathbf{X}_h \times M_h$ solves 
and the convergence property $ \| \nabla (\mathbf{y}-\mathbf{y}_{h}(u))\|_{\mathbf{L}^2(\Omega)} \rightarrow 0 $ as 
$h \rightarrow 0$.

\emph{Step 2:} \emph{The bound \eqref{eq:errorbound}}. Let $(\mathbf{y}(u_{h}),\mathsf{p}(u_{h}))\in\mathbf{H}_0^1(\Omega)\times L_0^2(\Omega)$ be the solution \eqref{eq:brinkamn_problem_state}, where $u$ is replaced by $u_{h}$. The control of $\| (\mathbf{y}(u_{h})-\mathbf{y}_{h}, \mathsf{p}(u_h)-\mathsf{p}_h) \|_{\Omega}$ follows directly from Theorem \ref{eq:theorem_error_estimates}:
$
\| (\mathbf{y}(u_{h})-\mathbf{y}_{h}, \mathsf{p}(u_h)-\mathsf{p}_h ) \|_{\Omega} \lesssim h\|\mathbf{f}\|_{\mathbf{L}^2(\Omega)}.
$
In the following, we bound $\|\nabla(\mathbf{y}-\mathbf{y}(u_{h}))\|_{\mathbf{L}^2(\Omega)}$. To this end, we note that $(\mathbf{y}-\mathbf{y}(u_h),\mathsf{p}-\mathsf{p}(u_h))$
solves
\begin{multline*}
(\nabla (\mathbf{y}-\mathbf{y}(u_{h})),\nabla \mathbf{v})_{\mathbf{L}^2(\Omega)}+( u(\mathbf{y}-\mathbf{y}(u_{h})), \mathbf{v})_{\mathbf{L}^2(\Omega)}-(\mathsf{p}-\mathsf{p}(u_{h}),\text{div }\mathbf{v})_{L^2(\Omega)} \\=((u_{h}-u)\mathbf{y}(u_{h}),\mathbf{v})_{\mathbf{L}^2(\Omega)},
\qquad 
(\mathsf{q},\text{div }(\mathbf{y}-\mathbf{y}(u_{h})))_{L^2(\Omega)}  =  0
\end{multline*}
for all $(\mathbf{v},\mathsf{q}) \in \mathbf{H}_0^1(\Omega) \times L_0^2(\Omega)$. If we set $\mathbf{v} = \mathbf{y}-\mathbf{y}(u_{h})$ and $\mathsf{q}=0$, we obtain $\|\nabla(\mathbf{y}-\mathbf{y}(u_{h}))\|_{\mathbf{L}^2(\Omega)}\lesssim \|(u-u_{h})\mathbf{y}(u_{h})\|_{\mathbf{H}^{-1}(\Omega)} \lesssim \|u-u_{h}\|_{L^{2}(\Omega)}$, and as a result
\begin{align}\label{eq:error_aux_est02}
\|\nabla(\mathbf{y}-\mathbf{y}_{h})\|_{\mathbf{L}^2(\Omega)}\lesssim h \|\mathbf{f}\|_{\mathbf{L}^2(\Omega)}+\|u-u_{h}\|_{L^2(\Omega)}.
\end{align}
To bound $\|\mathsf{p}-\mathsf{p}(u_h)\|_{L^2(\Omega)}$, we use the inf--sup condition \eqref{eq:infsup} and the problem that $(\mathbf{y}-\mathbf{y}(u_h),\mathsf{p}-\mathsf{p}(u_h))\in\mathbf{H}_0^1(\Omega)\times L_0^2(\Omega)$ solves. In fact, we have
\begin{align}\label{eq:error_aux_est03}
\|\mathsf{p}-\mathsf{p}(u_h)\|_{L^2(\Omega)}\lesssim \|\nabla(\mathbf{y}-\mathbf{y}(u_h))\|_{\mathbf{L}^2(\Omega)}+\|u-u_h\|_{L^2(\Omega)}
\lesssim
\|u-u_{h}\|_{L^2(\Omega)},
\end{align}
where we have used the bound $\|\nabla(\mathbf{y}-\mathbf{y}(u_{h}))\|_{\mathbf{L}^2(\Omega)} \lesssim \|u-u_{h}\|_{L^{2}(\Omega)}$ in the last step.
\end{proof}

\subsubsection{Auxiliary results: discretization of the adjoints equations}
Before deriving error bounds, we introduce auxiliary variables that are of particular importance for our analysis. Let $(\boldsymbol{\mathfrak{z}},\mathfrak{r})\in \mathbf{H}_0^1(\Omega)\times L_0^2(\Omega)$ be the solution of
\begin{multline}\label{eq:aux_problem_adj_1}
(\nabla \mathbf{v},\nabla \boldsymbol{\mathfrak{z}})_{\mathbf{L}^2(\Omega)}+( u_{h}\boldsymbol{\mathfrak{z}}, \mathbf{v})_{\mathbf{L}^2(\Omega)}+(\mathfrak{r},\text{div }\mathbf{v})_{L^2(\Omega)} =  	(\mathbf{y}_{h}-\mathbf{y}_{\Omega},\mathbf{v})_{\mathbf{L}^2(\Omega)},
\\
(\mathsf{q},\text{div }\boldsymbol{\mathfrak{z}})_{L^2(\Omega)}  =  0 \qquad \forall(\mathbf{v},\mathsf{q})\in \mathbf{H}_0^1(\Omega)\times L_0^2(\Omega).
\end{multline}
We also define $(\mathbf{z}_{h},\mathsf{r}_{h})\in \mathbf{X}_h\times M_h$ as the solution to
\begin{multline}\label{eq:aux_problem_adj_2}
(\nabla \mathbf{v}_{h},\nabla \mathbf{z}_{h})_{\mathbf{L}^2(\Omega)}+( u_{h}\mathbf{z}_{h}, \mathbf{v}_{h})_{\mathbf{L}^2(\Omega)}+(\mathsf{r}_{h},\text{div }\mathbf{v}_{h})_{L^2(\Omega)} \\=  	(\mathbf{y}_{h}-\mathbf{y}_{\Omega},\mathbf{v}_{h})_{\mathbf{L}^2(\Omega)},
\quad
(\mathsf{q}_{h},\text{div }\mathbf{z}_h)_{L^2(\Omega)}  =  0 \qquad \forall (\mathbf{v}_{h},\mathsf{q}_{h})\in \mathbf{X}_h\times M_h.
\end{multline}

\begin{proposition}[error bound]
\label{eq:theorem_adjoint_aux_estimates}
Let $\mathbf{f},\mathbf{y}_{\Omega}\in \mathbf{L}^2(\Omega)$, and let $\Omega \subset \mathbb{R}^d$ be a convex polytope. Let $(\mathbf{z},\mathsf{r})$ and $(\mathbf{z}_{h},\mathsf{r}_{h})$ be the solutions of \eqref{eq:brinkamn_problem_adjoint} and \eqref{eq:aux_problem_adj_2}, respectively. Then,
\begin{align}\label{eq:errorbound_adj}
\|\nabla(\mathbf{z}-\mathbf{z}_{h})\|_{\mathbf{L}^2(\Omega)}+\|\mathsf{r}-\mathsf{r}_{h}\|_{L^2(\Omega)}\lesssim h [ \|\mathbf{f}\|_{\mathbf{L}^2(\Omega)}+\|\mathbf{y}_{\Omega}\|_{\mathbf{L}^2(\Omega)}] + \|u-u_{h}\|_{L^2(\Omega)}.
\end{align}
\end{proposition}

\begin{proof}
We begin the proof by controlling the error $\| (\boldsymbol{\mathfrak{z}} - \mathbf{z}_{h}, \mathfrak{r}-\mathsf{r}_{h}) \|_{\Omega}$. To do this, we use the error bound \eqref{eq:error_estimate_Fem} and also that $\| \nabla \mathbf{y}_{h}\|_{\mathbf{L}^2(\Omega)} \lesssim \| \mathbf{f} \|_{ \mathbf{H}^{-1}(\Omega) }$.
In fact, we have
\begin{align}\label{eq:bound_aux_conv1}
\|\nabla(\boldsymbol{\mathfrak{z}}-\mathbf{z}_{h})\|_{\mathbf{L}^2(\Omega)}+\|\mathfrak{r}-\mathsf{r}_{h}\|_{L^2(\Omega)}
\lesssim 
h( \| \mathbf{f} \|_{ \mathbf{H}^{-1}(\Omega)} + \|\mathbf{y}_{\Omega}\|_{\mathbf{L}^2(\Omega)}).
\end{align}
Let us now note that $(\mathbf{z}- \boldsymbol{\mathfrak{z}},\mathsf{r}-\mathfrak{r})\in\mathbf{H}_0^1(\Omega)\times L_0^2(\Omega)$ solves the following problem:
\begin{multline*}
(\nabla (\mathbf{z}- \boldsymbol{\mathfrak{z}}),\nabla \mathbf{v})_{\mathbf{L}^2(\Omega)}+( u(\mathbf{z}- \boldsymbol{\mathfrak{z}}), \mathbf{v})_{\mathbf{L}^2(\Omega)}+(\mathsf{r}-\mathfrak{r},\text{div }\mathbf{v})_{L^2(\Omega)} \\=((u_{h}-u)\boldsymbol{\mathfrak{z}},\mathbf{v})_{\mathbf{L}^2(\Omega)}+(\mathbf{y}-\mathbf{y}_{h},\mathbf{v})_{\mathbf{L}^2(\Omega)},\quad (\mathsf{q},\text{div }(\mathbf{z}- \boldsymbol{\mathfrak{z}}))_{L^2(\Omega)}  =  0,
\end{multline*}
for all $(\mathbf{v},\mathsf{q}) \in \mathbf{H}_0^1(\Omega) \times L_0^2(\Omega)$. If we set $\mathbf{v} = \mathbf{z}- \boldsymbol{\mathfrak{z}}$ and $\mathsf{q}=0$, we obtain
\begin{equation}
\label{eq:bound_aux_conv2}
\begin{aligned}
\|\nabla(\mathbf{z}- \boldsymbol{\mathfrak{z}})\|_{\mathbf{L}^2(\Omega)}
& \lesssim  \|u-u_{h}\|_{L^2(\Omega)}\|\nabla\boldsymbol{\mathfrak{z}}\|_{\mathbf{L}^{2}(\Omega)}+\|\nabla (\mathbf{y}-\mathbf{y}_{h})\|_{\mathbf{L}^2(\Omega)}\\
& \lesssim  \|u-u_{h}\|_{L^2(\Omega)}[ 1 + \|\mathbf{f}\|_{\mathbf{H}^{-1}(\Omega)} +\|\mathbf{y}_{\Omega}\|_{\mathbf{L}^2(\Omega)}] +h\|\mathbf{f}\|_{\mathbf{L}^2(\Omega)}.
 \end{aligned}
\end{equation}
The bound for $\|\mathsf{r}-\mathfrak{r}\|_{L^2(\Omega)}$ follows from the inf-sup condition \eqref{eq:infsup}, \eqref{eq:errorbound}, and \eqref{eq:bound_aux_conv2}:
\begin{equation}\label{eq:bound_aux_conv3}
\|\mathsf{r}-\mathfrak{r}\|_{L^2(\Omega)}\lesssim 
\|u-u_{h}\|_{L^2(\Omega)}[ 1 + \|\mathbf{f}\|_{\mathbf{H}^{-1}(\Omega)} +\|\mathbf{y}_{\Omega}\|_{\mathbf{L}^2(\Omega)}] +h\|\mathbf{f}\|_{\mathbf{L}^2(\Omega)}.\end{equation}
The bound \eqref{eq:errorbound_adj} follows from \eqref{eq:bound_aux_conv1}, \eqref{eq:bound_aux_conv2}, and \eqref{eq:bound_aux_conv3}. This concludes the proof.
\end{proof}

\subsubsection{Convergence results for the fully discrete scheme}

In this section, we derive two convergence results for the fully discrete scheme. The first result ensures that a sequence of discrete global solutions contains a subsequence that converges to a global solution of the continuous problem \eqref{eq:min_functional}--\eqref{eq:brinkamn_problem_state} as $h \rightarrow 0$. As a second result, we obtain that strict local solutions of the optimal control problem \eqref{eq:min_functional}--\eqref{eq:brinkamn_problem_state} can be approximated by local solutions of the fully discrete optimal control problems.

\begin{theorem}[convergence of global solutions]
\label{theorem_convergence_global_fully}
Let $\mathbf{f}$ and $\mathbf{y}_{\Omega}$ be in $\mathbf{L}^2(\Omega)$. Let $h>0$, and let $\bar{u}_{h}\in \mathbb{U}_{ad,h}$ be a global solution of the fully discrete scheme. Then, there exist nonrelabeled subsequences of $\{\bar{u}_{h}\}_{h>0}$ such that $\bar{u}_{h}\mathrel{\ensurestackMath{\stackon[1pt]{\rightharpoonup}{\scriptstyle\ast}}} \bar{u}$ in the weak topology of $L^{\infty}(\Omega)$ as $h \rightarrow 0$, where $\bar{u}$ is a global solution of \eqref{eq:min_functional}--\eqref{eq:brinkamn_problem_state}. Moreover, we have
\begin{align}\label{eq:convergence_globalsolutions_fully}
\lim_{h \rightarrow 0} \|\bar{u} - \bar{u}_h\|_{L^2(\Omega)}=0,
\qquad \lim_{h \rightarrow 0} j_h(\bar{u}_{h})=j(\bar{u}).
\end{align}
\end{theorem}
\begin{proof}
	We begin the proof by noting that, since $\{\bar{u}_{h}\}_{h>0}\subset \mathbb{U}_{ad,h}$ is uniformly bounded in $L^{\infty}(\Omega)$, we can extract a nonrelabeled subsequence such that $\bar{u}_{h}\mathrel{\ensurestackMath{\stackon[1pt]{\rightharpoonup}{\scriptstyle\ast}}}\bar{u}$ in the weak topology of $L^{\infty}(\Omega)$ as $h\rightarrow 0$. We must now prove that $\bar{u}$ is optimal for \eqref{eq:min_functional}--\eqref{eq:brinkamn_problem_state} and that the limits in \eqref{eq:convergence_globalsolutions_fully} hold. We divide the proof into two steps.
	
	\emph{Step 1:} \emph{$\bar{u}\in \mathbb{U}_{ad}$ is a global solution of problem \eqref{eq:min_functional}--\eqref{eq:brinkamn_problem_state}}. Let $\mathfrak{u}\in \mathbb{U}_{ad}$ be a global solution of \eqref{eq:min_functional}--\eqref{eq:brinkamn_problem_state}. Let $\mathsf{P}_h: L^2(\Omega)\to \mathbb{U}_h$ be the orthogonal projection operator and define $\mathfrak{u}_{h}:=\mathsf{P}_h(\mathfrak{u})$. Since $\mathfrak{u}\in \mathbb{U}_{ad}$, it is clear that $\mathfrak{u}_{h}\in \mathbb{U}_{ad,h}$. On the other hand, the results of Theorem \ref{thm:regularity_optimal_control} guarantee that 	$\mathfrak{u}\in H^1(\Omega)$. Consequently, $\|\mathfrak{u}-\mathfrak{u}_{h}\|_{L^2(\Omega)}\lesssim h \to 0$ as $h \rightarrow 0$. We now use the global optimality of $\mathfrak{u}$, Proposition \ref{theorem_convergence_01}, the global optimality of $\bar{u}_{h}$, and the strong convergence $\mathfrak{u}_{h}\to\mathfrak{u}$ in $L^2(\Omega)$ as $h \rightarrow 0$ to obtain $j(\mathfrak{u})\leq j(\bar{u})\leq \liminf_{h\to 0} j_h (\bar{u}_{h}) \leq \limsup_{h\to 0}j_h (\bar{u}_{h})\leq \limsup_{h \to 0}j_h (\mathfrak{u}_{h})=j(\mathfrak{u})$. As a result, $\bar{u}$ is a global solution of \eqref{eq:min_functional}--\eqref{eq:brinkamn_problem_state} and $j_h (\bar{u}_{h}) \rightarrow j(\bar{u})$ as $h \rightarrow 0$.

\emph{Step 2:} \emph{$\| \bar{u} - \bar{u}_h \|_{L^2(\Omega)}\to 0$ as $h \rightarrow 0$.} In view of the results of Proposition \ref{theorem_convergence_01}, we have that $\bar{\mathbf{y}}_h\to \bar{\mathbf{y}}$ in $\mathbf{H}_0^1(\Omega)$ as $h \rightarrow 0$. This, the fact that $j_h(\bar{u}_{h})\to j(\bar{u})$ as $h\rightarrow 0$, and the weak convergence $\bar{u}_{h}\rightharpoonup \bar{u}$ in $L^2(\Omega)$ as $h\rightarrow 0$ allow us to conclude that $\{\bar{u}_h\}_{h>0}$ converges to $\bar{u}$ in $L^2(\Omega)$ as $h\rightarrow 0$. This completes the proof.
\end{proof}

\begin{theorem}[convergence of local solutions]
Let $\mathbf{f}$ and $\mathbf{y}_{\Omega}$ be in $\mathbf{L}^2(\Omega)$. Let $\bar{u} \in \mathbb{U}_{ad}$ be a strict local minimum of \eqref{eq:min_functional}--\eqref{eq:brinkamn_problem_state}. Then, there exists a sequence $\{\bar{u}_{h}\}_{0 < h \leq h_{\Box}}$ of local minima of the fully discrete scheme such that \eqref{eq:convergence_globalsolutions_fully} holds.
\end{theorem}

\begin{proof}
Since $\bar{u}$ is a strict local minimum of \eqref{eq:min_functional}--\eqref{eq:brinkamn_problem_state}, there exists $\epsilon>0$, so that the problem: Find
$
\min \{j(u) \,| \, u \in \mathbb{U}_{ad}: \|\bar{u}-u\|_{L^2(\Omega)}\leq\epsilon\},
$
admits $\bar{u}$ as the unique solution. We now introduce the discrete problem for $h>0$: Find
$
\min \{j_h(u_{h}) \,| \, u_{h}\in \mathbb{U}_{ad,h} : \|\bar{u}-u_{h}\|_{L^2(\Omega)}\leq\epsilon\}.
$
In the following, we show that this problem admits at least one solution. To this end, we must prove that the set in which the minimum is sought is nonempty. Since $\mathsf{P}_h(\bar{u}) \in \mathbb{U}_{ad,h}$ and $\|\bar{u}-\mathsf{P}_h(\bar{u})\|_{L^2(\Omega)} \to 0$ as $h \rightarrow 0$, there exists $h_{\epsilon}>0$ so that for every $h \leq h_{\epsilon}$ it holds that $\|\bar{u}-\mathsf{P}_h(\bar{u})\|_{L^2(\Omega)}\leq\epsilon$. Consequently, the set $\{ u_{h}\in \mathbb{U}_{ad,h}: \|\bar{u}-u_{h}\|_{L^2(\Omega)}\leq\epsilon\}$ is nonempty for every $h \leq h_{\epsilon}$.

Let $h \leq h_{\epsilon}$, and let $\bar{u}_{h}$ be a global solution to the discrete problem mentioned above. Applying the same arguments as the ones developed in the proof of Theorem \ref{theorem_convergence_global_fully} we obtain the existence of a nonrelabeled subsequence of $\{\bar{u}_{h}\}_{0 < h \leq h_{\epsilon}}$ such that it converges strongly in $L^2(\Omega)$ to a solution of the problem: Find
$
\min \{j(u) \,| \, u \in \mathbb{U}_{ad}: \|\bar{u}-u\|_{L^2(\Omega)}\leq\epsilon\}.
$ 
Since this problem admits $\bar{u}$ as the unique solution, we have that $\bar{u}_{h} \to \bar{u}$ in $L^2(\Omega)$ as $h \rightarrow 0$ for the whole sequence. This guarantees that the constraint $\|\bar{u}-\bar{u}_{h}\|_{L^2(\Omega)}\leq \epsilon$ is not active for $h$ sufficiently small and thus that $\bar{u}_{h}$ solves the fully discrete scheme and \eqref{eq:convergence_globalsolutions_fully} holds. This concludes the proof.
\end{proof}

\subsubsection{Convergence results for the semidiscrete scheme}

As in the case of the fully discrete scheme, we have the following results: Let $\mathbf{f}$ and $\mathbf{y}_{\Omega}$ in $\mathbf{L}^2(\Omega)$.

$\bullet$ Let $h>0$ and let $\bar{\mathsf{u}}_{h}\in \mathbb{U}_{ad}$ be a global solution of the semidiscrete scheme. Then, there exist nonrelabeled subsequences of $\{\bar{\mathsf{u}}_{h}\}_{h>0}$ such that $\bar{\mathsf{u}}_{h} \mathrel{\ensurestackMath{\stackon[1pt]{\rightharpoonup}{\scriptstyle\ast}}} \bar{u}$ in the weak topology of $L^{\infty}(\Omega)$ as $h \rightarrow 0$, where $\bar{u}$ corresponds to a global solution of the optimal control problem \eqref{eq:min_functional}--\eqref{eq:brinkamn_problem_state}. Moreover, $\|\bar{u}-\bar{\mathsf{u}}_{h}\|_{L^2(\Omega)} \rightarrow 0$ and $j_h(\bar{\mathsf{u}}_{h}) \rightarrow j(\bar{u})$ as $h \rightarrow 0$.

$\bullet$ Let $\bar{u}\in\mathbb{U}_{ad}$ be a strict local minimum of \eqref{eq:min_functional}--\eqref{eq:brinkamn_problem_state}. Then, there exists a sequence $\{\bar{\mathsf{u}}_{h}\}_{0 < h \leq h_{\circ}}$ of local minima of the semidiscrete scheme such that $\|\bar{u}-\bar{\mathsf{u}}_{h}\|_{L^2(\Omega)} \rightarrow 0$ and $j_h(\bar{\mathsf{u}}_{h}) \rightarrow j(\bar{u})$ as $h \rightarrow 0$.

%%%%%%%%%%%%%%%%%%%%%%%%%%%%%%%%%%%%%%%%%%%%%%%%%%%%%%%%%%%%%%%%
%%%%%%%%%%%%%%%%%%%%%%%%%%%%%%%%%%%%%%%%%%%%%%%%%%%%%%%%%%%%%%%%
%%%%%%%%%%%%%%%%%%%%%%%%%%%%%%%%%%%%%%%%%%%%%%%%%%%%%%%%%%%%%%%%
%%%%%%%%%%%%%%%%%%%%%%%%%%%%%%%%%%%%%%%%%%%%%%%%%%%%%%%%%%%%%%%%


\section{A priori error bounds}
The main goal of this section is to derive error bounds for the fully and semidiscrete methods presented in \S \ref{fully_discrete_framework} and \S \ref{semi_discrete_framework}, respectively.

\subsection{The fully discrete scheme}
Let $\{\bar{u}_{h}\}_{h>0}$ be a sequence of local minima of the fully discrete scheme such that $\bar{u}_{h} \to \bar{u}$ in $L^2(\Omega)$ as $h\rightarrow 0$, where $\bar{u}$ corresponds to a local solution of \eqref{eq:min_functional}--\eqref{eq:brinkamn_problem_state}. In this section, we will derive the bound
\begin{align}\label{eq:error_estimate_apriori_fully}
\|\bar{u}-\bar{u}_{h}\|_{L^2(\Omega)}\lesssim h \qquad\forall h\in(0, h_{\dagger}).
\end{align}

We begin our analysis with the following instrumental result.

\begin{lemma}[auxiliary error estimate]\label{lemma_auxiliary_error_fd}
Let $\boldsymbol{\mathbf{f}}$ and $\mathbf{y}_{\Omega}$ be in $\mathbf{L}^2(\Omega)$. Let $\Omega \subset \mathbb{R}^d$ be a convex polytope. Let us assume that $\bar{u}\in \mathbb{U}_{ad}$ satisfies the second-order optimality conditions \eqref{equivalence_condition}. If \eqref{eq:error_estimate_apriori_fully} is false, then there exists $h_{\star}>0$ such that
\begin{align}\label{eq:auxiliary_error_estimate_fd}
\mathfrak{C}\|\bar{u}-\bar{u}_{h}\|_{L^2(\Omega)}^2\leq [j'(\bar{u}_{h})-j'(\bar{u})](\bar{u}_{h}-\bar{u})
 \quad
\forall h<h_{\star},
\quad
\mathfrak{C}:=2^{-1}\min\{\mu,\alpha\}.
\end{align}
Here, $\alpha$ is the control cost, and $\mu$ denotes the constant appearing in \eqref{equivalence_condition}.
\end{lemma}
\begin{proof}
We proceed on the basis of a contradiction argument as in \cite[Section 7]{MR2272157}: Since \eqref{eq:error_estimate_apriori_fully} is false, there exists a subsequence $\{h_k\}_{k\in\mathbb{N}}\subset \mathbb{R}^{+}$ such that
\begin{align}\label{eq:auxuliary_error_estimate_fd}
\lim_{h_{k} \rightarrow 0} \|\bar{u}-\bar{u}_{h_k}\|_{L^2(\Omega)} = 0,
\qquad \lim_{h_{k} \rightarrow 0} h_k^{-1} \|\bar{u}-\bar{u}_{h_k}\|_{L^2(\Omega)}=+\infty.
\end{align}
In the following, we omit the subindex $k$ to simplify the notation and denote $\bar{u}_{h_k}=\bar{u}_{h}$. We note that $h \rightarrow 0$ as $k \uparrow \infty$. Define $w_h:=(\bar{u}_h-\bar{u})/\|\bar{u}_h-\bar{u}\|_{L^2(\Omega)}$ for each $h>0$ and note that $\{ w_h \}_{h>0}$ is uniformly bounded in $L^2(\Omega)$. Therefore, there exists a nonrelabeled subsequence such that $w_h \rightharpoonup w$ in $L^2(\Omega)$ as $h\rightarrow 0$. The rest of the proof is divided into two steps.

\emph{Step 1:} \emph{$w\in C_{\bar{u}}$}. Let us first recall that $C_{\bar{u}}$ is defined in \eqref{eq:cone}. Since $\bar{u}_{h} \in \mathbb{U}_{ad,h}$ for any $h>0$, we have that $w_h$ satisfies the sign conditions in \eqref{eq:cone_condition}. Consequently, the weak limit $w$ also satisfies \eqref{eq:cone_condition}. In what follows, we prove that $\bar{\mathfrak{d}}(x) \neq 0$ implies that $w(x)=0$ for a.e.~$x\in\Omega$. To this end, we introduce the variable $\bar{\mathfrak{d}}_h:=\alpha \bar{u}_h-\bar{\mathbf{y}}_h\cdot\bar{\mathbf{z}}_h$ and recall that $\bar{\mathfrak{d}}=\alpha \bar{u}-\bar{\mathbf{y}}\cdot\bar{\mathbf{z}}$. We now control $\|\bar{\mathfrak{d}} - \bar{\mathfrak{d}}_h\|_{L^2(\Omega)}$ in view of H\"older's inequality, the bounds \eqref{eq:errorbound_adj} and \eqref{eq:errorbound}, and property $\|\bar{u}-\bar{u}_h\|_{L^2(\Omega)}\to 0$ as $h\rightarrow 0$:
\begin{multline*}
\|\bar{\mathfrak{d}}-\bar{\mathfrak{d}}_h\|_{L^2(\Omega)}\lesssim \|\bar{u}-\bar{u}_h\|_{L^2(\Omega)}+\|\nabla\bar{\mathbf{y}}_h\|_{\mathbf{L}^2(\Omega)}\|\nabla( \bar{\mathbf{z}}-\bar{\mathbf{z}}_h)\|_{\mathbf{L}^2(\Omega)}
+ \|\nabla\bar{\mathbf{z}}\|_{\mathbf{L}^2(\Omega)}
\\
\cdot \|\nabla( \bar{\mathbf{y}}-\bar{\mathbf{y}}_h)\|_{\mathbf{L}^2(\Omega)}\lesssim \|\bar{u}-\bar{u}_h\|_{L^2(\Omega)} + h \|\mathbf{f}\|_{\mathbf{L}^2(\Omega)}[\|\mathbf{f}\|_{\mathbf{L}^2(\Omega)} + \| \mathbf{y}_{\Omega} \|_{\mathbf{L}^2(\Omega)}]
\to 0
\end{multline*}
as $h \rightarrow 0$. We have thus obtained that $\bar{\mathfrak{d}}_h \rightarrow \bar{\mathfrak{d}}$ in $L^2(\Omega)$ as $h \rightarrow 0$. Since $w_h \rightharpoonup w$ in $L^2(\Omega)$ as $h\rightarrow 0$, we have
\begin{eqnarray*}
\int_{\Omega}\bar{\mathfrak{d}}(x)w(x)\mathrm{d}x=
\lim_{h\rightarrow 0} \frac{1}{\|\bar{u}_h-\bar{u}\|_{L^2(\Omega)}}
\left[\int_{\Omega}\bar{\mathfrak{d}}_h(\mathsf{P}_h(\bar{u})-\bar{u})\mathrm{d}x+\int_{\Omega}\bar{\mathfrak{d}}_h(\bar{u}_h-\mathsf{P}_h(\bar{u}))\mathrm{d}x\right],
\end{eqnarray*}
where $\mathsf{P}_h: L^2(\Omega) \rightarrow \mathbb{U}_h$ denotes the $L^2$-orthogonal projection operator. Since $\mathsf{P}_h(\bar{u})\in \mathbb{U}_{ad,h}$, the discrete variational inequality \eqref{eq:variational_inequality_discrete} yields $(\bar{\mathfrak{d}}_h,\bar{u}_h-\mathsf{P}_h(\bar{u}))_{L^2(\Omega)}\leq 0$. Let us now note that $\|\bar{\mathfrak{d}}_h\|_{L^2(\Omega)}\leq\|\bar{\mathfrak{d}}_h-\bar{\mathfrak{d}}\|_{L^2(\Omega)}+\|\bar{\mathfrak{d}}\|_{L^2(\Omega)}\lesssim 1$ for every $h \leq \mathfrak{h}$, which implies that $\{\|\bar{\mathfrak{d}}_h\|_{L^2(\Omega)}\}_{0 < h \leq \mathfrak{h}}$ is uniformly bounded in $\mathbb{R}$. Thus, 
\begin{eqnarray*}
\int_{\Omega}\bar{\mathfrak{d}}(x)w(x)\mathrm{d}x
\leq
\lim_{h \rightarrow 0}
\frac{1}{\|\bar{u}_h-\bar{u}\|_{L^2(\Omega)}}
 \int_{\Omega}\bar{\mathfrak{d}}_h(\mathsf{P}_h(\bar{u})-\bar{u})\mathrm{d}x
\lesssim \lim_{h \rightarrow 0}\dfrac{\|\mathsf{P}_h(\bar{u})-\bar{u}\|_{L^2(\Omega)}}{\|\bar{u}_h-\bar{u}\|_{L^2(\Omega)}}=0.
\end{eqnarray*}
To obtain the last equality, we used the regularity property $\bar{u} \in H^1(\Omega)$ (cf.~Theorem \ref{thm:regularity_optimal_control}), the interpolation error bound $\| \bar{u} - \mathsf{P}_h(\bar{u}) \|_{L^2(\Omega)} \lesssim h | \bar{u} |_{H^1(\Omega)}$, and the right-hand side limit in \eqref{eq:auxuliary_error_estimate_fd}. On the other hand, since $w$ satisfies the sign conditions in \eqref{eq:cone_condition} and $\bar u$ satisfies the proyection formula \eqref{eq:representation}, we obtain that $\bar{\mathfrak{d}}(x)w(x)\geq 0$ for a.e.~$x\in\Omega$. Therefore, $\int_{\Omega}|\bar{\mathfrak{d}}(x)w(x)|\mathrm{d}x = \int_{\Omega}\bar{\mathfrak{d}}(x)w(x)\mathrm{d}x \leq 0$. Consequently, if $\bar{\mathfrak{d}}(x)\neq 0$, then $w(x)=0$ for a.e.~$x\in\Omega$. This proves that $w\in C_{\bar{u}}$.

\emph{Step 2:} \emph{The error bound \eqref{eq:auxiliary_error_estimate_fd}.} We apply the mean value theorem to obtain 
\begin{align}\label{eq:mean_value_theorem}
[j'(\bar{u}_h)-j'(\bar{u})](\bar{u}_h-\bar{u})=j''(\hat{u}_h)(\bar{u}_h-\bar{u})^2,\quad \hat{u}_h:=\bar{u}+\theta_h (\bar{u}_h - \bar{u}),
\end{align}
where $\theta_h \in (0,1)$. Let $(\mathbf{y}(\hat{u}_h),\mathsf{p}(\hat{u}_h))\in \mathbf{H}_0^1(\Omega)\times L_0^2(\Omega)$ be the solution to  \eqref{eq:brinkamn_problem_state}, where $u$ is replaced by $\hat{u}_h$, and let $(\mathbf{z}(\hat{u}_h),\mathsf{r}(\hat{u}_h))\in \mathbf{H}_0^1(\Omega)\times L_0^2(\Omega)$ be the solution to \eqref{eq:brinkamn_problem_adjoint}, where $\mathbf{y}$ and $u$ are replaced by $\mathbf{y}(\hat{u}_h)$ and $\hat{u}_h$, respectively.  We now note that $(\bar{\mathbf{y}}-\mathbf{y}(\hat{u}_h),\bar{\mathsf{p}}-\mathsf{p}(\hat{u}_h))\in\mathbf{H}_0^1(\Omega)\times L_0^2(\Omega)$ solves the following problem:
 \begin{multline*}
(\nabla (\bar{\mathbf{y}}-\mathbf{y}(\hat{u}_{h})),\nabla \mathbf{v})_{\mathbf{L}^2(\Omega)}+( \bar{u}(\bar{\mathbf{y}}-\mathbf{y}(\hat{u}_{h})), \mathbf{v})_{\mathbf{L}^2(\Omega)}-(\bar{\mathsf{p}}-\mathsf{p}(\hat{u}_{h}),\text{div }\mathbf{v})_{L^2(\Omega)} \\=((\hat{u}_{h}-\bar{u})\mathbf{y}(\hat{u}_{h}),\mathbf{v})_{\mathbf{L}^2(\Omega)},\quad (\mathsf{q},\text{div }(\bar{\mathbf{y}}-\mathbf{y}(\hat{u}_{h})))_{L^2(\Omega)}  =  0,
\end{multline*}
for all $(\mathbf{v},\mathsf{q}) \in \mathbf{H}_0^1(\Omega) \times L_0^2(\Omega)$. If we set $\mathbf{v} = \bar{\mathbf{y}}-\mathbf{y}(\hat{u}_{h})$ and $\mathsf{q}=0$, we obtain the bound $\|\nabla(\bar{\mathbf{y}}-\mathbf{y}(\hat{u}_{h}))\|_{\mathbf{L}^2(\Omega)}\lesssim \|\bar{u}-\hat{u}_{h}\|_{L^2(\Omega)}$. Since $\bar{u}_h \to \bar{u}$ in $L^2(\Omega)$ as $h \rightarrow 0$, we can conclude that $\mathbf{y}(\hat{u}_h)\to \bar{\mathbf{y}}$ in $\mathbf{H}_0^1(\Omega)$ and that $\mathsf{p}(\hat{u}_h)\to\bar{\mathsf{p}}$ in $L_0^2(\Omega)$ as $h \rightarrow 0$; the latter is a consequence of the inf--sup condition \eqref{eq:infsup}. Similarly, we can conclude that $\mathbf{z}(\hat{u}_h)\to \bar{\mathbf{z}}$ in $\mathbf{H}_0^1(\Omega)$ and that $\mathsf{r}(\hat{u}_h) \to \bar{\mathsf{r}}$ in $L_0^2(\Omega)$ as $h \rightarrow 0$. Let us now define $(\boldsymbol{\varphi}(w_h),\zeta(w_h))\in \mathbf{H}_0^1(\Omega)\times L_0^2(\Omega)$ as the solution to \eqref{eq:problem_S_first}, where $u$, $\mathbf{y}$, and $v$ are replaced by $\hat{u}_h$, $\mathbf{y}(\hat{u}_{h})$, and $w_h$, respectively. If we apply the arguments from step 2 of the proof of Theorem \ref{theorem_second_order_optimal}, we obtain that $w_h\rightharpoonup w$ in $L^2(\Omega)$ as $h \rightarrow 0$ implies that $\boldsymbol{\varphi}(w_h)\rightharpoonup \boldsymbol{\varphi}$ in $\mathbf{H}_0^1(\Omega)$ and that $\zeta(w_h)\rightharpoonup \zeta$ in $L_0^2(\Omega)$ as $h \rightarrow 0$. We thus use the identity \eqref{eq:identity_j}, the definition of $w_h$, and the second-order equivalent conditions in \eqref{equivalence_condition} to arrive at
\begin{multline*}
\lim_{h\to 0}j''(\hat{u}_h)w_h^2
=
\lim_{h\to 0}
\left[
\alpha\|w_h\|_{L^2(\Omega)}^2-2(w_h\boldsymbol{\varphi}(w_h),\mathbf{z}(\hat{u}_h))_{\mathbf{L}^2(\Omega)}+\|\boldsymbol{\varphi}(w_h)\|_{\mathbf{L}^2(\Omega)}^2
\right]
\\
=\alpha-2(w\boldsymbol{\varphi},\bar{\mathbf{z}})_{\mathbf{L}^2(\Omega)}+\|\boldsymbol{\varphi}\|_{\mathbf{L}^2(\Omega)}^2=\alpha+j''(\bar{u})w^2-\alpha \|w\|_{L^2(\Omega)}^2\geq \alpha + (\mu-\alpha) \|w\|_{L^2(\Omega)}^2,
\end{multline*}
where we have used that $\boldsymbol{\varphi}(w_h)\rightarrow \boldsymbol{\varphi}$ in $\mathbf{L}^4(\Omega)$ as $h \rightarrow 0$. Since $\|w\|_{L^2(\Omega)}\leq 1$, we can conclude that $\lim_{h\rightarrow 0}j''(\hat{u}_h)w_h^2\geq \min\{\mu,\alpha\}>0$. Thus, we deduce the existence of $h_{\star}>0$ such that such that $j''(\hat{u}_h)w_h^2\geq \min\{\mu,\alpha\}/2$ for every $h<h_{\star}$. The bound \eqref{eq:auxiliary_error_estimate_fd} therefore follows from the definition of $w_h$ and the identity \eqref{eq:mean_value_theorem}.
\end{proof}

We are now ready to prove the error bound \eqref{eq:error_estimate_apriori_fully}.

\begin{theorem}[a priori error bound]
\label{theorem_estimates_control}
Let $\boldsymbol{\mathbf{f}}$ and $\mathbf{y}_{\Omega}$ be in $\mathbf{L}^2(\Omega)$. Let $\Omega \subset \mathbb{R}^d$ be a convex polytope. Let us assume that $\bar{u}\in \mathbb{U}_{ad}$ satisfies the second-order optimality conditions \eqref{equivalence_condition}. Then, there exists $h_{\dagger}>0$ such that
\begin{align}\label{eq:error_estimate_fd}
\|\bar{u}-\bar{u}_{h}\|_{L^2(\Omega)}
\lesssim 
h\|\mathbf{f}\|_{\mathbf{L}^2(\Omega)}
\left[ 
\|\mathbf{f}\|_{\mathbf{L}^2(\Omega)}+\|\mathbf{y}_{\Omega}\|_{\mathbf{L}^2(\Omega)}
\right]
\quad\forall h\in(0, h_{\dagger}).
\end{align}
\end{theorem}
\begin{proof}
We proceed by contradiction: assume that the error bound \eqref{eq:error_estimate_fd} is false. We can therefore use the result of Lemma \ref{lemma_auxiliary_error_fd} to obtain that \eqref{eq:auxiliary_error_estimate_fd} holds for every $h<h_{\star}$. On the other hand, if we set $u=\bar{u}_h$ in \eqref{eq:ineq_variational} and $u_h=\mathsf{P}_h(\bar{u})$ in \eqref{eq:variational_inequality_discrete} we can conclude that $-j'(\bar{u})(\bar{u}_h-\bar{u})\leq 0$ and $j'_h(\bar{u}_h)(\mathsf{P}_h(\bar{u})-\bar{u}_h)\geq 0$, respectively. Thus,
\begin{multline}\label{eq:estimates_aux_error_fullydiscrete}
\|\bar{u}-\bar{u}_h\|_{L^2(\Omega)}^2\lesssim j'(\bar{u}_h)(\bar{u}_h-\bar{u})+j'_h (\bar{u}_h)(\mathsf{P}_h(\bar{u})-\bar{u}_h)\\
=j'(\bar{u}_h)(\mathsf{P}_h(\bar{u})-\bar{u})+\left[j'(\bar{u}_h)-j'_h (\bar{u}_h)\right](\bar{u}_h-\mathsf{P}_h(\bar{u}))=:\text{I}_h + \text{II}_h.
\end{multline}

To control the term $\text{I}_h$, we introduce $(\hat{\mathbf{y}},\hat{\mathsf{p}})\in \mathbf{H}_0^1(\Omega)\times L_0^2(\Omega)$ as the solution to
\begin{align}
\label{eq:aux_problem_adj_1_error}
(\nabla\hat{\mathbf{y}},\nabla \mathbf{v})_{\mathbf{L}^2(\Omega)}+( \bar{u}_{h}\hat{\mathbf{y}}, \mathbf{v})_{\mathbf{L}^2(\Omega)}-(\hat{\mathsf{p}},\text{div }\mathbf{v})_{L^2(\Omega)} =  	(\mathbf{f},\mathbf{v})_{\mathbf{L}^2(\Omega)}
\end{align}
and $(\mathsf{q},\text{div }\hat{\mathbf{y}})_{L^2(\Omega)} = 0$ for all $(\mathbf{v},\mathsf{q}) \in \mathbf{H}_0^1(\Omega) \times L_0^2(\Omega)$. We also introduce the pair $(\hat{\mathbf{z}},\hat{\mathsf{r}})\in \mathbf{H}_0^1(\Omega)\times L_0^2(\Omega)$ as the solution to
\begin{align}
\label{eq:aux_problem_adj_2_error}
(\nabla\hat{\mathbf{z}},\nabla \mathbf{v})_{\mathbf{L}^2(\Omega)}+( \bar{u}_{h}\hat{\mathbf{z}}, \mathbf{v})_{\mathbf{L}^2(\Omega)}+(\hat{\mathsf{r}},\text{div }\mathbf{v})_{L^2(\Omega)} = 	(\hat{\mathbf{y}}-\mathbf{y}_{\Omega},\mathbf{v})_{\mathbf{L}^2(\Omega)}
\end{align}
and $(\mathsf{q},\text{div }\hat{\mathbf{z}})_{L^2(\Omega)}  = 0$ for all $(\mathbf{v},\mathsf{q}) \in \mathbf{H}_0^1(\Omega) \times L_0^2(\Omega)$. We thus control $\text{I}_h$ as follows:
\begin{eqnarray*}
\text{I}_h 
&=&(\alpha \bar{u}_h-\hat{\mathbf{y}}\cdot\hat{\mathbf{z}},\mathsf{P}_h(\bar{u})-\bar{u})_{L^2(\Omega)}
=(\mathsf{P}_h(\hat{\mathbf{y}}\cdot\hat{\mathbf{z}})-\hat{\mathbf{y}}\cdot\hat{\mathbf{z}},\mathsf{P}_h(\bar{u})-\bar{u})_{L^2(\Omega)}
\\
&\leq & \|\mathsf{P}_h(\hat{\mathbf{y}}\cdot\hat{\mathbf{z}})-\hat{\mathbf{y}}\cdot\hat{\mathbf{z}}\|_{L^2(\Omega)} \|\mathsf{P}_h(\bar{u})-\bar{u}\|_{L^2(\Omega)}
\end{eqnarray*}
where we used that $(\alpha \bar{u}_h,\mathsf{P}_h(\bar{u})-\bar{u})_{L^2(\Omega)} = 0 = (\mathsf{P}_h(\hat{\mathbf{y}}\cdot\hat{\mathbf{z}}),\mathsf{P}_h(\bar{u})-\bar{u})_{L^2(\Omega)}$. The regularity results of Theorem \ref{thm:regularity_optimal_control} combined with a basic error bound for $\mathsf{P}_h$ show that $\|\mathsf{P}_h(\bar{u})-\bar{u}\|_{L^2(\Omega)} \lesssim \|\mathbf{f}\|_{\mathbf{L}^2(\Omega)}[\|\mathbf{f}\|_{\mathbf{L}^2(\Omega)}+\|\mathbf{y}_{\Omega}\|_{\mathbf{L}^2(\Omega)}]$. To bound $\|\mathsf{P}_h(\hat{\mathbf{y}}\cdot\hat{\mathbf{z}})-\hat{\mathbf{y}}\cdot\hat{\mathbf{z}}\|_{L^2(\Omega)}$, we first note that since $\mathbf{f}, \hat{\mathbf{y}}-\mathbf{y}_{\Omega} \in \mathbf{L}^2(\Omega)$, we have that $\hat{\mathbf{y}}, \hat{\mathbf{z}} \in \mathbf{C}(\bar \Omega)$, (cf.~Theorem \ref{theorem_regularity_Linfty}). Thus
$
\|\nabla(\hat{\mathbf{y}}\cdot \hat{\mathbf{z}})\|_{\mathbf{L}^2(\Omega)} 
\leq 
\|\hat{\mathbf{z}} \|_{\mathbf{C}(\bar \Omega)} \| \nabla \hat{\mathbf{y}}\|_{\mathbf{L}^2(\Omega)}
+
\|\hat{\mathbf{y}} \|_{\mathbf{C}(\bar \Omega)} \| \nabla \hat{\mathbf{z}}\|_{\mathbf{L}^2(\Omega)}
\lesssim
\|\mathbf{f}\|_{\mathbf{L}^2(\Omega)}
[
\|\mathbf{f}\|_{\mathbf{L}^2(\Omega)}+\|\mathbf{y}_{\Omega}\|_{\mathbf{L}^2(\Omega)}
].
$
A basic error estimate for $\mathsf{P}_h$ thus shows that
$
\|\mathsf{P}_h(\hat{\mathbf{y}}\cdot\hat{\mathbf{z}})-\hat{\mathbf{y}}\cdot\hat{\mathbf{z}}\|_{L^2(\Omega)}
\lesssim h \|\mathbf{f}\|_{\mathbf{L}^2(\Omega)}[\|\mathbf{f}\|_{\mathbf{L}^2(\Omega)}+\|\mathbf{y}_{\Omega}\|_{\mathbf{L}^2(\Omega)}].
$
As a result, we have obtained the following bound:
\[
\text{I}_h \lesssim h^2\|\mathbf{f}\|_{\mathbf{L}^2(\Omega)}^2
\left[ 
\|\mathbf{f}\|_{\mathbf{L}^2(\Omega)}+\|\mathbf{y}_{\Omega}\|_{\mathbf{L}^2(\Omega)}
\right]^2.
\]

We now bound $\text{II}_h$. To this end, we first note that
$
\text{II}_h = (\hat{\mathbf{y}}\cdot\hat{\mathbf{z}}-\bar{\mathbf{y}}_h\cdot\bar{\mathbf{z}}_h,\bar{u}_h-\mathsf{P}_h(\bar{u}))_{L^2(\Omega)}=(\hat{\mathbf{y}}\cdot\hat{\mathbf{z}}-\bar{\mathbf{y}}_h\cdot\bar{\mathbf{z}}_h,\mathsf{P}_h(\bar{u}_h-\bar{u}))_{L^2(\Omega)}.$
If we add and subtract $(\bar{\mathbf{y}}_h\cdot\hat{\mathbf{z}},\mathsf{P}_h(\bar{u}_h-\bar{u}))_{L^2(\Omega)}$ and use a suitable Sobolev embedding and Young's inequality, we obtain
\begin{multline*}
\text{II}_h \lesssim \left[ 
\|\nabla\hat{\mathbf{z}}\|_{\mathbf{L}^{2}(\Omega)}\|\nabla(\hat{\mathbf{y}}-\bar{\mathbf{y}}_h)\|_{\mathbf{L}^{2}(\Omega)} 
+ 
\|\nabla\bar{\mathbf{y}}_h\|_{\mathbf{L}^{2}(\Omega)}\|\nabla(\hat{\mathbf{z}}-\bar{\mathbf{z}}_h)\|_{\mathbf{L}^{2}(\Omega)} \right]
\|\bar{u}-\bar{u}_h\|_{L^2(\Omega)}
\\
\lesssim \|\nabla\hat{\mathbf{z}}\|_{\mathbf{L}^{2}(\Omega)}^2\|\nabla(\hat{\mathbf{y}}-\bar{\mathbf{y}}_h)\|_{\mathbf{L}^2(\Omega)}^2 
+
\|\nabla\bar{\mathbf{y}}_h\|_{\mathbf{L}^{2}(\Omega)}^2 \|\nabla(\hat{\mathbf{z}}-\bar{\mathbf{z}}_h)\|_{\mathbf{L}^2(\Omega)}^2+\tfrac{1}{2}\|\bar{u}-\bar{u}_h\|_{L^2(\Omega)}^2,
\end{multline*}
upon using that $\mathsf{P}_h$ is stable in $L^2(\Omega)$. Note that $\|\nabla\hat{\mathbf{z}}\|_{\mathbf{L}^{2}(\Omega)} \lesssim \| \mathbf{f} \|_{\mathbf{L}^2(\Omega)}+\|\mathbf{y}_{\Omega}\|_{\mathbf{L}^2(\Omega)}$ and that $\|\nabla\bar{\mathbf{y}}_h\|_{\mathbf{L}^{2}(\Omega)} \lesssim \| \mathbf{f} \|_{\mathbf{L}^2(\Omega)}$ for every $h>0$. On the other hand, an immediate application of the estimate \eqref{eq:error_estimate_Fem} shows that $\|\nabla(\hat{\mathbf{y}}-\bar{\mathbf{y}}_h)\|_{\mathbf{L}^2(\Omega)}\lesssim h \|\mathbf{f}\|_{\mathbf{L}^2(\Omega)}$. To control $\|\nabla(\hat{\mathbf{z}}-\bar{\mathbf{z}}_h)\|_{\mathbf{L}^2(\Omega)}$, we introduce the problem: Find $(\tilde{\mathbf{z}},\tilde{\mathsf{r}})\in \mathbf{H}_0^1(\Omega)\times L_0^2(\Omega)$ such that
\begin{align}
\label{eq:aux_problem_adj_3_error}
(\nabla\tilde{\mathbf{z}},\nabla \mathbf{v})_{\mathbf{L}^2(\Omega)}+( \bar{u}_{h}\tilde{\mathbf{z}}, \mathbf{v})_{\mathbf{L}^2(\Omega)}+(\tilde{\mathsf{r}},\text{div }\mathbf{v})_{L^2(\Omega)} = 	(\bar{\mathbf{y}}_h-\mathbf{y}_{\Omega},\mathbf{v})_{\mathbf{L}^2(\Omega)}
\end{align}
and $(\mathsf{q},\text{div }\tilde{\mathbf{z}})_{L^2(\Omega)} =  0$ for all $(\mathbf{v},\mathsf{q}) \in \mathbf{H}_0^1(\Omega) \times L_0^2(\Omega)$. A simple application of the triangular inequality yields $\|\nabla(\hat{\mathbf{z}}-\bar{\mathbf{z}}_h)\|_{\mathbf{L}^2(\Omega)}\leq \|\nabla(\hat{\mathbf{z}}-\tilde{\mathbf{z}})\|_{\mathbf{L}^2(\Omega)}+\|\nabla(\tilde{\mathbf{z}}-\bar{\mathbf{z}}_h)\|_{\mathbf{L}^2(\Omega)}$. The control of $\|\nabla(\hat{\mathbf{z}}-\tilde{\mathbf{z}})\|_{\mathbf{L}^2(\Omega)}$ follows from a stability bound for the problem that $(\hat{\mathbf{z}}-\tilde{\mathbf{z}},\hat{\mathsf{r}}-\tilde{\mathsf{r}})$ solves, while the control of $\|\nabla(\tilde{\mathbf{z}}-\bar{\mathbf{z}}_h)\|_{\mathbf{L}^2(\Omega)}$ follows from \eqref{eq:error_estimate_Fem}. As a result, we obtain the following bound:
$
\|\nabla(\hat{\mathbf{z}}-\bar{\mathbf{z}}_h)\|_{\mathbf{L}^2(\Omega)}
\lesssim \|\hat{\mathbf{y}}-\bar{\mathbf{y}}_h\|_{\mathbf{L}^2(\Omega)}+ h(\|\mathbf{f}\|_{\mathbf{L}^2(\Omega)}+\|\mathbf{y}_{\Omega}\|_{\mathbf{L}^2(\Omega)}).$
With all these ingredients, we can bound the term $\text{II}_h$ as follows:
\begin{align*}
\text{II}_h \lesssim h^2 \|\mathbf{f}\|_{\mathbf{L}^2(\Omega)}^2
\left[\|\mathbf{f}\|_{\mathbf{L}^2(\Omega)}+\|\mathbf{y}_{\Omega}\|_{\mathbf{L}^2(\Omega)}\right]^2
+\tfrac{1}{2}\|\bar{u}-\bar{u}_h\|_{L^2(\Omega)}^2.
\end{align*}

The error bound \eqref{eq:error_estimate_fd} thus follows from replacing the estimates obtained for the terms $\text{I}_h$ and $\text{II}_h$ into \eqref{eq:estimates_aux_error_fullydiscrete}. This, which is a contradiction, concludes the proof.
\end{proof}

\begin{corollary}[a priori error bounds]
Let the assumptions of Theorem \ref{theorem_estimates_control} hold. Then, there exists $h_{\dagger}>0$ such that
\begin{align*}\label{eq:estimates_error_fully_allvariables}
\|\nabla(\bar{\mathbf{y}}-\bar{\mathbf{y}}_{h})\|_{\mathbf{L}^2(\Omega)}+\|\bar{\mathsf{p}}-\bar{\mathsf{p}}_{h}\|_{L^2(\Omega)}
\lesssim h\|\mathbf{f}\|_{\mathbf{L}^2(\Omega)}( 1 + \|\mathbf{f}\|_{\mathbf{L}^2(\Omega)}+\|\mathbf{y}_{\Omega}\|_{\mathbf{L}^2(\Omega)}),
\\
\|\nabla(\bar{\mathbf{z}}-\bar{\mathbf{z}}_{h})\|_{\mathbf{L}^2(\Omega)}+\|\bar{\mathsf{r}}-\bar{\mathsf{r}}_{h}\|_{L^2(\Omega)}
\lesssim 
h ( \|\mathbf{f}\|_{\mathbf{L}^2(\Omega)} + 1 )( \|\mathbf{f}\|_{\mathbf{L}^2(\Omega)}+\|\mathbf{y}_{\Omega}\|_{\mathbf{L}^2(\Omega)}),
\end{align*}
for all $h<h_{\dagger}$. In both bounds the hidden constants do not depend on $\mathbf{f}$, $\mathbf{y}_{\Omega}$, and $h$.
\end{corollary}

\begin{proof}
The proof follows directly from a combination of the bounds given in Propositions \ref{theorem_convergence_01} and \ref{eq:theorem_adjoint_aux_estimates} and Theorem \ref{theorem_estimates_control}. This concludes the proof.
\end{proof}

\subsection{The semidiscrete scheme}
Let $\{\bar{\mathsf{u}}_{h}\}_{h>0}$ be a sequence of local minima of the semidiscrete scheme such that $\bar{\mathsf{u}}_{h}\to\bar{u}$ in $L^2(\Omega)$ as $h \rightarrow 0$, where $\bar{u}$ denotes a local solution of \eqref{eq:min_functional}--\eqref{eq:brinkamn_problem_state}. In this section, we will derive the error bound \eqref{eq:error_estimate_sd}.

As in the fully discrete scheme, the following instrumental result is important.

\begin{lemma}[auxiliary error estimate]
\label{lemma_auxiliary_error_sd}
Let $\boldsymbol{\mathbf{f}}$ and $\mathbf{y}_{\Omega}$ be in  $\mathbf{L}^2(\Omega)$. Let $\Omega \subset \mathbb{R}^d$ be a convex polytope. Let us assume that $\bar{u}\in \mathbb{U}_{ad}$ satisfies the second-order optimality conditions \eqref{equivalence_condition}. Then, there exists $h_{\star}>0$ such that
\begin{align}\label{eq:auxiliary_error_estimate_sd}
\mathfrak{C}\|\bar{u}-\bar{\mathsf{u}}_{h}\|_{L^2(\Omega)}^2\leq [j'(\bar{\mathsf{u}}_{h})-j'(\bar{u})](\bar{\mathsf{u}}_{h}-\bar{u})
\quad
\forall h<h_{\star},
\quad
\mathfrak{C}:=2^{-1}\min\{\mu,\alpha\}.
\end{align}
Here, $\alpha$ is the control cost, and $\mu$ denotes the constant appearing in \eqref{equivalence_condition}.
\end{lemma}
\begin{proof}
Define $w_h:=(\bar{\mathsf{u}}_h-\bar{u})/\|\bar{\mathsf{u}}_h-\bar{u}\|_{L^2(\Omega)}$. We note that $w_h \rightharpoonup w$ in $L^2(\Omega)$ up to a subsequence, if necessary, as $h \to 0$. The arguments developed in the proof of Lemma \ref{lemma_auxiliary_error_fd} show that $w$ satisfies the sign conditions in \eqref{eq:cone_condition} and that $\bar{\mathfrak{d}}_h:= \alpha \bar{\mathsf{u}}_h - \bar{\mathbf{y}}_h \cdot \bar{\mathbf{z}}_h \to \alpha \bar{u} - \bar{\mathbf{y}} \cdot \bar{\mathbf{z}} = \bar{\mathfrak{d}}$ in $L^2(\Omega)$ as $h \to 0$. We now set $\mathsf{u}_h = \bar{u}$ in \eqref{eq:variational_inequality_semi-discrete} to obtain
\[
 \int_{\Omega} \bar{\mathfrak{d}} w \mathrm{d}x = \lim_{h \to 0} \int_{\Omega} \bar{\mathfrak{d}}_h w_h \mathrm{d}x = \lim_{h \to 0} \|\bar{\mathsf{u}}_h-\bar{u}\|^{-1}_{L^2(\Omega)} \int_{\Omega} (\alpha \bar{\mathsf{u}}_h - \bar{\mathbf{y}}_h \cdot \bar{\mathbf{z}}_h)(\bar{\mathsf{u}}_h-\bar{u})  \mathrm{d}x \leq 0.
\]
As in the proof of Lemma \ref{lemma_auxiliary_error_fd}, this shows that if $\bar{\mathfrak{d}}(x)\neq 0$, then $w(x)=0$ for a.e.~$x\in\Omega$. This proves that $w\in C_{\bar{u}}$. The rest of the proof follows the arguments from the proof of Lemma \ref{lemma_auxiliary_error_fd}. For brevity, we omit the details.
\end{proof}

We now derive the error bound \eqref{eq:error_estimate_sd}. To this end, we present the following result which guarantees that discrete solutions to problem \eqref{eq:brinkamn_problem_state_sd} are uniformly bounded in $\mathbf{W}^{1,\kappa}(\Omega) \times L^{\kappa}(\Omega)$ for some $\kappa > 4$ when $d=2$ and $\kappa >3$ when $d=3$. Note that as a consequence of a standard Sobolev embedding, we have that such discrete velocities are uniformly bounded with respect to $h$ in $\mathbf{L}^{\infty}(\Omega)$: $\|\bar{\mathbf{y}}_h \|_{\mathbf{L}^{\infty}(\Omega)} \lesssim \| \mathbf{f} \|_{\mathbf{L}^2(\Omega)}$.

\begin{proposition}[discrete stability]
\label{pro:discrete_stability}
Let $\boldsymbol{\mathbf{f}} \in \mathbf{L}^2(\Omega)$, $\Omega$ be a convex polytope, and $\T_h$ be quasi-uniform. Then, there exists $\kappa > 4$ if $d=2$ and $\kappa >3$ if $d=3$ so that
 \begin{equation}
  \| \nabla \bar{\mathbf{y}}_h \|_{\mathbf{L}^{\kappa}(\Omega)} 
  +
  \| \bar{\mathsf{p}}_h \|_{L^\kappa(\Omega)}
  \lesssim \| \mathbf{f} \|_{\mathbf{L}^2(\Omega)}.
 \end{equation}
\end{proposition}
\begin{proof}
 The proof follows directly from \cite[Corollaries 4 and 5]{MR3422453} for $d=3$ and \cite[Theorems 8.2 and 8.4]{MR2121575} for $d=2$ combined with the regularity results of Theorem \ref{theorem_regularity_Linfty}.
\end{proof}


\begin{theorem}[a priori error bound]
\label{theorem_estimates_control_semi}
Let $\boldsymbol{\mathbf{f}}$ and $\mathbf{y}_{\Omega}$ be in $\mathbf{L}^2(\Omega)$. Let $\Omega \subset \mathbb{R}^d$ be a convex polytope, and let $\T_h$ be a quasi-uniform mesh. Let us assume that $\bar{u} \in \mathbb{U}_{ad}$ satisfies the optimality conditions \eqref{equivalence_condition}. Then, there exists $h_{\dagger}>0$ such that
\begin{align}\label{eq:error_estimate_sd}
\|\bar{u}-\bar{\mathsf{u}}_{h}\|_{L^2(\Omega)}\lesssim h^2\|\mathbf{f}\|_{\mathbf{L}^2(\Omega)}\left[\|\mathbf{f}\|_{\mathbf{L}^2(\Omega)}+\|\mathbf{y}_{\Omega}\|_{\mathbf{L}^2(\Omega)}\right]\quad\forall h\in(0, h_{\dagger}).
\end{align}
\end{theorem}
\begin{proof}
Set $u =\bar{\mathsf{u}}_h$ in \eqref{eq:ineq_variational} and $u=\bar{u}$ in \eqref{eq:variational_inequality_semi-discrete} to arrive at $-j'(\bar{u})(\bar{\mathsf{u}}_h-\bar{u})\leq 0$ and $j'_h(\bar{\mathsf{u}}_h)(\bar{u}-\bar{\mathsf{u}}_h)\geq 0$, respectively. With these estimates in hand, we use the instrumental error bound \eqref{eq:auxiliary_error_estimate_sd} to obtain
\begin{multline*}\label{eq:estimates_aux_error_semi-discrete}
\|\bar{u}-\bar{\mathsf{u}}_h\|_{L^2(\Omega)}^2
\lesssim [j'(\bar{\mathsf{u}}_h)-j_h'(\bar{\mathsf{u}}_h)](\bar{\mathsf{u}}_h-\bar{u})
=(\bar{\mathbf{y}}_h \cdot \bar{\mathbf{z}}_h-\hat{\mathbf{y}}\cdot\hat{\mathbf{z}},\bar{\mathsf{u}}_h-\bar{u})_{L^2(\Omega)}
\\
=(\hat{\mathbf{z}}\cdot(\bar{\mathbf{y}}_h-\hat{\mathbf{y}}),\bar{\mathsf{u}}_h-\bar{u})_{L^2(\Omega)}
+(\bar{\mathbf{y}}_h\cdot (\bar{\mathbf{z}}_h-\hat{\mathbf{z}}),\bar{\mathsf{u}}_h-\bar{u})_{L^2(\Omega)}
=:\text{I}_h +\text{II}_h.
\end{multline*}
Here, $(\hat{\mathbf{y}},\hat{\mathsf{p}})\in \mathbf{H}_0^1(\Omega)\times L_0^2(\Omega)$ and $(\hat{\mathbf{z}},\hat{\mathsf{r}})\in \mathbf{H}_0^1(\Omega)\times L_0^2(\Omega)$ correspond to the solutions of problems \eqref{eq:aux_problem_adj_1_error} and \eqref{eq:aux_problem_adj_2_error}, respectively, where $\bar{u}_h$ is replaced by $\bar{\mathsf{u}}_h$.

We now control the term $\text{I}_h$. To do this, we proceed as follows:
\begin{multline}\label{eq:estimates_I1_sd}
\text{I}_h \leq \| \hat{\mathbf{z}} \|_{\mathbf{L}^{\infty}(\Omega)}\|\bar{\mathbf{y}}_h-\hat{\mathbf{y}}\|_{\mathbf{L}^2(\Omega)}\|\bar{\mathsf{u}}_h-\bar{u}\|_{L^2(\Omega)}
\\
\lesssim h^2\|\mathbf{f}\|_{\mathbf{L}^2(\Omega)}\left[\|\mathbf{f}\|_{\mathbf{L}^2(\Omega)}+\|\mathbf{y}_{\Omega}\|_{\mathbf{L}^2(\Omega)}\right]\|\bar{\mathsf{u}}_h-\bar{u}\|_{L^2(\Omega)},
\end{multline}
where we used $\|\bar{\mathbf{y}}_h-\hat{\mathbf{y}}\|_{\mathbf{L}^2(\Omega)} \lesssim h^2 ( \| \hat{\mathbf{y}} \|_{\mathbf{H}^{2}(\Omega)} + \| \hat{\mathsf{p}} \|_{H^1(\Omega)}) \lesssim h^2 \| \mathbf{f} \|_{\mathbf{L}^{2}(\Omega)}$, which follows from \cite[Proposition 41.8]{Guermond-Ern} (see also \cite[Theorem 4.21]{Guermond-Ern} and \cite[Theorem 4.26]{Guermond-Ern}), and the estimate $ \| \hat{\mathbf{z}} \|_{\mathbf{L}^{\infty}(\Omega)} \lesssim \| \hat{\mathbf{y}} \|_{\mathbf{L}^{2}(\Omega)} + \| \mathbf{y}_{\Omega} \|_{\mathbf{L}^{2}(\Omega)}$, which follows from Theorem \ref{theorem_regularity_Linfty}.

We now control $\text{II}_h$. To this end, we introduce $(\tilde{\mathbf{z}},\tilde{\mathsf{r}})\in \mathbf{H}_0^1(\Omega)\times L_0^2(\Omega)$ as the solution to \eqref{eq:aux_problem_adj_3_error}, where $\bar{u}_h$ is replaced by $\mathsf{\bar{u}}_h$. With these ingredients in hand, we obtain: $\text{II}_h \leq \|\bar{\mathbf{y}}_h\|_{\mathbf{L}^{\infty}(\Omega)}\left(\| \hat{\mathbf{z}}-\tilde{\mathbf{z}} \|_{\mathbf{L}^2(\Omega)}+\|\tilde{\mathbf{z}}-\bar{\mathbf{z}}_h\|_{\mathbf{L}^2(\Omega)}\right)\|\bar{\mathsf{u}}_h-\bar{u}\|_{L^2(\Omega)}.$ To control $\| \hat{\mathbf{z}}-\tilde{\mathbf{z}} \|_{\mathbf{L}^2(\Omega)}$, we use a stability estimate for the problem that $(\hat{\mathbf{z}}-\tilde{\mathbf{z}},\hat{\mathsf{r}}-\tilde{\mathsf{r}}) \in \mathbf{H}_0^1(\Omega)\times L_0^2(\Omega)$ solves: $\|\hat{\mathbf{z}}-\tilde{\mathbf{z}}\|_{\mathbf{L}^2(\Omega)}\lesssim \|\hat{\mathbf{y}}-\bar{\mathbf{y}}_h \|_{\mathbf{L}^2(\Omega)}\lesssim h^2 \|\mathbf{f}\|_{\mathbf{L}^2(\Omega)}$. An estimate for $\| \tilde{\mathbf{z}}-\bar{\mathbf{z}}_h \|_{\mathbf{L}^2(\Omega)}$ follows from \cite[Proposition 41.8]{Guermond-Ern}: $\| \tilde{\mathbf{z}}-\bar{\mathbf{z}}_h \|_{\mathbf{L}^2(\Omega)}\lesssim h^2(\|\mathbf{f}\|_{\mathbf{L}^2(\Omega)}+\|\mathbf{y}_{\Omega}\|_{\mathbf{L}^2(\Omega)})$. On the other hand, Proposition \ref{pro:discrete_stability} guarantees that
$\| \bar{\mathbf{y}}_h\|_{\mathbf{L}^{\infty}(\Omega)} \lesssim \| \mathbf{f} \|_{\mathbf{L}^{2}(\Omega)}$. We thus obtain
\begin{align}\label{eq:estimates_I2_sd}
\text{II}_h \lesssim h^2 \| \mathbf{f} \|_{\mathbf{L}^{2}(\Omega)} [\|\mathbf{f}\|_{\mathbf{L}^2(\Omega)}+\|\mathbf{y}_{\Omega}\|_{\mathbf{L}^2(\Omega)}]\|\bar{\mathsf{u}}_h-\bar{u}\|_{L^2(\Omega)}.
\end{align}

The desired estimate \eqref{eq:error_estimate_sd} results from a combination of \eqref{eq:estimates_I1_sd} and \eqref{eq:estimates_I2_sd}. 
\end{proof}

We conclude this section with the following error estimates.
\begin{corollary}
Let the assumptions of Theorem \ref{theorem_estimates_control_semi} hold. Then, there exists $h_{\dagger}>0$ such that we have the following error bounds for every  $h<h_{\dagger}$:
\begin{align*}\label{eq:estimates_error_semi_allvariables}
\|\nabla(\bar{\mathbf{y}}-\bar{\mathbf{y}}_{h})\|_{\mathbf{L}^2(\Omega)}
+
\|\bar{\mathsf{p}}-\bar{\mathsf{p}}_{h}\|_{L^2(\Omega)}\lesssim h\|\mathbf{f}\|_{\mathbf{L}^2(\Omega)}(1+\|\mathbf{f}\|_{\mathbf{L}^2(\Omega)}+\|\mathbf{y}_{\Omega}\|_{\mathbf{L}^2(\Omega)}),
\\
\|\nabla(\bar{\mathbf{z}}-\bar{\mathbf{z}}_{h})\|_{\mathbf{L}^2(\Omega)}+\|\bar{\mathsf{r}}-\bar{\mathsf{r}}_{h}\|_{L^2(\Omega)}
\lesssim
h( \|\mathbf{f}\|_{\mathbf{L}^2(\Omega)} + 1 )(\|\mathbf{f}\|_{\mathbf{L}^2(\Omega)}+\|\mathbf{y}_{\Omega}\|_{\mathbf{L}^2(\Omega)}).
\end{align*}
In both bounds the hidden constants do not depend on $\mathbf{f}$, $\mathbf{y}_{\Omega}$, and $h$.
\end{corollary}

\begin{proof}
The proof follows directly from a combination of the bounds given in Propositions \ref{theorem_convergence_01} and \ref{eq:theorem_adjoint_aux_estimates} and Theorem \ref{theorem_estimates_control_semi}. This concludes the proof.
\end{proof}

\section{A posteriori error bounds}
In this section, we design a posteriori error estimators for the fully and semidiscrete methods presented in \S\ref{fully_discrete_framework} and \S\ref{semi_discrete_framework}, and obtain global reliability bounds.

%%%%%%%%%%%%%%%%%%%%%%%%%%%%%%%%%%%%%%%%%%%%%%%%%%%%%%%%%%
%%%%%%%%%%%%%%%%%%%%%%%%%%%%%%%%%%%%%%%%%%%%%%%%%%%%%%%%%%
%%%%%%%%%%%%%%%%%%%%%%%%%%%%%%%%%%%%%%%%%%%%%%%%%%%%%%%%%%
%%%%%%%%%%%%%%%%%%%%%%%%%%%%%%%%%%%%%%%%%%%%%%%%%%%%%%%%%%

\subsection{An auxiliary result}
To present an a posteriori error estimator for the fully and semidiscrete methods, we introduce the auxiliary variable $\breve{u}$ as follows:
\begin{align}\label{eq:representation_aux}
\breve{u}:= \Pi_{[\mathsf{a},\mathsf{b}]}(\alpha^{-1}\bar{\mathbf{y}}_{h}\cdot\bar{\mathbf{z}}_{h}).
\end{align}
Here, $(\bar{\mathbf{y}}_{h},\bar{\mathsf{p}}_{h})$ denotes the solution to \eqref{eq:brinkamn_problem_state_fd}, where $\mathbf{u}_h$ is replaced by $\bar{\mathbf{u}}_h$, and $(\bar{\mathbf{z}}_{h},\bar{\mathsf{r}}_{h})$ denotes the solution to \eqref{eq:brinkman_problem_adjoint_fully}. We note that $\breve{u}\in\mathbb{U}_{ad}$ satisfies \cite[Lemma 2.26]{MR2583281}
\begin{eqnarray}\label{eq:variational_inequality_discrete_aux}
(\alpha \breve{u}-\bar{\mathbf{y}}_{h}\cdot \bar{\mathbf{z}}_{h},u-\breve{u})_{L^2(\Omega)}\geq 0\quad\forall u\in \mathbb{U}_{ad}.
\end{eqnarray}

The following bound is crucial for our a posteriori error analysis and extends \cite{MR4450052}.

\begin{theorem}[auxiliary bound] 
Let $\bar{u}$ be a local solution of \eqref{eq:min_functional}--\eqref{eq:brinkamn_problem_state} such that \eqref{equivalence_condition} holds for $\tau > 0$, and let 
$(\bar{\mathbf{y}},\bar{\mathsf{p}})$ and $(\bar{\mathbf{z}},\bar{\mathsf{r}})$ be the corresponding state and adjoint state variables. Let $\bar{u}_{h}$ be a local minimum of the fully discrete scheme, and let $(\bar{\mathbf{y}}_h,\bar{\mathsf{p}}_h)$ and $(\bar{\mathbf{z}}_h,\bar{\mathsf{r}}_h)$ be the corresponding discrete state and adjoint state variables. If
\begin{align}\label{eq:bound_aux}
\|\bar{\mathbf{y}}\cdot \bar{\mathbf{z}}-\bar{\mathbf{y}}_h\cdot \bar{\mathbf{z}}_h\|_{L^2(\Omega)}\leq \alpha \mu (2\mathfrak{D})^{-1},
\qquad
\|\bar{\mathbf{y}}\cdot \bar{\mathbf{z}}-\bar{\mathbf{y}}_h\cdot \bar{\mathbf{z}}_h\|_{L^{\infty}(\Omega)} < 2^{-1} \tau,
\end{align}
then we have the following error bound:
\begin{align}\label{eq:bound_aux_j_final}
\mu \|\bar{u}-\breve{u}\|_{L^2(\Omega)}^2 \leq 2(j'(\breve{u})-j'(\bar{u}))(\breve{u}-\bar{u}),
\end{align}
where the constants $\mu$ and $\mathfrak{D}$ are given as in \eqref{equivalence_condition} and \eqref{eq:lemma_j}, respectively.
\end{theorem}
\begin{proof} We divide the proof into two steps.

\emph{Step 1.} \emph{The bound \eqref{eq:bound_aux_j_final}}. Let us assume for the moment that $\breve{u}-\bar{u} \in C_{\bar u}^{\tau}$. We can therefore set $v = \breve{u}-\bar{u}$ in \eqref{equivalence_condition} and use the mean value theorem to obtain
\[
\mu \| \breve{u}-\bar{u} \|_{L^2(\Omega)}^2 \leq j''(\bar{u})(\breve{u}-\bar{u})^2 = (j'(\breve{u})-j'(\bar{u}))(\breve{u}-\bar{u})+(j''(\bar{u})-j''(\xi))(\breve{u}-\bar{u})^2,
\]
where $\xi=\bar{u}+\theta (\breve{u}-\bar{u})$ and $\theta \in (0,1)$. We now use \eqref{eq:lemma_j}, the projection formulas \eqref{eq:representation} and \eqref{eq:representation_aux}, the Lipschitz property of $\Pi_{[\mathsf{a},\mathsf{b}]}$, and the assumption \eqref{eq:bound_aux} to obtain
\begin{multline}
 (j''(\bar{u})-j''(\xi))(\breve{u}-\bar{u})^2 \leq \mathfrak{D} \| \breve{u} -\bar{u} \|_{L^2(\Omega)} \|\breve{u}-\bar{u}\|_{L^2(\Omega)}^2
 \\
 \leq \mathfrak{D} \alpha^{-1} \|\bar{\mathbf{y}}\cdot \bar{\mathbf{z}} - \bar{\mathbf{y}}_h\cdot \bar{\mathbf{z}}_h  \|_{L^2(\Omega)} \|\breve{u}-\bar{u}\|_{L^2(\Omega)}^2 \leq (\mu/2) \|\breve{u}-\bar{u}\|_{L^2(\Omega)}^2.
\end{multline}
If this bound is replaced in the one derived for $\mu \| \breve{u}-\bar{u} \|_{L^2(\Omega)}^2$, \eqref{eq:bound_aux_j_final} is obtained.

\emph{Step 2.} $\breve{u}-\bar{u} \in C_{\bar u}^{\tau}$. Since by definition $\breve{u} \in \mathbb{U}_{ad}$, we can conclude that $\breve{u}-\bar{u}$ satisfies the conditions in \eqref{eq:cone_condition}. To prove the remaining condition in \eqref{eq:cone}, we proceed as follows: Let $x\in \Omega$ be such that $\bar{\mathfrak{d}}(x) = \alpha \bar{u}(x) - \bar{\mathbf{y}}(x) \cdot \bar{\mathbf{z}}(x) > \tau$. The fact that $\tau > 0$ guarantees that $\bar{u}(x) > \alpha^{-1} \bar{\mathbf{y}}(x) \cdot \bar{\mathbf{z}}(x)$. The projection formula \eqref{eq:representation} therefore shows that $\bar{u}(x) = \mathfrak{a}$. On the other hand, we note that assumption \eqref{eq:bound_aux} shows that 
\[
 \| (\alpha \bar{u} - \bar{\mathbf{y}} \cdot \bar{\mathbf{z}}) - ( \alpha \breve{u}- \bar{\mathbf{y}}_h \cdot \bar{\mathbf{z}}_h )  \|_{L^{\infty}(\Omega)} 
 \leq 2 \| \bar{\mathbf{y}} \cdot \bar{\mathbf{z}} -  \bar{\mathbf{y}}_h \cdot \bar{\mathbf{z}}_h  \|_{L^{\infty}(\Omega)} < \tau.
\]
Since $\alpha \bar{u}(x) - \bar{\mathbf{y}}(x) \cdot \bar{\mathbf{z}}(x) > \tau$, we can conclude that $\alpha \breve{u}(x)- \bar{\mathbf{y}}_h(x) \cdot \bar{\mathbf{z}}_h(x) >0$ and therefore that $\breve{u}(x) > \alpha^{-1} \bar{\mathbf{y}}_h(x) \cdot \bar{\mathbf{z}}_h(x)$. This implies that $\breve{u}(x) = \mathfrak{a}$. As a result, we have proved that $(\breve{u} - \bar{u})(x) = 0$ if $x \in \Omega$ is such that $\bar{\mathfrak{d}}(x) > \tau$. Similar calculations show that $(\breve{u} - \bar{u})(x) = 0$ if $x \in \Omega$ is such that $\bar{\mathfrak{d}}(x) < -\tau$. This concludes the proof.
\end{proof}

%%%%%%%%%%%%%%%%%%%%%%%%%%%%%%%%%%%%%%%%%%%%%%%%%%%%%%%%%%
%%%%%%%%%%%%%%%%%%%%%%%%%%%%%%%%%%%%%%%%%%%%%%%%%%%%%%%%%%
%%%%%%%%%%%%%%%%%%%%%%%%%%%%%%%%%%%%%%%%%%%%%%%%%%%%%%%%%%
%%%%%%%%%%%%%%%%%%%%%%%%%%%%%%%%%%%%%%%%%%%%%%%%%%%%%%%%%%

\subsection{Reliability analysis: the fully discrete scheme} 
\label{Real:fully_discrete} 
In this section, we propose an a posteriori error estimator for the fully discrete scheme and obtain a global reliability estimate. To this end, we introduce the following auxiliary variables. First, we introduce $(\check{\mathbf{y}},\check{\mathsf{p}})\in \mathbf{H}_0^1(\Omega)\times L_0^2(\Omega)$ as the solution of
\begin{equation}\label{eq:brinkamn_problem_state_aux}
(\nabla \check{\mathbf{y}},\nabla \mathbf{v})_{\mathbf{L}^2(\Omega)}+( \bar{u}_{h}\check{\mathbf{y}}, \mathbf{v})_{\mathbf{L}^2(\Omega)}-(\check{\mathsf{p}},\text{div }\mathbf{v})_{L^2(\Omega)} = 	(\mathbf{f},\mathbf{v})_{\mathbf{L}^2(\Omega)},
\end{equation}
and $(\mathsf{q},\text{div }\check{\mathbf{y}})_{L^2(\Omega)}  = 0$ for all $\mathbf{v} \in \mathbf{H}_0^1(\Omega)$ and $\mathsf{q} \in L_0^2(\Omega)$, respectively. We note that problem \eqref{eq:brinkamn_problem_state_aux} is well-posed. We also note that $(\bar{\mathbf{y}}_{h},\bar{\mathsf{p}}_{h})$, which solves \eqref{eq:brinkamn_problem_state_fd} where $u_{h}$ is replaced by $\bar{u}_{h}$, corresponds to the finite element approximation of $(\check{\mathbf{y}},\check{\mathsf{p}})$ in the setting of \S\ref{fem}. Therefore, it is natural to define the following \emph{a posteriori error estimator}:
\begin{equation}\label{eq:estimator_st}
\mathcal{E}_{st,\T}^2:= \sum_{K\in\T}\mathcal{E}_{st,K}^2,
\end{equation}
where $\mathcal{E}_{st,K}^2:=h_K^2  \|\mathcal{R}_{K}^{st}\|_{\mathbf{L}^2(K)}^2
+h_K  \|\mathcal{J}_{\gamma}^{st}\|_{\mathbf{L}^2(\partial K\setminus\partial \Omega)}^2
+\|\text{div }\bar{\mathbf{y}}_{h}\|_{L^2(K)}^2$. Here,
\begin{align}\label{eq:residual_state}
\mathcal{R}_{K}^{st}:=(\mathbf{f}+\Delta \bar{\mathbf{y}}_{h}-\bar{u}_{h}\bar{\mathbf{y}}_{h}-\nabla\bar{\mathsf{p}}_{h})|_K,
\qquad 
\mathcal{J}_{\gamma}^{st}:=\llbracket(\nabla \bar{\mathbf{y}}_{h}-\bar{\mathsf{p}}_{h}\mathbf{I})\cdot \mathbf{n}\rrbracket.
\end{align}
An application of Theorem \ref{eq:reliability}, with $\mathfrak{u}$ and $\boldsymbol{\mathfrak{f}}$ replaced by $\bar{u}_{h}$ and $\mathbf{f}$, shows that
\begin{align}\label{eq:bound_realiability_1}
\|\nabla(\check{\mathbf{y}}-\bar{\mathbf{y}}_{h})\|_{\mathbf{L}^2(\Omega)}^2+\|\check{\mathsf{p}}-\bar{\mathsf{p}}_{h}\|_{L^2(\Omega)}^2\lesssim \mathcal{E}_{st,\T}^2.
\end{align}

We also introduce the pair $(\check{\mathbf{z}},\check{\mathsf{r}})\in\mathbf{H}_0^1(\Omega)\times L_0^2(\Omega)$ as the solution of
\begin{equation}
\label{eq:brinkamn_problem_adjoint_aux}
(\nabla \mathbf{v},\nabla \check{\mathbf{z}})_{\mathbf{L}^2(\Omega)}+( \bar{u}_{h}\check{\mathbf{z}}, \mathbf{v})_{\mathbf{L}^2(\Omega)}+(\check{\mathsf{r}},\text{div }\mathbf{v})_{L^2(\Omega)} = 	(\bar{\mathbf{y}}_{h}-\mathbf{y}_{\Omega},\mathbf{v})_{\mathbf{L}^2(\Omega)}
\end{equation}
and $(\mathsf{q},\text{div }\check{\mathbf{z}})_{L^2(\Omega)}  = 0$ for all $\mathbf{v} \in  \mathbf{H}_0^1(\Omega)$ and $\mathsf{q} \in L_0^2(\Omega)$, respectively. We note that \eqref{eq:brinkamn_problem_adjoint_aux} is well-posed. Since $(\bar{\mathbf{z}}_{h},\bar{\mathsf{r}}_{h})$ corresponds to the finite element approximation of $(\check{\mathbf{z}},\check{\mathsf{r}})$ in the setting of \S \ref{fem}, we introduce the following a posteriori error estimator:
% On the other hand, we indicate that the pair $(\bar{\mathbf{z}}_{h},\bar{\mathsf{r}}_{h})$, which solves \eqref{eq:brinkman_problem_adjoint_fully}, can be seen as the finite element approximation, within the finite element spaces $\mathbf{V}_h$ and $\mathsf{P}_h$ defined in \eqref{ME:vel_space}--\eqref{TH:press_space}, for $(\check{\mathbf{z}},\check{\mathsf{r}})\in\mathbf{H}_0^1(\Omega)\times L_0^2(\Omega)$. 
% We now present the following a posteriori error estimator for $K\in\T$. 
\begin{equation}
\label{eq:estimator_AD}
\mathcal{E}_{adj,\T}^2:= \sum_{K\in\T}\mathcal{E}_{adj,K}^2, 
% \quad \displaystyle\mathcal{E}_{adj,K}^2:=h_K^2  \|\mathcal{R}_{K}^{adj}\|_{\mathbf{L}^2(K)}^2\\+h_K  \|\mathcal{J}_{\gamma}^{adj}\|_{\mathbf{L}^2(\partial K\setminus\partial \Omega)}^2+\|\text{div }\bar{\mathbf{z}}_{h}\|_{L^2(K)}^2.
\end{equation}
where $\mathcal{E}_{adj,K}^2:=h_K^2  \|\mathcal{R}_{K}^{adj}\|_{\mathbf{L}^2(K)}^2+h_K  \|\mathcal{J}_{\gamma}^{adj}\|_{\mathbf{L}^2(\partial K\setminus\partial \Omega)}^2 + \|\text{div }\bar{\mathbf{z}}_{h}\|_{L^2(K)}^2$. Here,
% Here, we consider the element residual $\mathcal{R}_{K}^{adj}$ and the interelement residual $\mathcal{J}_{\gamma}^{adj}$ as
\begin{align}\label{eq:residual_adjoint}
\mathcal{R}_{K}^{adj} := (\bar{\mathbf{y}}_{h}-\mathbf{y}_{\Omega}+\Delta \bar{\mathbf{z}}_{h}\!-\bar{u}_{h}\bar{\mathbf{z}}_{h} + \nabla\bar{\mathsf{r}}_{h})|_K,\qquad 
\mathcal{J}_{\gamma}^{adj} := \llbracket(\nabla \bar{\mathbf{z}}_{h}+\bar{\mathsf{r}}_{h}\mathbf{I})\cdot \mathbf{n}\rrbracket.
\end{align}
An application of Theorem \ref{eq:reliability}, with $\mathfrak{u}$ and $\boldsymbol{\mathfrak{f}}$ replaced by $\bar{u}_{h}$ and $\bar{\mathbf{y}}_{h}-\mathbf{y}_{\Omega}$, yields
\begin{eqnarray}\label{eq:bound_realiability_2}
\|\nabla(\check{\mathbf{z}}-\bar{\mathbf{z}}_{h})\|_{\mathbf{L}^2(\Omega)}^2+\|\check{\mathsf{r}}-\bar{\mathsf{r}}_{h}\|_{L^2(\Omega)}^2\lesssim \mathcal{E}_{adj,\T}^2.
\end{eqnarray}

After introducing a posteriori error estimators associated with the discretization of the state and adjoint equations, we define the following a posteriori local error indicator and error estimator associated with the discretization of the control variable:
\begin{eqnarray}\label{eq:estimator_control}
\mathcal{E}_{ct,K}:=\|\breve{u}-\bar{u}_{h}\|_{L^2(K)}, \qquad \mathcal{E}_{ct,\T}^2:= \sum_{K\in\T}\mathcal{E}_{ct,K}^2 = \|\breve{u}-\bar{u}_{h}\|_{L^2(\Omega)}^{2}.
\end{eqnarray}

Finally, we introduce an a posteriori error estimator for the fully discrete scheme, which is the sum of three contributions:
\begin{align}\label{eq:error_contributions}
\mathcal{E}_{ocp,\T}^2:=\mathcal{E}_{st,\T}^2 + \mathcal{E}_{adj,\T}^2 + \mathcal{E}_{ct,\T}^2.
\end{align}
% where the estimators $\mathcal{E}_{st,\T}$, $\mathcal{E}_{adj,\T}$, and $\mathcal{E}_{ct,\T}$ are defined as in \eqref{eq:estimator_st}, \eqref{eq:estimator_AD}, and \eqref{eq:estimator_control}, respectively.

To simplify the presentation and proof of the next result, we define the errors $\mathbf{e}_{\bar{\mathbf{y}}}:=\bar{\mathbf{y}}-\bar{\mathbf{y}}_{h}$, $e_{\bar{\mathsf{p}}}:=\bar{\mathsf{p}}-\bar{\mathsf{p}}_{h}$, $\mathbf{e}_{\bar{\mathbf{z}}}:=\bar{\mathbf{z}}-\bar{\mathbf{z}}_{h}$, $e_{\bar{\mathsf{r}}}:=\bar{\mathsf{r}}-\bar{\mathsf{r}}_{h}$, and $e_{\bar{u}}:=\bar{u}-\bar{u}_{h}$. We also define $(\breve{\mathbf{y}},\breve{\mathsf{p}})\in \mathbf{H}_0^1(\Omega)\times L_0^2(\Omega)$ as the solution of
\begin{equation}\label{eq:brinkamn_problem_state_auxA}
(\nabla \breve{\mathbf{y}},\nabla \mathbf{v})_{\mathbf{L}^2(\Omega)}+( \breve{u}\breve{\mathbf{y}}, \mathbf{v})_{\mathbf{L}^2(\Omega)}-(\breve{\mathsf{p}},\text{div }\mathbf{v})_{L^2(\Omega)} =(\mathbf{f},\mathbf{v})_{\mathbf{L}^2(\Omega)}
% \\
% (\mathsf{q},\text{div }\breve{\mathbf{y}})_{L^2(\Omega)}  = 0\quad\forall (\mathbf{v},\mathsf{q})\in  \mathbf{H}_0^1(\Omega)\times L_0^2(\Omega),
\end{equation}
and $(\mathsf{q},\text{div }\breve{\mathbf{y}})_{L^2(\Omega)}  = 0$ for all $\mathbf{v} \in \mathbf{H}_0^1(\Omega)$ and $\mathsf{q}\in L_0^2(\Omega)$, respectively. Finally, we define
$(\breve{\mathbf{z}},\breve{\mathsf{r}})\in\mathbf{H}_0^1(\Omega)\times L_0^2(\Omega)$ as the solution to
\begin{equation}\label{eq:brinkamn_problem_adjoint_auxA}
(\nabla \breve{\mathbf{z}},\nabla \mathbf{v})_{\mathbf{L}^2(\Omega)}+( \breve{u}\breve{\mathbf{z}}, \mathbf{v})_{\mathbf{L}^2(\Omega)}+(\breve{\mathsf{r}},\text{div }\mathbf{v})_{L^2(\Omega)} = 	(\breve{\mathbf{y}}-\mathbf{y}_{\Omega},\mathbf{v})_{\mathbf{L}^2(\Omega)}
% \\
% (\mathsf{q},\text{div }\breve{\mathbf{z}})_{L^2(\Omega)}  = 0 \quad\forall (\mathbf{v},\mathsf{q})\in  \mathbf{H}_0^1(\Omega)\times L_0^2(\Omega). 
\end{equation}
and $(\mathsf{q},\text{div }\breve{\mathbf{z}})_{L^2(\Omega)}  = 0$ for all $\mathbf{v} \in \mathbf{H}_0^1(\Omega)$ and $\mathsf{q} \in L_0^2(\Omega)$, respectively.


\begin{theorem}[global reliability of $\mathcal{E}_{ocp,\T}$]
\label{global_reliability_fully} 
Let $\bar{u}$ be a local solution of \eqref{eq:min_functional}--\eqref{eq:brinkamn_problem_state} such that \eqref{equivalence_condition} holds for $\tau > 0$, and let 
$(\bar{\mathbf{y}},\bar{\mathsf{p}})$ and $(\bar{\mathbf{z}},\bar{\mathsf{r}})$ be the corresponding state and adjoint state variables. Let $\bar{u}_{h}$ be a local minimum of the fully discrete scheme, and let $(\bar{\mathbf{y}}_h,\bar{\mathsf{p}}_h)$ and $(\bar{\mathbf{z}}_h,\bar{\mathsf{r}}_h)$ be the corresponding discrete state and adjoint state variables. If \eqref{eq:bound_aux} holds, then
\begin{align}\label{eq:estimates_global_reliability}
\|\nabla \mathbf{e}_{\bar{\mathbf{y}}} \|_{\mathbf{L}^2(\Omega)}^2
+
\|\nabla \mathbf{e}_{\bar{\mathbf{z}}} \|_{\mathbf{L}^2(\Omega)}^2
+
\| e_{\bar{\mathsf{p}}} \|_{L^2(\Omega)}^2
+
\| e_{\bar{\mathsf{r}}} \|_{L^2(\Omega)}^2 
+
\|e_{\bar{u}}\|_{L^2(\Omega)}^2
\lesssim \mathcal{E}_{ocp,\T}^2,
\end{align}
where the hidden constant is independent of the continuous and discrete optimal variables, the size of the elements in the mesh $\T$, and $\#\T$.
\end{theorem}
\begin{proof}
We divide the proof into six steps.

\emph{Step 1.} \emph{A bound for $\|e_{\bar{u}}\|_{L^2(\Omega)}$.} We use the variable $\breve{u}$ defined in \eqref{eq:representation_aux}, the definitions in \eqref{eq:estimator_control}, and a simple application of the triangle inequality to obtain
\begin{align}\label{eq:bound_00}
\|e_{\bar{u}}\|_{L^2(\Omega)}\leq \|\bar{u}-\breve{u}\|_{L^2(\Omega)}+\mathcal{E}_{ct,\T}.
\end{align}
It is therefore sufficient to bound $\|\bar{u}-\breve{u}\|_{L^2(\Omega)}$. To do this, we set $u=\breve{u}$ in \eqref{eq:ineq_variational} and $u=\bar{u}$ in \eqref{eq:variational_inequality_discrete_aux} to obtain $j'(\bar{u})(\breve{u}-\bar{u})\geq 0$ and $(\alpha \breve{u}-\bar{\mathbf{y}}_{h}\cdot \bar{\mathbf{z}}_{h},\bar{u}-\breve{u})_{L^2(\Omega)}\geq 0$.
%\begin{align*}
%j'(\bar{u})(\breve{u}-\bar{u})\geq 0,\quad (\alpha \breve{u}-\bar{\mathbf{y}}_{h}\cdot \bar{\mathbf{z}}_{h},\bar{u}-\breve{u})_{L^2(\Omega)}\geq 0.
%\end{align*}
With these estimates in hand, we use \eqref{eq:bound_aux_j_final} to obtain
\begin{multline}\label{eq:bound_01}
\mu\|\bar{u}-\breve{u}\|_{L^2(\Omega)}^2
\leq 
2 (j'(\breve{u})-j'(\bar{u}))(\breve{u}-\bar{u})
\leq 
2 j'(\breve{u})(\breve{u}-\bar{u})\\
= 
2 (\alpha \breve{u}-\breve{\mathbf{y}}\cdot \breve{\mathbf{z}},\breve{u}-\bar{u})_{L^2(\Omega)}
\leq 
2 (\bar{\mathbf{y}}_{h}\cdot \bar{\mathbf{z}}_{h}-\breve{\mathbf{y}}\cdot \breve{\mathbf{z}},\breve{u}-\bar{u})_{L^2(\Omega)},
\end{multline}
where $(\breve{\mathbf{y}},\breve{\mathsf{p}})$ and $(\breve{\mathbf{z}},\breve{\mathsf{r}})$ are the solutions of \eqref{eq:brinkamn_problem_state_auxA} and \eqref{eq:brinkamn_problem_adjoint_auxA}, respectively. Add and subtract $\bar{\mathbf{y}}_{h}\cdot \breve{\mathbf{z}}$ in \eqref{eq:bound_01} and use $\mathbf{H}_0^1(\Omega) \hookrightarrow \mathbf{L}^4(\Omega)$ to obtain
% to obtain the bound $\|\bar{u}-\breve{u}\|_{L^2(\Omega)}\lesssim \|\bar{\mathbf{z}}_{h}-\breve{\mathbf{z}}\|_{\mathbf{L}^4(\Omega)}\| \bar{\mathbf{y}}_{h}\|_{\mathbf{L}^4(\Omega)}+ \|\bar{\mathbf{y}}_{h}-\breve{\mathbf{y}}\|_{\mathbf{L}^4(\Omega)}\| \breve{\mathbf{z}}\|_{\mathbf{L}^4(\Omega)}$.
%\begin{align*}
%	\|\bar{u}-\breve{u}\|_{L^2(\Omega)}\lesssim \|\bar{\mathbf{z}}_{h}-\breve{\mathbf{z}}\|_{\mathbf{L}^4(\Omega)}\| \bar{\mathbf{y}}_{h}\|_{\mathbf{L}^4(\Omega)}+ \|\bar{\mathbf{y}}_{h}-\breve{\mathbf{y}}\|_{\mathbf{L}^4(\Omega)}\| \breve{\mathbf{z}}\|_{\mathbf{L}^4(\Omega)}.
%\end{align*}
% In view of a standard Sobolev embedding, we conclude that
$
	\|\bar{u} - \breve{u}\|_{L^2(\Omega)}\lesssim \|\nabla(\bar{\mathbf{z}}_{h} - \breve{\mathbf{z}})\|_{\mathbf{L}^2(\Omega)}\| \nabla\bar{\mathbf{y}}_{h}\|_{\mathbf{L}^2(\Omega)} +  \|\nabla(\bar{\mathbf{y}}_{h} - \breve{\mathbf{y}})\|_{\mathbf{L}^2(\Omega)}\| \nabla\breve{\mathbf{z}}\|_{\mathbf{L}^2(\Omega)}.
$
Note that $\| \nabla\bar{\mathbf{y}}_{h}\|_{\mathbf{L}^2(\Omega)}\lesssim \|\mathbf{f}\|_{\mathbf{L}^2(\Omega)}$ and that
% In what follows, we estimate the remaining terms $\| \nabla\bar{\mathbf{y}}_{h}\|_{\mathbf{L}^2(\Omega)}$ and $\| \nabla\breve{\mathbf{z}}\|_{\mathbf{L}^2(\Omega)}$ in \eqref{eq:bound_02}. To accomplish this task, we set $\mathbf{v}_{h}=\bar{\mathbf{y}}_{h}$ and $\mathsf{q}_h=0$ in \eqref{eq:brinkamn_problem_state_fd} and use that $\bar{u}_{h}\in\mathbb{U}_{ad}$ to obtain the estimate $\| \nabla\bar{\mathbf{y}}_{h}\|_{\mathbf{L}^2(\Omega)}\lesssim \|\mathbf{f}\|_{\mathbf{L}^2(\Omega)}$. Similar arguments yield to
% \begin{align*}
$
\|\nabla \breve{\mathbf{z}}\|_{\mathbf{L}^2(\Omega)}
% \lesssim 
% \|\breve{\mathbf{y}}-\mathbf{y}_{\Omega}\|_{\mathbf{L}^2(\Omega)}\lesssim \|\nabla\breve{\mathbf{y}}\|_{\mathbf{L}^2(\Omega)}+\|\mathbf{y}_{\Omega}\|_{\mathbf{L}^2(\Omega)}
\lesssim \|\mathbf{f}\|_{\mathbf{L}^2(\Omega)}+\|\mathbf{y}_{\Omega}\|_{\mathbf{L}^2(\Omega)}$.
% \end{align*}
From this follows
\begin{align}\label{eq:bound_03}
	\|\bar{u}-\breve{u}\|_{L^2(\Omega)}\lesssim \|\nabla(\bar{\mathbf{z}}_{h}-\breve{\mathbf{z}})\|_{\mathbf{L}^2(\Omega)}+ \|\nabla(\bar{\mathbf{y}}_{h}-\breve{\mathbf{y}})\|_{\mathbf{L}^2(\Omega)}.
\end{align}
%with a hidden constant independent of the continuous and discrete optimal variables, the size of the elements in the mesh $\T$, and $\#\T$ but depends on the continuous problem data.

Let us now bound $\|\nabla(\bar{\mathbf{y}}_{h}-\breve{\mathbf{y}})\|_{\mathbf{L}^2(\Omega)}$. For this purpose, we use the pair $(\check{\mathbf{y}},\check{\mathsf{p}})$ defined as the solution to \eqref{eq:brinkamn_problem_state_aux} and the a posteriori error bound \eqref{eq:bound_realiability_1} to obtain
\begin{align}\label{eq:bound_04}
\|\nabla(\bar{\mathbf{y}}_{h}-\breve{\mathbf{y}})\|_{\mathbf{L}^2(\Omega)}\lesssim \|\nabla(\breve{\mathbf{y}}-\check{\mathbf{y}})\|_{\mathbf{L}^2(\Omega)}+\mathcal{E}_{st,\T}.
\end{align}
% We now to control the term $\|\nabla(\breve{\mathbf{y}}-\check{\mathbf{y}})\|_{\mathbf{L}^2(\Omega)}$. 
We now observe that $(\breve{\mathbf{y}}-\check{\mathbf{y}},\breve{\mathsf{p}}-\check{\mathsf{p}})\in \mathbf{H}_0^1(\Omega)\times L_0^2(\Omega)$ solves the following problem
\begin{multline*}\label{eq:brinkamn_problem_state_aux_error}
(\nabla (\breve{\mathbf{y}} - \check{\mathbf{y}}),\nabla \mathbf{v})_{\mathbf{L}^2(\Omega)}
+
( \breve{u}(\breve{\mathbf{y}}\!-\!\check{\mathbf{y}}), \mathbf{v})_{\mathbf{L}^2(\Omega)} 
-
(\breve{\mathsf{p}}\!-\!\check{\mathsf{p}},\text{div }\mathbf{v})_{L^2(\Omega)} 
\\
=  	(\check{\mathbf{y}}(\bar{u}_h\!-\!\breve{u}),\mathbf{v})_{\mathbf{L}^2(\Omega)},\quad
(\mathsf{q},\text{div }(\breve{\mathbf{y}}-\check{\mathbf{y}}))_{L^2(\Omega)} = 0\quad\forall (\mathbf{v},\mathsf{q})\in  \mathbf{H}_0^1(\Omega)\times L_0^2(\Omega).
\end{multline*}
Set $(\mathbf{v},\mathsf{q}) = (\breve{\mathbf{y}}-\check{\mathbf{y}},0)$ and use \eqref{eq:estimator_control} to get
$\|\nabla (\breve{\mathbf{y}}-\check{\mathbf{y}})\|_{\mathbf{L}^2(\Omega)} \lesssim \|\nabla \check{\mathbf{y}}\|_{\mathbf{L}^2(\Omega)}\|\bar{u}_{h}-\breve{u}\|_{L^2(\Omega)}\lesssim \mathcal{E}_{ct,\T}$.
%invoke the fact that $\breve{\mathbf{y}}-\check{\mathbf{y}}\in \mathbf{V}(\Omega)$ to arrive at
%\begin{align*}
%\|\nabla (\breve{\mathbf{y}}-\check{\mathbf{y}})\|_{\mathbf{L}^2(\Omega)}^2+( \breve{u}(\breve{\mathbf{y}}-\check{\mathbf{y}}), \breve{\mathbf{y}}-\check{\mathbf{y}})_{\mathbf{L}^2(\Omega)} =  	(\check{\mathbf{y}}(\bar{u}_h-\breve{u}),\breve{\mathbf{y}}-\check{\mathbf{y}})_{\mathbf{L}^2(\Omega)},
%\end{align*}
%\begin{align*}
%\|\nabla (\breve{\mathbf{y}}-\check{\mathbf{y}})\|_{\mathbf{L}^2(\Omega)}\lesssim \|\nabla \check{\mathbf{y}}\|_{\mathbf{L}^2(\Omega)}\|\bar{u}_{h}-\breve{u}\|_{L^2(\Omega)}\lesssim \|\mathbf{f}\|_{\mathbf{L}^2(\Omega)}\|\bar{u}_{h}-\breve{u}\|_{L^2(\Omega)}\lesssim \mathcal{E}_{ct,\T}.
%\end{align*}
Substituting this bound into \eqref{eq:bound_04} and the one obtained in \eqref{eq:bound_03}, we obtain
\begin{align}\label{eq:bound_05}
	\|\bar{u}-\breve{u}\|_{L^2(\Omega)}\lesssim \|\nabla(\bar{\mathbf{z}}_{h}-\breve{\mathbf{z}})\|_{\mathbf{L}^2(\Omega)}+ \mathcal{E}_{ct,\T}+\mathcal{E}_{st,\T}.
\end{align}

In the following, we bound $\|\nabla(\bar{\mathbf{z}}_{h}-\breve{\mathbf{z}})\|_{\mathbf{L}^2(\Omega)}$ in \eqref{eq:bound_05}. To do this, we use the pair $(\check{\mathbf{z}},\check{\mathsf{r}})$ defined as the solution of \eqref{eq:brinkamn_problem_adjoint_aux} and the a posteriori error estimate \eqref{eq:bound_realiability_2} to obtain 
\begin{align}\label{eq:bound_06}
\|\nabla(\bar{\mathbf{z}}_{h}-\breve{\mathbf{z}})\|_{\mathbf{L}^2(\Omega)}\lesssim \mathcal{E}_{adj,\T}+\|\nabla(\breve{\mathbf{z}}-\check{\mathbf{z}})\|_{\mathbf{L}^2(\Omega)}.
\end{align}
To bound $\|\nabla(\breve{\mathbf{z}}-\check{\mathbf{z}})\|_{\mathbf{L}^2(\Omega)}$, we note that $(\breve{\mathbf{z}}-\check{\mathbf{z}},\breve{\mathsf{r}}-\check{\mathsf{r}})\in\mathbf{H}_0^1(\Omega)\times L_0^2(\Omega)$ solves
\begin{multline*}
(\nabla (\breve{\mathbf{z}}-\check{\mathbf{z}}),\nabla \mathbf{v})_{\mathbf{L}^2(\Omega)}+( \breve{u}(\breve{\mathbf{z}}-\check{\mathbf{z}}), \mathbf{v})_{\mathbf{L}^2(\Omega)}+(\breve{\mathsf{r}}-\check{\mathsf{r}},\text{div }\mathbf{v})_{L^2(\Omega)} = ((\bar{u}_{h}-\breve{u})\check{\mathbf{z}},\mathbf{v})_{\mathbf{L}^2(\Omega)}\\+(\breve{\mathbf{y}}-\bar{\mathbf{y}}_{h},\mathbf{v})_{\mathbf{L}^2(\Omega)},\quad
(\mathsf{q},\text{div }(\breve{\mathbf{z}}-\check{\mathbf{z}}))_{L^2(\Omega)}  =  0 \quad\forall (\mathbf{v},\mathsf{q})\in  \mathbf{H}_0^1(\Omega)\times L_0^2(\Omega). 
\end{multline*}
If we set $\mathbf{v} = \breve{\mathbf{z}}-\check{\mathbf{z}}$ and $\mathsf{q}=0$, we obtain
$
\|\nabla(\breve{\mathbf{z}}-\check{\mathbf{z}})\|_{\mathbf{L}^2(\Omega)}\lesssim \|\nabla\check{\mathbf{z}}\|_{\mathbf{L}^2(\Omega)}\|\bar{u}_{h}-\breve{u}\|_{L^2(\Omega)}+\|\nabla(\breve{\mathbf{y}}-\bar{\mathbf{y}}_{h})\|_{\mathbf{L}^2(\Omega)}.
$
Note that $\|\nabla \check{\mathbf{z}}\|_{\mathbf{L}^2(\Omega)} \lesssim \|\mathbf{f}\|_{\mathbf{L}^2(\Omega)}+\|\mathbf{y}_{\Omega}\|_{\mathbf{L}^2(\Omega)}$. We now use the definitions in \eqref{eq:estimator_control} and $\|\nabla (\breve{\mathbf{y}}-\bar{\mathbf{y}}_{h})\|_{\mathbf{L}^2(\Omega)}\lesssim \mathcal{E}_{st,\T}+\mathcal{E}_{ct,\T}$ to obtain 
$
\|\nabla(\breve{\mathbf{z}}-\check{\mathbf{z}})\|_{\mathbf{L}^2(\Omega)}\lesssim \mathcal{E}_{st,\T}+\mathcal{E}_{ct,\T}.
$
Substitute this estimate into \eqref{eq:bound_06} and the one obtained in \eqref{eq:bound_05} to obtain $\|\bar{u}-\breve{u}\|_{L^2(\Omega)}\lesssim \mathcal{E}_{st,\T}+\mathcal{E}_{adj,\T}+\mathcal{E}_{ct,\T}$. In view of \eqref{eq:bound_00}, we finally obtain
\begin{align}\label{eq:bound_10}
\|e_{\bar{u}}\|_{L^2(\Omega)}\lesssim \mathcal{E}_{st,\T}+\mathcal{E}_{adj,\T}+\mathcal{E}_{ct,\T}.
\end{align}

\emph{Step 2.} \emph{A bound for} $\|\nabla \mathbf{e}_{\bar{\mathbf{y}}}\|_{\mathbf{L}^2(\Omega)}$. With the auxiliary pair $(\check{\mathbf{y}},\check{\mathsf{p}})$ and the a posteriori error bound \eqref{eq:bound_realiability_1} in hand, we bound $\|\nabla \mathbf{e}_{\bar{\mathbf{y}}}\|_{\mathbf{L}^2(\Omega)}$ as follows:
\begin{align}\label{eq:bound_11}
\|\nabla \mathbf{e}_{\bar{\mathbf{y}}}\|_{\mathbf{L}^2(\Omega)}\leq \|\nabla(\bar{\mathbf{y}}-\check{\mathbf{y}})\|_{\mathbf{L}^2(\Omega)}+\mathcal{E}_{st,\T}.
\end{align}
To bound $\|\nabla(\bar{\mathbf{y}}-\check{\mathbf{y}})\|_{\mathbf{L}^2(\Omega)}$, we note that $(\bar{\mathbf{y}}-\check{\mathbf{y}},\bar{\mathsf{p}}-\check{\mathsf{p}})\times \mathbf{H}_0^1(\Omega)\times L_0^2(\Omega)$ solves 
\begin{multline}\label{eq:brinkamn_problem_state_aux_errorB}
(\nabla (\bar{\mathbf{y}}-\check{\mathbf{y}}),\nabla \mathbf{v})_{\mathbf{L}^2(\Omega)}
+
( \bar{u}(\bar{\mathbf{y}}-\check{\mathbf{y}}), \mathbf{v})_{\mathbf{L}^2(\Omega)}
-
(\bar{\mathsf{p}}-\check{\mathsf{p}},\text{div }\mathbf{v})_{L^2(\Omega)}
\\
=
((\bar{u}_h-\bar{u})\check{\mathbf{y}},\mathbf{v})_{\mathbf{L}^2(\Omega)},
\quad
(\mathsf{q},\text{div }(\bar{\mathbf{y}}-\check{\mathbf{y}}))_{L^2(\Omega)}  =  0 
\quad
\forall (\mathbf{v},\mathsf{q})\in  \mathbf{H}_0^1(\Omega)\times L_0^2(\Omega). 
\end{multline}
If we set $\mathbf{v}=\bar{\mathbf{y}}-\check{\mathbf{y}}$ and $\mathsf{q}=0$ and use \eqref{eq:bound_10}, we immediately obtain the bound
$
\|\nabla (\bar{\mathbf{y}}-\check{\mathbf{y}})\|_{\mathbf{L}^2(\Omega)}\lesssim 
\|\mathbf{f}\|_{\mathbf{L}^2(\Omega)} \|e_{\bar{u}}\|_{L^2(\Omega)}\lesssim \mathcal{E}_{st,\T}+\mathcal{E}_{adj,\T}+\mathcal{E}_{ct,\T}.
$
Finally, we replace this bound into \eqref{eq:bound_11} to obtain the a posteriori error estimate
\begin{align}\label{eq:bound_12}
\|\nabla \mathbf{e}_{\bar{\mathbf{y}}}\|_{\mathbf{L}^2(\Omega)}\lesssim \mathcal{E}_{st,\T}+\mathcal{E}_{adj,\T}+\mathcal{E}_{ct,\T}.
\end{align}

\emph{Step 3.} \emph{A bound for $\|\nabla \mathbf{e}_{\bar{\mathbf{z}}}\|_{\mathbf{L}^2(\Omega)}$}. Using similar arguments as in Step 2, we obtain
% % \EO{We recall that $(\check{\mathbf{z}},\check{\mathsf{r}})$ is defined as the solution to \eqref{eq:brinkamn_problem_adjoint_aux}. A simple application of the triangle inequality combined with \eqref{eq:bound_realiability_2} yield}
% % \begin{align}\label{eq:bound_13}
% % \|\nabla \mathbf{e}_{\bar{\mathbf{z}}}\|_{\mathbf{L}^2(\Omega)}\leq \|\nabla(\bar{\mathbf{z}}-\check{\mathbf{z}})\|_{\mathbf{L}^2(\Omega)}+\mathcal{E}_{adj,\T}.
% % \end{align}
% % To control $\|\nabla(\bar{\mathbf{z}}-\check{\mathbf{z}})\|_{\mathbf{L}^2(\Omega)}$, we note that $(\bar{\mathbf{z}}-\check{\mathbf{z}},\bar{\mathsf{r}}-\check{\mathsf{r}})\in \mathbf{H}_0^1(\Omega)\times L_0^2(\Omega)$ solves
% % \begin{multline}\label{eq:brinkamn_problem_state_aux_errorC}
% % (\nabla (\bar{\mathbf{z}}-\check{\mathbf{z}}),\nabla \mathbf{v})_{\mathbf{L}^2(\Omega)}+( \bar{u}(\bar{\mathbf{z}}-\check{\mathbf{z}}), \mathbf{v})_{\mathbf{L}^2(\Omega)}+(\bar{\mathsf{r}}-\check{\mathsf{r}},\text{div }\mathbf{v})_{L^2(\Omega)} \\=  	((\bar{u}_h-\bar{u})\check{\mathbf{z}},\mathbf{v})_{\mathbf{L}^2(\Omega)}+ (\bar{\mathbf{y}}-\bar{\mathbf{y}}_{h},\mathbf{v})_{\mathbf{L}^2(\Omega)},\quad
% % (\mathsf{q},\text{div }(\bar{\mathbf{z}}-\check{\mathbf{z}}))_{L^2(\Omega)}  =  0,
% % \end{multline}
% % for all $(\mathbf{v},\mathsf{q})\in\mathbf{H}_0^1(\Omega) \times L_0^2(\Omega)$. \EO{If we set} $\mathbf{v}=\bar{\mathbf{z}}-\check{\mathbf{z}}$ and $\mathsf{q}=0$, we obtain
% % \begin{align*}
% % \|\nabla(\bar{\mathbf{z}}-\check{\mathbf{z}})\|_{\mathbf{L}^2(\Omega)}\lesssim \|\nabla\check{\mathbf{z}}\|_{\mathbf{L}^2(\Omega)}\|e_{\bar{u}}\|_{L^2(\Omega)}+\|\mathbf{e}_{\bar{\mathbf{y}}}\|_{\mathbf{L}^2(\Omega)}.
% % \end{align*}
% % Recall that $\|\nabla \check{\mathbf{z}}\|_{\mathbf{L}^2(\Omega)} \lesssim \|\mathbf{f}\|_{\mathbf{L}^2(\Omega)}+\|\mathbf{y}_{\Omega}\|_{\mathbf{L}^2(\Omega)}$. We therefore refer to the a posteriori error bounds \eqref{eq:bound_10} and \eqref{eq:bound_12} to obtain $\|\nabla(\bar{\mathbf{z}}-\check{\mathbf{z}})\|_{\mathbf{L}^2(\Omega)}\lesssim \mathcal{E}_{st,\T}+\mathcal{E}_{adj,\T}+\mathcal{E}_{ct,\T}$. If this bound is replaced in \eqref{eq:bound_13}, the desired bound is obtained:
$
\|\nabla \mathbf{e}_{\bar{\mathbf{z}}}\|_{\mathbf{L}^2(\Omega)}\lesssim \mathcal{E}_{st,\T}+\mathcal{E}_{adj,\T}+\mathcal{E}_{ct,\T}.
$
For the sake of simplicity, we omit the details.

\emph{Step 4.} \emph{A bound for} $\|e_{\bar{\mathsf{p}}}\|_{L^2(\Omega)}$. A basic application of the bound \eqref{eq:bound_realiability_1} yields
% with the estimate \eqref{eq:bound_realiability_1} yield
\begin{align}\label{eq:bound_15}
\|e_{\bar{\mathsf{p}}}\|_{L^2(\Omega)}\lesssim \|\bar{\mathsf{p}}-\check{\mathsf{p}}\|_{L^2(\Omega)}+\mathcal{E}_{st,\T}.
\end{align}
We now control $\|\bar{\mathsf{p}}-\check{\mathsf{p}}\|_{L^2(\Omega)}$ using the inf--sup condition \eqref{eq:infsup}, the fact that $(\bar{\mathbf{y}}-\check{\mathbf{y}},\bar{\mathsf{p}}-\check{\mathsf{p}})$ solves problem \eqref{eq:brinkamn_problem_state_aux_errorB}, and the standard embedding $\mathbf{H}_0^1(\Omega) \hookrightarrow \mathbf{L}^4(\Omega)$. In fact,
% the H\"older's inequality, and a standard Sobolev embedding to arrive at
\begin{multline}\label{eq:infsup_state_aux}
\|\bar{\mathsf{p}}-\check{\mathsf{p}}\|_{L^2(\Omega)}\lesssim \sup_{\mathbf{v}\in \mathbf{H}_0^1(\Omega)}\dfrac{(\bar{\mathsf{p}}-\check{\mathsf{p}},\text{div }\mathbf{v})_{L^2(\Omega)}}{\|\nabla \mathbf{v}\|_{\mathbf{L}^2(\Omega)}}\lesssim \|\nabla (\bar{\mathbf{y}}-\check{\mathbf{y}}) \|_{\mathbf{L}^2(\Omega)}\\
+\|\nabla (\bar{\mathbf{y}}-\check{\mathbf{y}}) \|_{\mathbf{L}^2(\Omega)} \|\bar{u}\|_{L^2(\Omega)}+\|\nabla \check{\mathbf{y}}\|_{\mathbf{L}^2(\Omega)}\|\bar{u}_{h}-\bar{u}\|_{L^2(\Omega)}.
\end{multline}
We note that $\|\nabla \check{\mathbf{y}}\|_{\mathbf{L}^2(\Omega)}\lesssim \| \mathbf{f}\|_{\mathbf{L}^2(\Omega)}$ and that $\| \bar{u} \|_{L^2(\Omega)} \leq C(\mathfrak{a},\mathfrak{b},\Omega)$. We now refer to the bound $\|\nabla (\bar{\mathbf{y}}-\check{\mathbf{y}})\|_{\mathbf{L}^2(\Omega)} \lesssim \mathcal{E}_{ocp,\T}$, which was derived in Step 2, and the error estimate \eqref{eq:bound_10} to obtain $\|\bar{\mathsf{p}}-\check{\mathsf{p}}\|_{L^2(\Omega)}\lesssim \mathcal{E}_{ocp,\T}$. In view of  \eqref{eq:bound_15}, this estimate implies 
\begin{align}\label{eq:bound_16}
\|e_{\bar{\mathsf{p}}}\|_{L^2(\Omega)}\lesssim \mathcal{E}_{st,\T}+\mathcal{E}_{adj,\T}+\mathcal{E}_{ct,\T}.
\end{align}

\emph{Step 5.} \emph{A bound for $\|e_{\bar{\mathsf{r}}}\|_{L^2(\Omega)}$.} Using similar arguments as in Step 4, we obtain.
% A basic application of the triangle inequality combined with the estimate \eqref{eq:bound_realiability_2} allow us to conclude that
% \begin{align}\label{eq:bound_17}
% \|e_{\bar{\mathsf{r}}}\|_{L^2(\Omega)}\lesssim \|\bar{\mathsf{r}}-\check{\mathsf{r}}\|_{L^2(\Omega)}+\mathcal{E}_{adj,\T}.
% \end{align}
% \EO{To control $\|\bar{\mathsf{r}}-\check{\mathsf{r}}\|_{L^2(\Omega)}$ we use the inf--sup condition \eqref{eq:infsup}, the problem that $(\bar{\mathbf{z}}-\check{\mathbf{z}},\bar{\mathsf{r}}-\check{\mathsf{r}})$ solves, and the standard embedding $\mathbf{H}_0^1(\Omega) \hookrightarrow \mathbf{L}^4(\Omega)$. In fact, we have}
% \begin{multline}\label{eq:infsup_adjoint_aux}
% \|\bar{\mathsf{r}}-\check{\mathsf{r}}\|_{L^2(\Omega)}\lesssim \sup_{\mathbf{v}\in \mathbf{H}_0^1(\Omega)}\dfrac{(\bar{\mathsf{r}}-\check{\mathsf{r}},\text{div }\mathbf{v})_{L^2(\Omega)}}{\|\nabla \mathbf{v}\|_{\mathbf{L}^2(\Omega)}}\lesssim \|\nabla (\bar{\mathbf{z}}-\check{\mathbf{z}}) \|_{\mathbf{L}^2(\Omega)}\\
% +\|\nabla (\bar{\mathbf{z}}-\check{\mathbf{z}}) \|_{\mathbf{L}^2(\Omega)} \|\bar{u}\|_{L^2(\Omega)}+\|\nabla \check{\mathbf{z}}\|_{\mathbf{L}^2(\Omega)}\|\bar{u}_{h}-\bar{u}\|_{L^2(\Omega)}+\|\nabla (\bar{\mathbf{y}}-\bar{\mathbf{y}}_{h})\|_{\mathbf{L}^2(\Omega)}.
% \end{multline}
% We thus apply similar arguments as the ones developed in \emph{Step 4} to obtain the estimate $\|\bar{\mathsf{r}}-\check{\mathsf{r}}\|_{L^2(\Omega)}\lesssim \mathcal{E}_{st,\T}+\mathcal{E}_{adj,\T}+\mathcal{E}_{ct,\T}.$ Therefore, in view of the bound \eqref{eq:bound_17}, we obtain
$
\|e_{\bar{\mathsf{r}}}\|_{L^2(\Omega)}\lesssim \mathcal{E}_{st,\T}+\mathcal{E}_{adj,\T}+\mathcal{E}_{ct,\T}
$
For the sake of simplicity, we omit the details.

\emph{Step 6.} The desired estimate \eqref{eq:estimates_global_reliability} results from the combination of the bounds derived in Steps 1, 2, 3, 4, and 5. This concludes the proof.
\end{proof}

%%%%%%%%%%%%%%%%%%%%%%%%%%%%%%%%%%%%%%%%%%%%%%%%%%%%%%%%%%
%%%%%%%%%%%%%%%%%%%%%%%%%%%%%%%%%%%%%%%%%%%%%%%%%%%%%%%%%%
%%%%%%%%%%%%%%%%%%%%%%%%%%%%%%%%%%%%%%%%%%%%%%%%%%%%%%%%%%
%%%%%%%%%%%%%%%%%%%%%%%%%%%%%%%%%%%%%%%%%%%%%%%%%%%%%%%%%%

\subsection{Reliability analysis: the semidiscrete scheme} 
In this section, we propose an a posteriori error estimator for the semidiscrete scheme \eqref{eq:brinkamn_problem_state_sd}--\eqref{eq:brinkman_problem_adjoint_semi} and obtain a global reliability bound. 
% To do this, we consider the following auxiliary variables, which are also of particular importance for our reliability analysis. 
As in the previous section, we introduce $(\check{\mathbf{y}},\check{\mathsf{p}})\in \mathbf{H}_0^1(\Omega)\times L_0^2(\Omega)$ as the solution to the following problem
\begin{equation}
\label{eq:brinkamn_problem_state_aux_semi}
(\nabla \check{\mathbf{y}},\nabla \mathbf{v})_{\mathbf{L}^2(\Omega)}+(\bar{\mathsf{u}}\check{\mathbf{y}}, \mathbf{v})_{\mathbf{L}^2(\Omega)}-(\check{\mathsf{p}},\text{div }\mathbf{v})_{L^2(\Omega)} = 	(\mathbf{f},\mathbf{v})_{\mathbf{L}^2(\Omega)}
% \\
% (\mathsf{q},\text{div }\check{\mathbf{y}})_{L^2(\Omega)}  = 0,\quad \forall (\mathbf{v},\mathsf{q})\in\mathbf{H}_0^1(\Omega)\times L_0^2(\Omega).
\end{equation}
and $(\mathsf{q},\text{div }\check{\mathbf{y}})_{L^2(\Omega)}  = 0$ for all $\mathbf{v} \in \mathbf{H}_0^1(\Omega)$ and $\mathsf{q} \in L_0^2(\Omega)$, respectively.
We define the following error estimator associated to the discretization of the state equations:
\begin{equation*}\label{eq:estimator_st_semi}
\mathfrak{E}_{st,\T}^2:= \sum_{K\in\T}\mathfrak{E}_{st,K}^2, 
\end{equation*}
where $\mathfrak{E}_{st,K}^2:=h_K^2  \|\mathfrak{R}_{K}^{st}\|_{\mathbf{L}^2(K)}^2 + h_K  \|\mathfrak{J}_{\gamma}^{st}\|_{\mathbf{L}^2(\partial K\setminus\partial \Omega)}^2+\|\text{div }\bar{\mathbf{y}}_{h}\|_{L^2(K)}^2$. Here,
% where we consider the element residual $\mathfrak{R}_{K}^{st}$ and the interelement residual $\mathfrak{J}_{\gamma}^{st}$ as
\begin{align*}\label{eq:residual_state_semi}
\mathfrak{R}_{K}^{st}\ = (\mathbf{f}+\Delta \bar{\mathbf{y}}_{h} - \bar{\mathsf{u}}\bar{\mathbf{y}}_{h} - \nabla\bar{\mathsf{p}}_{h})|_K,
\qquad 
\mathfrak{J}_{\gamma}^{st} := \llbracket(\nabla \bar{\mathbf{y}}_{h}-\bar{\mathsf{p}}_{h}\mathbf{I})\cdot \mathbf{n}\rrbracket.
\end{align*}
We note that the result of Theorem \ref{eq:reliability} allows us to conclude the following a posteriori error bound:
% \begin{eqnarray}\label{eq:bound_realiability_1_semi}
$\|\nabla(\check{\mathbf{y}}-\bar{\mathbf{y}}_{h})\|_{\mathbf{L}^2(\Omega)} +\|\check{\mathsf{p}}-\bar{\mathsf{p}}_{h}\|_{L^2(\Omega)} \lesssim \mathfrak{E}_{st,\T}$.
% \end{eqnarray}

We also introduce the pair $(\check{\mathbf{z}},\check{\mathsf{r}})\in\mathbf{H}_0^1(\Omega)\times L_0^2(\Omega)$ as the solution to 
\begin{equation}\label{eq:brinkamn_problem_adjoint_aux_semi}
(\nabla \check{\mathbf{z}},\nabla \mathbf{v})_{\mathbf{L}^2(\Omega)}+( \bar{\mathsf{u}}\check{\mathbf{z}}, \mathbf{v})_{\mathbf{L}^2(\Omega)}+(\check{\mathsf{r}},\text{div }\mathbf{v})_{L^2(\Omega)} = (\bar{\mathbf{y}}_{h}-\mathbf{y}_{\Omega},\mathbf{v})_{\mathbf{L}^2(\Omega)}
\end{equation}
and $(\mathsf{q},\text{div }\check{\mathbf{z}})_{L^2(\Omega)}  =  0$ for all $\mathbf{v} \in \mathbf{H}_0^1(\Omega)$ and $\mathsf{q} \in L_0^2(\Omega)$, respectively.
We define the following error estimator associated to the discretization of the adjoint equations:
\begin{equation*}\label{eq:estimator_AD_semi}
\mathfrak{E}_{adj,\T}^2:= \sum_{K\in\T}\mathfrak{E}_{adj,K}^2, 
% \text{ with } \displaystyle\mathfrak{E}_{adj,K}^2:=h_K^2  \|\mathfrak{R}_{K}^{adj}\|_{\mathbf{L}^2(K)}^2\\+h_K  \|\mathfrak{J}_{\gamma}^{adj}\|_{\mathbf{L}^2(\partial K\setminus\partial \Omega)}^2+\|\text{div }\bar{\mathbf{z}}_{h}\|_{L^2(K)}^2,
\end{equation*}
where $\mathfrak{E}_{adj,K}^2:=h_K^2  \|\mathfrak{R}_{K}^{adj}\|_{\mathbf{L}^2(K)}^2+h_K  \|\mathfrak{J}_{\gamma}^{adj}\|_{\mathbf{L}^2(\partial K\setminus\partial \Omega)}^2+\|\text{div }\bar{\mathbf{z}}_{h}\|_{L^2(K)}^2$. Here,
% where we consider the element residual $\mathfrak{R}_{K}^{adj}$ and the interelement residual $\mathfrak{J}_{\gamma}^{adj}$ as
\begin{align*}\label{eq:residual_adjoint_semi}
\mathfrak{R}_{K}^{adj} := (\bar{\mathbf{y}}_{h}-\mathbf{y}_{\Omega}+\Delta \bar{\mathbf{z}}_{h} -\bar{\mathsf{u}}\bar{\mathbf{z}}_{h} + \nabla\bar{\mathsf{r}}_{h})|_K,
\qquad \mathfrak{J}_{\gamma}^{adj} := \llbracket(\nabla \bar{\mathbf{z}}_{h} +\bar{\mathsf{r}}_{h}\mathbf{I})\cdot \mathbf{n}\rrbracket.
\end{align*}
An application of Theorem \ref{eq:reliability} immediately yields the following a posteriori error estimate:
$
\|\nabla(\check{\mathbf{z}}-\bar{\mathbf{z}}_{h})\|_{\mathbf{L}^2(\Omega)} +\|\check{\mathsf{r}}-\bar{\mathsf{r}}_{h}\|_{L^2(\Omega)} \lesssim \mathfrak{E}_{adj,\T}.
$

To present the following result, we define $\mathsf{e}_{\bar{u}}:=\bar{u}-\bar{\mathsf{u}}$ and $\mathfrak{E}_{ocp,\T}:=\mathfrak{E}_{st,\T}+\mathfrak{E}_{adj,\T}$.

\begin{theorem}[global reliability $\mathfrak{E}_{ocp,\T}$] 
Let $\bar{u}$ be a local solution of \eqref{eq:min_functional}--\eqref{eq:brinkamn_problem_state} such that \eqref{equivalence_condition} holds for $\tau > 0$, and let 
$(\bar{\mathbf{y}},\bar{\mathsf{p}})$ and $(\bar{\mathbf{z}},\bar{\mathsf{r}})$ be the corresponding state and adjoint state variables. Let $\bar{\mathsf{u}}$ be a local minimum of the semidiscrete scheme, and let $(\bar{\mathbf{y}}_h,\bar{\mathsf{p}}_h)$ and $(\bar{\mathbf{z}}_h,\bar{\mathsf{r}}_h)$ be the corresponding discrete state and adjoint state variables. If \eqref{eq:bound_aux} holds, then
% 
% Let $(\bar{\mathbf{y}},\bar{\mathsf{p}},\bar{\mathbf{z}},\bar{\mathsf{r}},\bar{u})\in \mathbf{H}_0^1(\Omega)\times L_0^2(\Omega)\times\mathbf{H}_0^1(\Omega)\times L_0^2(\Omega)\times\mathbb{U}_{ad}$ be a local solution to \eqref{eq:min_functional}--\eqref{eq:brinkamn_problem_state} that satisfies the sufficient second order optimality condition \eqref{eq:second_order_condition}. Let $\bar{\mathsf{u}}$ be a local minimum of the associated semi-discrete optimal control problem with $(\bar{\mathbf{y}}_h,\bar{\mathsf{p}}_{h})$ and $(\bar{\mathbf{z}}_h,\bar{\mathsf{r}}_{h})$ being the corresponding state and adjoint state, respectively. If $\bar{\mathbf{y}}_h$ and $\bar{\mathbf{z}}_h$ satisfy, on the mesh $\T$, the bound \eqref{eq:bound_aux}. Then
\begin{align}\label{eq:estimates_global_reliability_semi}
\|\nabla \mathbf{e}_{\bar{\mathbf{y}}} \|_{\mathbf{L}^2(\Omega)}^2 + \|\nabla \mathbf{e}_{\bar{\mathbf{z}}} \|_{\mathbf{L}^2(\Omega)}^2 + \| e_{\bar{\mathsf{p}}} \|_{L^2(\Omega)}^2 + \| e_{\bar{\mathsf{r}}} \|_{L^2(\Omega)}^2 + \|\mathsf{e}_{\bar{u}}\|_{L^2(\Omega)}^2\lesssim \mathfrak{E}_{ocp,\T}^2,
\end{align}
where the hidden constant is independent of the continuous and discrete optimal variables, the size of the elements in the mesh $\T$, and $\#\T$.
\end{theorem}
\begin{proof}
We note that under the particular conditions entailed by the semidiscrete scheme, the auxiliary variable $\breve{u}$ defined in \eqref{eq:representation_aux} coincides with $\bar{\mathsf{u}}$. As a result, the variable $(\breve{\mathbf{y}},\breve{\mathsf{p}})$ defined as the solution \eqref{eq:brinkamn_problem_state_auxA} coincides with $(\check{\mathbf{y}},\check{\mathsf{p}})$ defined as the solution \eqref{eq:brinkamn_problem_state_aux_semi}. If we apply the same arguments as in the proof of Theorem \ref{global_reliability_fully} to this simplified situation, we can derive the bound \eqref{eq:estimates_global_reliability_semi}.  For the sake of brevity, we omit the details.
% The proof follows from an adaption of the arguments elaborated in the proof of the Theorem \ref{global_reliability_fully}.
\end{proof}

%%%%%%%%%%%%%%%%%%%%%%%%%%%%%%%%%%%%%%%%%%%%%%%%%%%%%%%%%%
%%%%%%%%%%%%%%%%%%%%%%%%%%%%%%%%%%%%%%%%%%%%%%%%%%%%%%%%%%
%%%%%%%%%%%%%%%%%%%%%%%%%%%%%%%%%%%%%%%%%%%%%%%%%%%%%%%%%%
%%%%%%%%%%%%%%%%%%%%%%%%%%%%%%%%%%%%%%%%%%%%%%%%%%%%%%%%%%

\subsection{Efficiency analysis: the fully discrete scheme}\label{efficiency_fully}

In this section, we derive local efficiency estimates for the local a posteriori error indicators $\mathcal{E}_{st,K}$ and $\mathcal{E}_{adj,K}$ and a global efficiency estimate for the a posteriori error estimator $\mathcal{E}_{ct,\T}$. To accomplish this task, we will proceed on the basis of standard bubble function arguments. 

We start our analysis by introducing the following notation: for an edge, triangle, or tetrahedron $G$, let $\mathcal{V}(G)$ be the set of vertices of $G$. Using this notation, we introduce the standard element and edge bubble functions: For $K\in\mathscr{T}$ and $\gamma\in\mathscr{S}$, 
\begin{equation}\label{def:standard_bubbles}
\varphi^{}_{K}=
(d+1)^{d+1}\prod_{\textsc{v} \in \mathcal{V}(K)} \lambda^{}_{\textsc{v}},
\qquad
\varphi^{}_{\gamma}=
d^d \prod_{\textsc{v} \in \mathcal{V}(\gamma)}\lambda^{}_{\textsc{v}}|^{}_{K'},
\quad 
K' \subset \mathcal{N}_{\gamma}.
\end{equation} 
Here, $\lambda_{\textsc{v}}$ denotes the barycentric coordinate function associated with the vertex $\textsc{v}$, and $\mathcal{N}_{\gamma}$ corresponds to the patch composed of the two elements of $\mathscr{T}$ that share $\gamma$.

The following identities are essential for performing an efficiency analysis. Since the pair $(\bar{\mathbf{y}},\bar{\mathsf{\mathsf{p}}})\in\mathbf{H}_0^1(\Omega)\times L_0^2(\Omega)$ solves \eqref{eq:brinkamn_problem_state}, where $u$ is replaced by $\bar{u}$, an elementwise integration by parts formula yields
\begin{multline}\label{eq:brinkamn_problem01_aux_error}
(\nabla \mathbf{e}_{\bar{\mathbf{y}}},\nabla \mathbf{v})_{\mathbf{L}^2(\Omega)}+( e_{\bar{u}}\bar{\mathbf{y}}, \mathbf{v})_{\mathbf{L}^2(\Omega)}-(e_{\bar{\mathsf{p}}},\text{div }\mathbf{v})_{L^2(\Omega)} + (\mathsf{q},\text{div }\mathbf{e}_{\bar{\mathbf{y}}})_{L^2(\Omega)}\\= -(\bar{u}_{h}\mathbf{e}_{\bar{\mathbf{y}}},\mathbf{v})_{\mathbf{L}^2(\Omega)}+ \sum_{\gamma\in \mathscr{S}} (\llbracket(\nabla \bar{\mathbf{y}}_{h}-\bar{\mathsf{p}}_{h}\mathbf{I})\cdot \mathbf{n}\rrbracket,\mathbf{v})_{\mathbf{L}^2(\gamma)}-\sum_{K\in\T}(\mathsf{q},\text{div }\bar{\mathbf{y}}_{h})_{\mathbf{L}^2(K)}\\
+
\sum_{K\in\T}
\left[(\mathsf{P}_{K}\mathbf{f}+\Delta \bar{\mathbf{y}}_{h}-\bar{u}_{h}\bar{\mathbf{y}}_{h}-\nabla\bar{\mathsf{p}}_{h},\mathbf{v})_{\mathbf{L}^2(K)}+(\mathbf{f}-\mathsf{P}_{K}\mathbf{f},\mathbf{v})_{\mathbf{L}^2(K)}\right],
\end{multline}
for all $\mathbf{v}\in \mathbf{H}_0^1(\Omega)$ and $\mathsf{q}\in L_0^2(\Omega)$. Here, $\mathsf{P}_K$ denotes the $L^2$-projection operator onto piecewise constant functions over $K\in\T$. Similarly, since $(\bar{\mathbf{z}},\bar{\mathsf{\mathsf{r}}})\in\mathbf{H}_0^1(\Omega)\times L_0^2(\Omega)$ solves the adjoint problem \eqref{eq:brinkamn_problem_adjoint} with $\mathbf{y}=\bar{\mathbf{y}}$ and $u=\bar{u}$, we have
\begin{multline}\label{eq:brinkamn_problem02_aux_error}
(\nabla \mathbf{e}_{\bar{\mathbf{z}}},\nabla \mathbf{v})_{\mathbf{L}^2(\Omega)}+( e_{\bar{u}}\bar{\mathbf{z}}, \mathbf{v})_{\mathbf{L}^2(\Omega)}+(e_{\bar{\mathsf{r}}},\text{div }\mathbf{v})_{L^2(\Omega)} + (\mathsf{q},\text{div }\mathbf{e}_{\bar{\mathbf{z}}})_{L^2(\Omega)}\\= -(\bar{u}_{h}\mathbf{e}_{\bar{\mathbf{z}}},\mathbf{v})_{\mathbf{L}^2(\Omega)}+ (\mathbf{e}_{\bar{\mathbf{y}}},\mathbf{v})+ \sum_{\gamma\in \mathscr{S}} (\llbracket(\nabla \bar{\mathbf{z}}_{h}+\bar{\mathsf{r}}_{h}\mathbf{I})\cdot \mathbf{n}\rrbracket,\mathbf{v})_{\mathbf{L}^2(\gamma)}-\sum_{K\in\T}(\mathsf{q},\text{div }\bar{\mathbf{z}}_{h})_{\mathbf{L}^2(K)}
\\
+
\sum_{K\in\T}
\left[
(\bar{\mathbf{y}}_{h}-\mathsf{P}_{K}\mathbf{y}_{\Omega}+\Delta \bar{\mathbf{z}}_{h}-\bar{u}_{h}\bar{\mathbf{z}}_{h}+\nabla\bar{\mathsf{r}}_{h},\mathbf{v})_{\mathbf{L}^2(K)}+(\mathsf{P}_{K}\mathbf{y}_{\Omega}-\mathbf{y}_{\Omega},\mathbf{v})_{\mathbf{L}^2(K)}
\right]
\end{multline}
for all $\mathbf{v}\in \mathbf{H}_0^1(\Omega)$ and $\mathsf{q}\in L_0^2(\Omega)$.
% We observe that $\mathsf{P}_K$ denotes the $L^2$-projection operator onto piecewise constant functions over $K\in\T$. 
As a final element, we introduce, for $\mathbf{v}\in \mathbf{L}^2(\Omega)$ and $\mathcal{M}\subset\T$, the oscillation term
\begin{align*}
\text{osc}_{\mathscr{T}}(\mathbf{v},\mathcal{M}):= \left( \sum_{K\in\mathcal{M}}h_{K}^2 \|\mathbf{v}-\mathsf{P}_K \mathbf{v}\|_{\mathbf{L}^2(K)}^2 \right)^{\frac{1}{2}}.
\end{align*}

With these ingredients in hand, we now proceed to obtain local efficiency properties for the indicator $\mathcal{E}_{st,K}$ defined in \S\ref{Real:fully_discrete}.
% defined in \eqref{eq:estimator_st}.

\begin{theorem}[local efficiency of $\mathcal{E}_{st,K}$]\label{eq:efficiency_est}
Let $\bar{u}$ be a local solution of \eqref{eq:min_functional}--\eqref{eq:brinkamn_problem_state} such that \eqref{equivalence_condition} holds for $\tau > 0$, and let $(\bar{\mathbf{y}},\bar{\mathsf{p}})$ and $(\bar{\mathbf{z}},\bar{\mathsf{r}})$ be the corresponding state and ajoint state variables. Let $\bar{u}_h$ be a local minimum of the fully discrete scheme, and let $(\bar{\mathbf{y}}_h,\bar{\mathsf{p}}_h)$ and $(\bar{\mathbf{z}}_h,\bar{\mathsf{r}}_h)$ be the corresponding discrete state and adjoint state variables, respectively. Then, for $K\in\T$, the local error indicator $\mathcal{E}_{st,K}$ satisfies the bound
\begin{multline}\label{eq:efficiency_estimates_est}
\mathcal{E}_{st,K} \lesssim\|\nabla \mathbf{e}_{\bar{\mathbf{y}}} \|_{\mathbf{L}^2(\mathcal{N}_{K})}+\| e_{\bar{\mathsf{p}}} \|_{L^2(\mathcal{N}_{K})}
+ h_K\|e_{\bar{u}}\|_{L^2(\mathcal{N}_{K})}\\
+h_{K}\|\mathbf{e}_{\bar{\mathbf{y}}}\|_{\mathbf{L}^2(\mathcal{N}_{K})}+\mathrm{osc}_{\mathscr{T}}(\mathbf{f},\mathcal{N}_{K}),
\end{multline}
where $\mathcal{N}_{K}$ is defined in \eqref{eq:patch}. The hidden constant is independent of the continuous and discrete optimal variables, the size of the elements in the mesh $\T$, and $\#\T$.
\end{theorem}
\begin{proof}
We proceed in four steps.

\emph{Step 1.} Let $K\in\T$. Define $\tilde{\mathcal{R}}_{K}^{st}:=(\mathsf{P}_{K}\mathbf{f}+\Delta \bar{\mathbf{y}}_{h}-\bar{u}_{h}\bar{\mathbf{y}}_{h}-\nabla\bar{\mathsf{p}}_{h})|_K$. A first bound for the term $h_K^2\|\mathcal{R}_{K}^{st}\|_{\mathbf{L}^2(K)}^2$ follows from the triangle inequality:
\begin{align}\label{eq:bound01_eff}
h_K^2\|\mathcal{R}_{K}^{st}\|_{\mathbf{L}^2(K)}^2\lesssim h_K^2 \|\tilde{\mathcal{R}}_{K}^{st}\|_{\mathbf{L}^2(K)}^2 + \text{osc}_{\mathscr{T}}(\mathbf{f},K)^2.
\end{align}
To bound $h_K^2 \|\tilde{\mathcal{R}}_{K}^{st}\|_{\mathbf{L}^2(K)}^2$, we set $\mathbf{v}=\varphi_K \tilde{\mathcal{R}}_{K}^{st}$ and $\mathsf{q}=0$ in \eqref{eq:brinkamn_problem01_aux_error}, where $\varphi_K $ denotes the bubble function defined in \eqref{def:standard_bubbles}, and use standard properties of $\varphi_K$ and inverse inequalities. These arguments show that
\begin{multline}\label{eq:bound02_eff}
\|\tilde{\mathcal{R}}_{K}^{st}\|_{\mathbf{L}^2(K)} \lesssim h_{K}^{-1} [\|\nabla \mathbf{e}_{\bar{\mathbf{y}}}\|_{\mathbf{L}^2(K)}
+
\| e_{\bar{\mathsf{p}}} \|_{L^2(K)}]
\\
+
\|\bar{\mathbf{y}}\|_{\mathbf{L}^{\infty}(K)}\|e_{\bar{u}}\|_{L^2(K)}
+
\|\bar{u}_{h}\|_{L^{\infty}(K)}\| \mathbf{e}_{\bar{\mathbf{y}}}\|_{\mathbf{L}^2(K)}+\|\mathbf{f}-\mathsf{P}_{K}\mathbf{f}\|_{\mathbf{L}^2(K)}.
\end{multline}
We note that, as a consequence of Theorem \ref{theorem_regularity_Linfty}, we have $\|\bar{\mathbf{y}}\|_{\mathbf{L}^{\infty}(K)}\lesssim \|\mathbf{f}\|_{\mathbf{L}^2(\Omega)}$. The fact that $\bar{u}_h\in \mathbb{U}_{ad}$ guarantees that $\bar{u}_h \in L^{\infty}(\Omega)$. Consequently,
\begin{multline}\label{eq:bound03_eff}
h_K\|\tilde{\mathcal{R}}_{K}^{st}\|_{\mathbf{L}^2(K)}  
\lesssim  
\|\nabla \mathbf{e}_{\bar{\mathbf{y}}}\|_{\mathbf{L}^2(K)}
+\| e_{\bar{\mathsf{p}}} \|_{L^2(K)}
\\
+
h_K\| \mathbf{e}_{\bar{\mathbf{y}}}\|_{\mathbf{L}^2(K)}
+
h_K \|e_{\bar{u}}\|_{L^2(K)}
+
\text{osc}_{\mathscr{T}}(\mathbf{f},K).
\end{multline}
If we replace this bound into \eqref{eq:bound01_eff}, we obtain the desired bound for the term $h_K^2 \|\mathcal{R}_{K}^{st}\|_{\mathbf{L}^2(K)}^2$.

\emph{Step 2.} Let $K\in\T$ and $\gamma\in \mathscr{S}$. We now control $h_K^{\frac{1}{2}}\|\mathcal{J}_{\gamma}^{st}\|_{\mathbf{L}^2(\gamma)}$. To do so, we set $\mathbf{v}=\varphi_{\gamma}\mathcal{J}_{\gamma}^{st}$ and $\mathsf{q}=0$ in \eqref{eq:brinkamn_problem01_aux_error} and continue with similar arguments as we developed for \eqref{eq:bound03_eff}:
\begin{multline*}
\|\mathcal{J}_{\gamma}^{st}\|_{\mathbf{L}^2(\gamma)}^2\lesssim \sum_{K'\in \mathcal{N}_{\gamma}} 
\Big(h_{K'}^{-1}
\left(\|\nabla \mathbf{e}_{\bar{\mathbf{y}}}\|_{\mathbf{L}^2(K')}
+
\| e_{\bar{\mathsf{p}}} \|_{L^2(K')}\right)
+
\|\bar{\mathbf{y}}\|_{\mathbf{L}^{\infty}(K')} \|e_{\bar{u}}\|_{L^2(K')}
\\
+
\|\bar{u}_{h}\|_{L^{\infty}(K')}\| \mathbf{e}_{\bar{\mathbf{y}}}\|_{\mathbf{L}^2(K')}+\|\mathbf{f}-\mathsf{P}_{K'}\mathbf{f}\|_{\mathbf{L}^2(K')}+\|\tilde{\mathcal{R}}_{K'}^{st}\|_{\mathbf{L}^2(K')}\Big)h_{K'}^{\frac{1}{2}}\|\mathcal{J}_{\gamma}^{st}\|_{\mathbf{L}^2(\gamma)}.
\end{multline*}
Using \eqref{eq:bound03_eff}, the bound $\|\bar{\mathbf{y}}\|_{\mathbf{L}^{\infty}(K')}\lesssim \|\mathbf{f}\|_{\mathbf{L}^2(\Omega)}$, and the fact that $\bar{u}_{h}\in \mathbb{U}_{ad}$, we obtain
\begin{multline}\label{eq:bound04_eff}
	h_{K}^{\frac{1}{2}}\|\mathcal{J}_{\gamma}^{st}\|_{\mathbf{L}^2(\gamma)}\lesssim \sum_{K'\in \mathcal{N}_{\gamma}} 
	\Big(\|\nabla \mathbf{e}_{\bar{\mathbf{y}}}\|_{\mathbf{L}^2(K')}
	+\| e_{\bar{\mathsf{p}}} \|_{L^2(K')}
\\
+h_{K'} \|e_{\bar{u}}\|_{L^2(K')} + h_{K'}\| \mathbf{e}_{\bar{\mathbf{y}}}\|_{\mathbf{L}^2(K')}+\text{osc}_{\mathscr{T}}(\mathbf{f},K')\Big).
\end{multline}

\emph{Step 3.} Let $K\in \T$. The control of the term $\|\text{div }\bar{\mathbf{y}}_{h}\|_{L^2(K)}^2$ follows easily from the incompressibility condition $\text{div }\bar{\mathbf{y}}=0$. In fact,
\begin{align}\label{eq:bound05_eff}
	\|\text{div }\bar{\mathbf{y}}_{h}\|_{L^2(K)}=\|\text{div }\mathbf{e}_{\bar{\mathbf{y}}}\|_{L^2(K)}\lesssim \|\nabla\mathbf{e}_{\bar{\mathbf{y}}}\|_{\mathbf{L}^2(K)}.
\end{align}

\emph{Step 4.} The proof concludes by collecting the estimates \eqref{eq:bound03_eff}, \eqref{eq:bound04_eff}, and \eqref{eq:bound05_eff}.
\end{proof}

In the following, we study local efficiency properties for the indicator $\mathcal{E}_{adj,K}$ defined in \S\ref{Real:fully_discrete}.

\begin{theorem}[local efficiency of $\mathcal{E}_{adj,K}$]\label{eq:efficiency_adj}
In the framework of Theorem \ref{eq:efficiency_est}, we have that, for $K\in\T$, the local error indicator $\mathcal{E}_{adj,K}$ satisfies the bound 
\begin{multline}\label{eq:efficiency_estimates_adj}
\mathcal{E}_{adj,K} \lesssim\|\nabla \mathbf{e}_{\bar{\mathbf{z}}} \|_{\mathbf{L}^2(\mathcal{N}_{K})}+\| e_{\bar{\mathsf{r}}} \|_{L^2(\mathcal{N}_{K})}
+
h_K \|e_{\bar{u}}\|_{L^2(\mathcal{N}_{K})}\\+h_{K}\|\mathbf{e}_{\bar{\mathbf{z}}}\|_{\mathbf{L}^2(\mathcal{N}_{K})}+h_{K}\|\mathbf{e}_{\bar{\mathbf{y}}}\|_{\mathbf{L}^2(\mathcal{N}_{K})}+\mathrm{osc}_{\mathscr{T}}(\mathbf{y}_{\Omega},\mathcal{N}_{K}),
\end{multline}
where $\mathcal{N}_{K}$ is defined in \eqref{eq:patch}. The hidden constant is independent of the continuous and discrete optimal variables, the size of the elements in the mesh $\T$, and $\#\T$.
\end{theorem}
\begin{proof}
We proceed in four steps.

\emph{Step 1:} Let $K\in\T$. Define $\tilde{\mathcal{R}}_{K}^{adj}:=(\bar{\mathbf{y}}_{h}-\mathsf{P}_{K}\mathbf{y}_{\Omega}+\Delta \bar{\mathbf{z}}_{h}\!-\bar{u}_{h}\bar{\mathbf{z}}_{h}\!+\!\nabla\bar{\mathsf{r}}_{h})|_K$. To bound $h_K^2\|\mathcal{R}_{K}^{adj}\|_{\mathbf{L}^2(K)}^2$, we first apply the triangle inequality and obtain the basic bound
\begin{align}\label{eq:bound06_eff}
h_K^2\|\mathcal{R}_{K}^{adj}\|_{\mathbf{L}^2(K)}^2\lesssim h_K^2 \|\tilde{\mathcal{R}}_{K}^{adj}\|_{\mathbf{L}^2(K)}^2 + \text{osc}_{\mathscr{T}}(\mathbf{y}_{\Omega},K)^2.
\end{align}
To bound $h_K^2 \|\tilde{\mathcal{R}}_{K}^{adj}\|_{\mathbf{L}^2(K)}^2$, we set $\mathbf{v}=\varphi_K \tilde{\mathcal{R}}_{K}^{st}$ and $\mathsf{q}=0$ in \eqref{eq:brinkamn_problem02_aux_error} and apply standard properties of the bubble function $\varphi_K$ and basic inequalities to arrive at 
\begin{multline}\label{eq:bound07_eff}
\|\tilde{\mathcal{R}}_{K}^{adj}\|_{\mathbf{L}^2(K)} \lesssim h_{K}^{-1} \left[
\|\nabla \mathbf{e}_{\bar{\mathbf{z}}}\|_{\mathbf{L}^2(K)}
+
\| e_{\bar{\mathsf{r}}} \|_{L^2(K)}
\right]
+
\|\bar{\mathbf{z}}\|_{\mathbf{L}^{\infty}(K)} \|e_{\bar{u}}\|_{L^2(K)}
\\
+
\|\bar{u}_{h}\|_{L^{\infty}(K)}\| \mathbf{e}_{\bar{\mathbf{z}}}\|_{\mathbf{L}^2(K)}+\| \mathbf{e}_{\bar{\mathbf{y}}}\|_{\mathbf{L}^2(K)}+\|\mathbf{y}_{\Omega}-\mathsf{P}_{K}\mathbf{y}_{\Omega}\|_{\mathbf{L}^2(K)}.
\end{multline}
We note that, as a consequence of Theorem \ref{theorem_regularity_Linfty}, we have the bound $\|\bar{\mathbf{z}}\|_{\mathbf{L}^{\infty}(K)}\lesssim \|\mathbf{f}\|_{\mathbf{L}^2(\Omega)}+\|\mathbf{y}_{\Omega}\|_{\mathbf{L}^2(\Omega)}$. The fact that $\bar{u}_h \in \mathbb{U}_{ad}$ guarantees that $\bar{u}_h \in L^{\infty}\Omega)$. Thus,
\begin{multline}\label{eq:bound08_eff}
h_K\|\tilde{\mathcal{R}}_{K}^{adj}\|_{\mathbf{L}^2(K)}  \lesssim  \|\nabla \mathbf{e}_{\bar{\mathbf{z}}}\|_{\mathbf{L}^2(K)}
+\| e_{\bar{\mathsf{r}}} \|_{L^2(K)}
+h_K \|e_{\bar{u}}\|_{L^2(K)}
\\+h_K\| \mathbf{e}_{\bar{\mathbf{z}}}\|_{\mathbf{L}^2(K)}+h_K\| \mathbf{e}_{\bar{\mathbf{y}}}\|_{\mathbf{L}^2(K)}+\text{osc}_{\mathscr{T}}(\mathbf{y}_{\Omega},K).
\end{multline}
If we replace this bound into \eqref{eq:bound06_eff}, we obtain the desired bound for $h_K^2 \|\mathcal{R}_{K}^{adj}\|_{\mathbf{L}^2(K)}^2$.
%Replacing this bound into the estimate \eqref{eq:bound06_eff}, we conclude the desired control for the residual term $h_K^2 \|\mathcal{R}_{K}^{adj}\|_{\mathbf{L}^2(K)}^2$.

\emph{Step 2.} Let $K\in\T$ and $\gamma\in \mathscr{S}$. We now control $h_K^{\frac{1}{2}}\|\mathcal{J}_{\gamma}^{adj}\|_{\mathbf{L}^2(\gamma)}$. To do so, we set $\mathbf{v}=\varphi_{\gamma}\mathcal{J}_{\gamma}^{adj}$ and $\mathsf{q}=0$ in \eqref{eq:brinkamn_problem02_aux_error} and continue with similar arguments as we developed for \eqref{eq:bound08_eff}. In fact, we have the bound
\begin{multline}\label{eq:bound09_eff}
h_{K}^{\frac{1}{2}}\|\mathcal{J}_{\gamma}^{adj}\|_{\mathbf{L}^2(\gamma)}\lesssim \sum_{K'\in \mathcal{N}_{\gamma}} \Big(\|\nabla \mathbf{e}_{\bar{\mathbf{z}}}\|_{\mathbf{L}^2(K')} +\| e_{\bar{\mathsf{r}}} \|_{L^2(K')} +
h_{K'} \|e_{\bar{u}}\|_{L^2(K')}
\\
+
h_{K'}\| \mathbf{e}_{\bar{\mathbf{y}}}\|_{\mathbf{L}^2(K')}
+
h_{K'}\| \mathbf{e}_{\bar{\mathbf{z}}}\|_{\mathbf{L}^2(K')}+\text{osc}_{\mathscr{T}}(\mathbf{y}_{\Omega},K')\Big).
\end{multline}

\emph{Step 3:} Let $K\in \T$. Since $\text{div }\bar{\mathbf{z}}=0$, we immediately obtain that
%The control of the term $\|\text{div }\bar{\mathbf{z}}_{h}\|_{L^2(K)}^2$ in \eqref{eq:estimator_AD} follows from the incompressibility condition $\text{div }\bar{\mathbf{z}}=0$. In fact,
\begin{align}\label{eq:bound10_eff}
	\|\text{div }\bar{\mathbf{z}}_{h}\|_{L^2(K)}=\|\text{div }\mathbf{e}_{\bar{\mathbf{z}}}\|_{L^2(K)}\lesssim \|\nabla\mathbf{e}_{\bar{\mathbf{z}}}\|_{\mathbf{L}^2(K)}.
\end{align}

\emph{Step 4:} The proof concludes by gathering the bounds \eqref{eq:bound08_eff}, \eqref{eq:bound09_eff}, and \eqref{eq:bound10_eff}.
%The desired estimate is obtained by collecting the bounds \eqref{eq:bound08_eff}, \eqref{eq:bound09_eff}, and \eqref{eq:bound10_eff}. This completes the proof.
\end{proof}

Finally, we present a global efficiency result for the estimator $\mathcal{E}_{ocp,\T}$.

\begin{theorem}[global efficiency of $\mathcal{E}_{ocp,\T}$]\label{eq:efficiency_ocp}
In the framework of Theorem \ref{eq:efficiency_est}, we have that, for $K\in\T$, the estimator $\mathcal{E}_{ocp,\T}$ satisfies the global bound
\begin{multline}\label{eq:efficiency_estimates_ocp}
\mathcal{E}_{ocp,\T} \lesssim\|\nabla \mathbf{e}_{\bar{\mathbf{y}}} \|_{\mathbf{L}^2(\Omega)}+\|\nabla \mathbf{e}_{\bar{\mathbf{z}}} \|_{\mathbf{L}^2(\Omega)}+\| e_{\bar{\mathsf{p}}} \|_{L^2(\Omega)}+\| e_{\bar{\mathsf{r}}} \|_{L^2(\Omega)}\\+\|e_{\bar{u}}\|_{L^2(\Omega)}
+
\mathrm{osc}_{\mathscr{T}}(\mathbf{f},\T)
+
\mathrm{osc}_{\mathscr{T}}(\mathbf{y}_{\Omega},\T).
\end{multline}
The hidden constant is independent of the continuous and discrete optimal variables, the size of the elements in the mesh $\T$, and $\#\T$.
\end{theorem}
\begin{proof}
We begin the proof by considering the definition of the error estimator $\mathcal{E}_{st,\T}$ given in \eqref{eq:estimator_st} and the local efficiency estimate \eqref{eq:efficiency_estimates_est} to conclude that
\begin{multline}\label{eq:efficiency_estimates_est_global}
\mathcal{E}_{st,\T} \lesssim\|\nabla \mathbf{e}_{\bar{\mathbf{y}}} \|_{\mathbf{L}^2(\Omega)}+\| e_{\bar{\mathsf{p}}} \|_{L^2(\Omega)}
+
\text{diam}(\Omega) \|e_{\bar{u}}\|_{L^2(\Omega)}
\\
+\text{diam}(\Omega)\|\mathbf{e}_{\bar{\mathbf{y}}}\|_{\mathbf{L}^2(\mathcal{N}_{K})}+\text{osc}_{\mathscr{T}}(\mathbf{f},\T).
\end{multline}
Similarly, considering the definition of the error estimator $\mathcal{E}_{adj,\T}$ in \eqref{eq:estimator_AD} and the local efficiency estimate \eqref{eq:efficiency_estimates_adj}, we can come to the following conclusion
\begin{multline}\label{eq:efficiency_estimates_adj_global}
\mathcal{E}_{adj,\T} \lesssim\|\nabla \mathbf{e}_{\bar{\mathbf{z}}} \|_{\mathbf{L}^2(\Omega)}
+
\| e_{\bar{\mathsf{r}}} \|_{L^2(\Omega)}
+
\text{diam}(\Omega) \|e_{\bar{u}}\|_{L^2(\Omega)}\\
+
\text{diam}(\Omega)\|\mathbf{e}_{\bar{\mathbf{z}}}\|_{\mathbf{L}^2(\Omega)}+\text{diam}(\Omega)\|\mathbf{e}_{\bar{\mathbf{y}}}\|_{\mathbf{L}^2(\Omega)}+\text{osc}_{\mathscr{T}}(\mathbf{y}_{\Omega},\T).
\end{multline}
Our aim now is to control the term $\mathcal{E}_{ct,\T}$ with the above estimates. To this end, we invoke \eqref{eq:estimator_control} and use a simple application of a triangle inequality to obtain
\begin{eqnarray*}
\mathcal{E}_{ct,\T}	&\leq &\|\breve{u}-\bar{u}\|_{L^2(\Omega)}+\|e_{\bar{u}}\|_{L^2(\Omega)}
\\
&\leq &\|\Pi_{[\mathsf{a},\mathsf{b}]}(\alpha^{-1}\bar{\mathbf{y}}_{h}\cdot\bar{\mathbf{z}}_{h})-\Pi_{[\mathsf{a},\mathsf{b}]}(\alpha^{-1}\bar{\mathbf{y}}\cdot\bar{\mathbf{z}})\|_{L^2(\Omega)}+\|e_{\bar{u}}\|_{L^2(\Omega)}.
\end{eqnarray*}
The application of the Lipschitz property of the projection $\Pi_{[\mathsf{a},\mathsf{b}]}$, the Cauchy--Schwarz inequality, and the standard Sobolev embedding $\mathbf{H}_0^1(\Omega)\hookrightarrow\mathbf{L}^4(\Omega)$ in the previous estimate show that
\begin{align*}
\mathcal{E}_{ct,\T}\leq \alpha^{-1} (\|\nabla \mathbf{e}_{\bar{\mathbf{z}}}\|_{\mathbf{L}^2(\Omega)}\|\nabla \bar{\mathbf{y}}_{h}\|_{\mathbf{L}^2(\Omega)}+\|\nabla \mathbf{e}_{\bar{\mathbf{y}}}\|_{\mathbf{L}^2(\Omega)}\|\nabla \bar{\mathbf{z}}\|_{\mathbf{L}^2(\Omega)})+\|e_{\bar{u}}\|_{L^2(\Omega)}.
\end{align*}
The bounds $\|\nabla \bar{\mathbf{y}}_{h}\|_{\mathbf{L}^2(\Omega)}\lesssim \|\mathbf{f}\|_{\mathbf{L}^2(\Omega)}$ and $\|\nabla \bar{\mathbf{z}}\|_{\mathbf{L}^2(\Omega)}\lesssim  \|\mathbf{f}\|_{\mathbf{L}^2(\Omega)}+ \|\mathbf{y}_{\Omega}\|_{\mathbf{L}^2(\Omega)}$ therefore result in the bound
\begin{align}\label{eq:efficiency_estimates_ct_global}
\mathcal{E}_{ct,\T}\lesssim \|\nabla \mathbf{e}_{\bar{\mathbf{z}}}\|_{\mathbf{L}^2(\Omega)}+\|\nabla \mathbf{e}_{\bar{\mathbf{y}}}\|_{\mathbf{L}^2(\Omega)}+\|e_{\bar{u}}\|_{L^2(\Omega)}.
\end{align}
The desired estimate \eqref{eq:efficiency_estimates_ocp} follows from the application of Poincare's inequality and a collection of the bounds \eqref{eq:efficiency_estimates_est_global}, \eqref{eq:efficiency_estimates_adj_global}, and \eqref{eq:efficiency_estimates_ct_global}. This concludes the proof.
\end{proof}

\subsection{Efficiency analysis: the semidiscrete scheme}\label{efficiency_semi}
We begin the section by defining the local error indicator $\mathfrak{E}_{ocp,K}^2:=\mathfrak{E}_{st,K}^2+\mathfrak{E}_{adj,K}^2$ for $K\in\T$.

We now present an analysis of local efficiency obtained by adapting the arguments elaborated in \S\ref{efficiency_fully}. For the sake of brevity, we omit the details of the proof.

\begin{theorem}[local efficiency of $\mathfrak{E}_{ocp,K}$]\label{eq:efficiency_est_semi}
Let $\bar{u}$ be a local solution of \eqref{eq:min_functional}--\eqref{eq:brinkamn_problem_state} such that \eqref{equivalence_condition} holds for $\tau > 0$, and let $(\bar{\mathbf{y}},\bar{\mathsf{p}})$ and $(\bar{\mathbf{z}},\bar{\mathsf{r}})$ be the corresponding state and ajoint state variables. Let $\bar{\mathsf{u}}$ be a local minimum of the semidiscrete scheme, and let $(\bar{\mathbf{y}}_h,\bar{\mathsf{p}}_h)$ and $(\bar{\mathbf{z}}_h,\bar{\mathsf{r}}_h)$ be the corresponding discrete state and adjoint state variables, respectively. Then, for $K\in\T$, the local error indicator $\mathfrak{E}_{ocp,K}$ satisfies the bound 
\begin{multline}\label{eq:efficiency_estimates_ocp_semi}
\mathfrak{E}_{ocp,K} \lesssim\|\nabla \mathbf{e}_{\bar{\mathbf{y}}} \|_{\mathbf{L}^2(\mathcal{N}_{K})}+\|\nabla \mathbf{e}_{\bar{\mathbf{z}}} \|_{\mathbf{L}^2(\mathcal{N}_{K})}
+
\| e_{\bar{\mathsf{p}}} \|_{L^2(\mathcal{N}_{K})}
+
\| e_{\bar{\mathsf{r}}} \|_{L^2(\mathcal{N}_{K})}
\\
+
h_K \| \mathbf{e}_{\bar{\mathbf{y}}} \|_{\mathbf{L}^2(\mathcal{N}_{K})}
+
h_K\| \mathbf{e}_{\bar{\mathbf{z}}} \|_{\mathbf{L}^2(\mathcal{N}_{K})}
+
h_K \|e_{\bar{u}}\|_{L^2(\mathcal{N}_{K})}
+
\mathrm{osc}_{\mathscr{T}}(\mathbf{f},\mathcal{N}_{K})\\
+
\mathrm{osc}_{\mathscr{T}}(\mathbf{y}_{\Omega},\mathcal{N}_{K})
+
\mathrm{osc}_{\mathscr{T}}(\bar{\mathsf{u}}\bar{\mathbf{y}}_h,\mathcal{N}_{K})
+
\mathrm{osc}_{\mathscr{T}}(\bar{\mathsf{u}}\bar{\mathbf{z}}_h,\mathcal{N}_{K}),
\end{multline}
where $\mathcal{N}_{K}$ is defined in \eqref{eq:patch}. The hidden constant is independent of the continuous and discrete optimal variables, the size of the elements in the mesh $\T$, and $\#\T$.
\end{theorem}

%%%%%%%%%%%%%%%%%%%%%%%%%%%%%%%%%%%%%%%%%%%%%%%%%%%%%%%%%%
%%%%%%%%%%%%%%%%%%%%%%%%%%%%%%%%%%%%%%%%%%%%%%%%%%%%%%%%%%
%%%%%%%%%%%%%%%%%%%%%%%%%%%%%%%%%%%%%%%%%%%%%%%%%%%%%%%%%%
%%%%%%%%%%%%%%%%%%%%%%%%%%%%%%%%%%%%%%%%%%%%%%%%%%%%%%%%%%

\section{Numerical experiments}
\label{sec:numericalexperiments}

In this section, we present a series of numerical experiments that illustrate the performance of the fully discrete scheme of \S \ref{fully_discrete_framework}. The experiments were performed with a code implemented in \texttt{C++}: the global linear systems were solved using the multifrontal massively parallel sparse direct solver (MUMPS) \cite{MUMPS2} and a quadrature formula, which guarantees accuracy up to polynomials of degree 19, was used to compute all integrals.

The framework for the discrete scheme is as follows: For a given partition $\T_h$, we seek $(\bar{\mathbf{y}}_h,\bar{\mathsf{p}}_h,\bar{\mathbf{z}}_h,\bar{\mathsf{r}}_h,\bar{u}_h)\in \mathbf{X}_h \times M_h\times \mathbf{X}_h \times M_h\times \mathbb{U}_{ad,h}$, which solves the discrete problem \eqref{eq:brinkamn_problem_state_fd}--\eqref{eq:brinkman_problem_adjoint_fully}.
% , which will be referred as the fully discrete solution. 
This problem was solved with an adaptation of the semi-smooth Newton method described in \cite[Appendix A.1]{MR2971171}. After the Newton iteration is solved, an adaptive strategy is performed. An element $K\in\T_{h}$ is refined using the following maximum strategy: Refine $K\in\T_h$ if the error indicator $\mathrm{E}_{K}$ is such that $\mathrm{E}_{K} \geq 0.5 \max_{K'\in\T_h}\mathrm{E}_{K'}$. Here, $\mathrm{E}_{K}^2:=\mathcal{E}_{st,K}^2 + \mathcal{E}_{adj,K}^2 + \mathcal{E}_{ct,K}^2$ (see \S  \ref{Real:fully_discrete} for details). 

To present our results, we also introduce the \emph{total number of degrees of freedom} as $\mathsf{Ndof}=2\text{dim} (\mathbf{X}_h)+2\text{dim}(M_h)+\dim(\mathbb{U}_{ad,h})$. The initial meshes we used for the resolution are shown in Fig. \ref{fig:meshes}.

\begin{figure}[!h]
\centering
\begin{minipage}[b]{0.4\textwidth}
\centering
\includegraphics[trim={0 0 0 0},clip,width=1.5cm,height=1.5cm,scale=0.66]{fig/mesh_square.eps} 
\end{minipage}
\begin{minipage}[b]{0.4\textwidth}
\centering
\includegraphics[trim={0 0 0 0},clip,width=2cm,height=1.5cm,scale=0.66]{fig/mesh_L.eps} 
\end{minipage}
\caption{Initial meshes: $\Omega=(0,1)^{2}$ (left) and $\Omega=(-1, 1) \times (-1, 1) \setminus [0, 1] \times [-1,0]$ (right).}
\label{fig:meshes}
\end{figure}

\subsection{A sharp interior layer: uniform vs. adaptive refinement}

In this section, we provide a numerical example for which an exact solution is known: we fix the state and the associated adjoint velocities and pressures and obtain an exact optimal control. With these ingredients, we can compute $\mathbf{f}$ and the desired state $\mathbf{y}_{\Omega}$.

\textbf{Example 1.} We set $\Omega=(0,1)^{2}$, $\mathsf{a}=0.1$, $\mathsf{b}=0.2$, and $\alpha=1.0$. The state and adjoint variables are given by: $\bar{\mathbf{y}}=50( \partial_{y}\xi, -\partial_{x}\xi )$, $\bar{\mathsf{p}}=xy(x - 1)(y - 1) - 1/36$, $\bar{\mathbf{z}}=50( \partial_{y}\varsigma , -\partial_{x}\varsigma )$, and $\bar{\mathsf{r}} =\sin(2\pi x)\sin(2\pi y)$, where $\xi=(xy(1-x)(1-y))^2$ and $\varsigma=(xy(1-x)(1-y))^2\tan^{-1}((x-0.5)/0.01)$. Due to the structure of the problem, we emphasize that the adjoint velocity field has a sharp interior layer at $x=0.5$.

In Fig. \ref{fig:test_01_1}, we show the velocity, pressure, and control errors for uniform and adaptive refinement. We note that the adjoint velocity error does not achieve an optimal convergence rate (see panel (A.1)) and that the error estimator only achieves optimal convergence when adaptive refinement is performed (see panel (A.4)). In Fig. \ref{fig:test_01_2}, we present a mesh that contains 8124 triangles and discrete velocities and control solutions: the sharp layer of the adjoint velocity is properly recovered by the adaptive procedure (see panel (B.3)).   

\begin{figure}[!h]
\centering
\psfrag{Seminorma-y-uniforme----}{$\|\nabla \mathbf{e}_{\bar{\mathbf{y}}} \|_{\mathbf{L}^2(\Omega)}$ (unif)}
\psfrag{Seminorma-z-uniforme----}{$\|\nabla \mathbf{e}_{\bar{\mathbf{z}}} \|_{\mathbf{L}^2(\Omega)}$ (unif)}
\psfrag{Seminorma-y-adaptiva----}{$\|\nabla \mathbf{e}_{\bar{\mathbf{y}}} \|_{\mathbf{L}^2(\Omega)}$ (adap)}
\psfrag{Seminorma-z-adaptiva----}{$\|\nabla \mathbf{e}_{\bar{\mathbf{z}}} \|_{\mathbf{L}^2(\Omega)}$ (adap)}
\psfrag{norma-p-uniforme----}{$\| e_{\bar{\mathsf{p}}} \|_{L^2(\Omega)}$ (unif)}
\psfrag{norma-r-uniforme----}{$\| e_{\bar{\mathsf{r}}} \|_{L^2(\Omega)}$ (unif)}
\psfrag{norma-p-adaptiva----}{$\| e_{\bar{\mathsf{p}}} \|_{L^2(\Omega)}$ (adap)}
\psfrag{norma-r-adaptiva----}{$\| e_{\bar{\mathsf{r}}} \|_{L^2(\Omega)}$ (adap)}
\psfrag{norma-u-uniforme----}{$\| e_{\bar{\mathsf{u}}} \|_{L^2(\Omega)}$ (unif)}
\psfrag{norma-u-adaptiva----}{$\| e_{\bar{\mathsf{u}}} \|_{L^2(\Omega)}$ (adap)}
\psfrag{estimador-uniforme----}{$\mathcal{E}_{ocp,\T}$ (unif)}
\psfrag{estimador-adaptiva----}{$\mathcal{E}_{ocp,\T}$ (adap)}
\psfrag{Ndof(-1)}{\small $\mathsf{Ndof}^{-1}$}
\psfrag{Ndof(-0.5)}{\small $\mathsf{Ndof}^{-1/2}$}
\begin{minipage}[b]{0.25\textwidth}
\centering
\includegraphics[trim={0 0 0 0},clip,width=3cm,height=3cm,scale=0.66]{fig/convergecia_velocidad_exam_1.eps} \\
\tiny{(A.1)}
\end{minipage}
\begin{minipage}[b]{0.25\textwidth}
\centering
\includegraphics[trim={0 0 0 0},clip,width=3cm,height=3cm,scale=0.66]{fig/convergecia_presion_exam_1.eps} \\
\tiny{(A.2)}
\end{minipage}
\begin{minipage}[b]{0.25\textwidth}
\centering
\includegraphics[trim={0 0 0 0},clip,width=3cm,height=3cm,scale=0.66]{fig/convergecia_control_exam_1.eps} \\
\tiny{(A.3)}
\end{minipage}
\begin{minipage}[b]{0.25\textwidth}
\centering
\includegraphics[trim={0 0 0 0},clip,width=3cm,height=3cm,scale=0.66]{fig/convergecia_estimador_exam_1.eps} \\
\tiny{(A.4)}
\end{minipage}
\caption{Example 1: Experimental convergence rates for the fully discrete scheme with uniform and adaptive refinements for the state and adjoint velocity errors (A.1); the state and adjoint pressure errors (A.2), the control errors (A.3), and the error estimators (A.4).}
\label{fig:test_01_1}
\end{figure}

\begin{figure}[!h]
\centering
\begin{minipage}[b]{0.25\textwidth}
\centering
\includegraphics[trim={0 0 0 0},clip,width=3cm,height=3cm,scale=0.66]{fig/mesh_1.eps} \\
\tiny{(B.1)}
\end{minipage}
\begin{minipage}[b]{0.25\textwidth}
\centering
\includegraphics[trim={0 0 0 0},clip,width=3cm,height=3cm,scale=0.66]{fig/state_vel_1.eps} \\
\tiny{(B.2)}
\end{minipage}
\begin{minipage}[b]{0.25\textwidth}
\centering
\includegraphics[trim={0 0 0 0},clip,width=3cm,height=3cm,scale=0.66]{fig/adjoint_vel_1.eps} \\
\tiny{(B.3)}
\end{minipage}
\begin{minipage}[b]{0.25\textwidth}
\centering
\includegraphics[trim={0 0 0 0},clip,width=3cm,height=3cm,scale=0.66]{fig/control_1.eps} \\
\tiny{(B.4)}
\end{minipage}
\caption{Example 1: For the fully discrete scheme, we present the mesh obtained after 30 adaptive refinements (8.124 triangles and 4.083 vertices) (B.1), the discrete optimal state velocity $|\bar{\mathbf{y}}_{h}|$ (B.2), the discrete optimal adjoint velocity $|\bar{\mathbf{z}}_{h}|$ (B.3), and the discrete optimal control $\bar{u}_{h}$ (B.4).}
\label{fig:test_01_2}
\end{figure}

\subsection{A non-convex domain} 

 In this section, we perform the adaptive strategy over an L-shaped domain for which no analytical solution is known.
 
\textbf{Example 2.} We set $\Omega=(-1, 1) \times (-1, 1) \setminus [0,1] \times [-1,0]$, $\mathsf{a}=1.0$, $\mathsf{b}=5.0$, $\alpha=1.0$, $\mathbf{f}=1000((x+y)^4,\sin(2\pi x)\sin(2\pi y))^{\mathrm{T}}$, and $\mathbf{y}_{\Omega}=1000(\sin(2\pi x)\sin(2\pi y),xy(1-x)(1-y))^{\mathrm{T}}$.

In Fig. \ref{fig:test_02}, we show the experimental convergence of the error estimator $\mathcal{E}_{ocp,\T}$ and the finite element approximations of the optimal state velocity and pressure, as well as the control variable. We can see that optimal convergence is achieved for the error estimator (see panel (C.1)) and that most of the refinement is concentrated in the reentrant corner of the domain (see panel (C.5)). 

\begin{figure}[!h]
\centering
\psfrag{estimador------}{$\mathcal{E}_{ocp,\T}$}
\psfrag{Ndof(-1)}{\small $\mathsf{Ndof}^{-1}$}
\begin{minipage}[b]{0.3\textwidth}
\centering
\includegraphics[trim={0 0 0 0},clip,width=3cm,height=3cm,scale=0.66]{fig/estimator-L.eps} \\
\tiny{(C.1)}
\end{minipage}
\begin{minipage}[b]{0.3\textwidth}
\centering
\includegraphics[trim={0 0 0 0},clip,width=4cm,height=3cm,scale=0.66]{fig/velocity_state_L_3.eps} \\
\tiny{(C.2)}
\end{minipage}
\begin{minipage}[b]{0.3\textwidth}
\centering
\includegraphics[trim={0 0 0 0},clip,width=4cm,height=3cm,scale=0.66]{fig/pressure_state_L_3.eps} \\
\tiny{(C.3)}
\end{minipage}
\\
\begin{minipage}[b]{0.4\textwidth}
\centering
\includegraphics[trim={0 0 0 0},clip,width=4cm,height=3cm,scale=0.66]{fig/control_L_3.eps} \\
\tiny{(C.4)}
\end{minipage}
\begin{minipage}[b]{0.4\textwidth}
\centering
\includegraphics[trim={0 0 0 0},clip,width=4cm,height=3cm,scale=0.66]{fig/mesh_L_final.eps} \\
\tiny{(C.5)}
\end{minipage}
\caption{Example 2: For the fully discrete scheme, we present the experimental convergence rate for the error estimator $\mathcal{E}_{ocp,\T}$ (C.1) and the discrete optimal state velocity $|\bar{\mathbf{y}}_{h}|$ (C.2), the discrete optimal state pressure $\bar{\mathbf{p}}_{h}$ (C.3), and the discrete optimal control $\bar{\mathbf{u}}_{h}$ (C.4) over a mesh with 8.326 elements and 24.978 vertices obtained after 40 adaptive iterations (C.5).}
\label{fig:test_02}
\end{figure}

%\section{Numerical experiments}
%\label{sec:numericalexperiments}
%
%In this section, we present a series of numerical experiments that illustrate the performance of the fully discrete method presented in \S \ref{fully_discrete_framework}. To accomplish this task, we implemented a \texttt{C++} code to carry out the numerical examples. In doing so, the global linear systems were solved using the multifrontal massively parallel sparse direct solver (MUMPS) \cite{MUMPS2}. A quadrature formula was used to compute all integrals, which guarantees accuracy up to polynomials of degree 19. 
%
%We now present the setting for the discrete scheme. For a given partition $\T_h$, we seek $(\bar{\mathbf{y}}_h,\bar{\mathsf{p}}_h,\bar{\mathbf{z}}_h,\bar{\mathsf{r}}_h,\bar{u}_h)\in \mathbf{X}_h \times M_h\times \mathbf{X}_h \times M_h\times \mathbb{U}_{ad,h}$ that solves the discrete problem \eqref{eq:brinkamn_problem_state_fd}--\eqref{eq:brinkman_problem_adjoint_fully}, which will be referred as the fully discrete solution. This problem was solved using an adaptation of the semi-smooth Newton method described in \cite[Appendix A.1]{MR2971171}. After the Newton iteration is resolved, an adaptive strategy is performed, where elements $K\in\T_{h}$ are refined by using the following maximum strategy: refine element $K\in\T_h$ if the error indicators satisfy $\mathrm{E}_{K} \geq  \frac{1}{2}\max_{K'\in\T_h}\mathrm{E}_{K'}$, where the error indicator is given by $\mathrm{E}_{K}:=\mathcal{E}_{st,K}^2 + \mathcal{E}_{adj,K}^2 + \mathcal{E}_{ct,K}^2$ (see Section \ref{Real:fully_discrete}).
%
%To the present the numerical experiments, the total number of degrees of freedom is $\mathsf{Ndof}=2\text{dim} (\mathbf{X}_h)+2\text{dim}(M_h)+\dim(\mathbb{U}_{ad,h})$ and the initial meshes for the problems are shown in Fig. \ref{fig:meshes}.
%
%\begin{figure}[!h]
%\centering
%\begin{minipage}[b]{0.4\textwidth}
%\centering
%\includegraphics[trim={0 0 0 0},clip,width=1.5cm,height=1.5cm,scale=0.66]{fig/mesh_square.eps} 
%\end{minipage}
%\begin{minipage}[b]{0.4\textwidth}
%\centering
%\includegraphics[trim={0 0 0 0},clip,width=2cm,height=1.5cm,scale=0.66]{fig/mesh_L.eps} 
%\end{minipage}
%\caption{Initial meshes: $\Omega=(0,1)^{2}$ (left) and $\Omega=(-1, 1) \times (-1, 1) \setminus [0, 1] \times [-1,0]$ (right).}
%\label{fig:meshes}
%\end{figure}
%
%\subsection{A sharp interior layer: uniform vs adaptive refinement on a convex domain}
%
%In this section, we provide a numerical example where exact solutions are known. In doing so, we fix the velocity and pressure state and adjoint variables to obtain the exact optimal control. With this at hand, we compute the source term $\mathbf{f}$ and the desired state $\mathbf{y}_{\Omega}$.
%~\\~\\
%\textbf{Example 1.} We set $\Omega=(0,1)^{2}$, $\mathsf{a}=0.1$, $\mathsf{b}=0.2$ and $\alpha=1.0$. The state and adjoint variables are given by: $\bar{\mathbf{y}}=50( \partial_{y}\xi, -\partial_{x}\xi )$, $\bar{\mathsf{p}}=xy(x - 1)(y - 1) - 1/36$, $\bar{\mathbf{z}}=50( \partial_{y}\varsigma , -\partial_{x}\varsigma )$, and $\bar{\mathsf{r}} =\sin(2\pi x)\sin(2\pi y)$, where $\xi=(xy(1-x)(1-y))^2$ and $\varsigma=(xy(1-x)(1-y))^2\tan^{-1}((x-0.5)/0.01)$. Due to the structure of the problem, we highlight that the adjoint velocity field presents a sharp interior layer at $x=0.5$.
%
%In Fig. \ref{fig:test_01_1}, we show velocity, pressure, and control errors for uniform and adaptive refinement. We can observe that the adjoint velocity error does not attain optimal convergence (see panel (A.1)), and the error estimator only attains optimal convergence when adaptive refinement is performed (see panel (A.4)). In Fig. \ref{fig:test_01_2}, we present a final mesh, velocities and final control discrete solutions, where the sharp layer of the adjoint velocity is properly recovered by the adaptive procedure (see panel (B.3)).   
%
%\begin{figure}[!h]
%\centering
%\psfrag{Seminorma-y-uniforme----}{$\|\nabla \mathbf{e}_{\bar{\mathbf{y}}} \|_{\mathbf{L}^2(\Omega)}$ (unif)}
%\psfrag{Seminorma-z-uniforme----}{$\|\nabla \mathbf{e}_{\bar{\mathbf{z}}} \|_{\mathbf{L}^2(\Omega)}$ (unif)}
%\psfrag{Seminorma-y-adaptiva----}{$\|\nabla \mathbf{e}_{\bar{\mathbf{y}}} \|_{\mathbf{L}^2(\Omega)}$ (adap)}
%\psfrag{Seminorma-z-adaptiva----}{$\|\nabla \mathbf{e}_{\bar{\mathbf{z}}} \|_{\mathbf{L}^2(\Omega)}$ (adap)}
%\psfrag{norma-p-uniforme----}{$\| e_{\bar{\mathsf{p}}} \|_{L^2(\Omega)}$ (unif)}
%\psfrag{norma-r-uniforme----}{$\| e_{\bar{\mathsf{r}}} \|_{L^2(\Omega)}$ (unif)}
%\psfrag{norma-p-adaptiva----}{$\| e_{\bar{\mathsf{p}}} \|_{L^2(\Omega)}$ (adap)}
%\psfrag{norma-r-adaptiva----}{$\| e_{\bar{\mathsf{r}}} \|_{L^2(\Omega)}$ (adap)}
%\psfrag{norma-u-uniforme----}{$\| e_{\bar{\mathsf{u}}} \|_{L^2(\Omega)}$ (unif)}
%\psfrag{norma-u-adaptiva----}{$\| e_{\bar{\mathsf{u}}} \|_{L^2(\Omega)}$ (adap)}
%\psfrag{estimador-uniforme----}{$\mathcal{E}_{ocp,\T}$ (unif)}
%\psfrag{estimador-adaptiva----}{$\mathcal{E}_{ocp,\T}$ (adap)}
%\psfrag{Ndof(-1)}{\small $\mathsf{Ndof}^{-1}$}
%\psfrag{Ndof(-0.5)}{\small $\mathsf{Ndof}^{-1/2}$}
%\begin{minipage}[b]{0.25\textwidth}
%\centering
%\includegraphics[trim={0 0 0 0},clip,width=3cm,height=3cm,scale=0.66]{fig/convergecia_velocidad_exam_1.eps} \\
%\tiny{(A.1)}
%\end{minipage}
%\begin{minipage}[b]{0.25\textwidth}
%\centering
%\includegraphics[trim={0 0 0 0},clip,width=3cm,height=3cm,scale=0.66]{fig/convergecia_presion_exam_1.eps} \\
%\tiny{(A.2)}
%\end{minipage}
%\begin{minipage}[b]{0.25\textwidth}
%\centering
%\includegraphics[trim={0 0 0 0},clip,width=3cm,height=3cm,scale=0.66]{fig/convergecia_control_exam_1.eps} \\
%\tiny{(A.3)}
%\end{minipage}
%\begin{minipage}[b]{0.25\textwidth}
%\centering
%\includegraphics[trim={0 0 0 0},clip,width=3cm,height=3cm,scale=0.66]{fig/convergecia_estimador_exam_1.eps} \\
%\tiny{(A.4)}
%\end{minipage}
%\caption{Example 1: Experimental rates of convergence for the fully discrete scheme with uniform and adaptive refinements for: state and adjoint velocity errors (A.1); state and adjoint pressure errors (A.2); control errors (A.3); error estimators (A.4).}
%\label{fig:test_01_1}
%\end{figure}
%
%\begin{figure}[!h]
%\centering
%\begin{minipage}[b]{0.25\textwidth}
%\centering
%\includegraphics[trim={0 0 0 0},clip,width=3cm,height=3cm,scale=0.66]{fig/mesh_1.eps} \\
%\tiny{(B.1)}
%\end{minipage}
%\begin{minipage}[b]{0.25\textwidth}
%\centering
%\includegraphics[trim={0 0 0 0},clip,width=3cm,height=3cm,scale=0.66]{fig/state_vel_1.eps} \\
%\tiny{(B.2)}
%\end{minipage}
%\begin{minipage}[b]{0.25\textwidth}
%\centering
%\includegraphics[trim={0 0 0 0},clip,width=3cm,height=3cm,scale=0.66]{fig/adjoint_vel_1.eps} \\
%\tiny{(B.3)}
%\end{minipage}
%\begin{minipage}[b]{0.25\textwidth}
%\centering
%\includegraphics[trim={0 0 0 0},clip,width=3cm,height=3cm,scale=0.66]{fig/control_1.eps} \\
%\tiny{(B.4)}
%\end{minipage}
%\caption{Example 1: for the fully discrete scheme, we present the final mesh after 30 adaptive refinement that contains 8.124 triangles and 4.083 vertices (B.1). The discrete optimal state velocity $|\bar{\mathbf{y}}_{h}|$ (B.2). The discrete optimal adjoint velocity $|\bar{\mathbf{z}}_{h}|$ (B.3). The discrete optimal control $\bar{u}_{h}$ (B.4).}
%\label{fig:test_01_2}
%\end{figure}
%
%\subsection{Unknown solution over a nonconvex L-shape domain} 
%
%In this section, we conduct the adaptive strategy over a L-shape domain where no analytical solution is known. 
%~\\~\\
%\textbf{Example 2.} We let $\Omega=(-1, 1) \times (-1, 1) \setminus [0,1] \times [-1,0]$, $\mathsf{a}=1.0$, $\mathsf{b}=5.0$, $\alpha=1.0$, $\mathbf{f}=1000((x+y)^4,\sin(2\pi x)\sin(2\pi y))^{\mathrm{T}}$, and $\mathbf{y}_{\Omega}=1000(\sin(2\pi x)\sin(2\pi y),xy(1-x)(1-y))^{\mathrm{T}}$.
%
%In Fig. \ref{fig:test_02}, we present the experimental convergence of the error estimator $\mathcal{E}_{ocp,\T}$ and the finite element solution for the optimal state velocity and pressure, and the discrete optimal control. We can see that optimal convergence for the error estimator is attained (see panel (C.1)), and most of the refinement is concentrated in the reentrant corner of the domain (see panel (C.5)). 
%
%\begin{figure}[!h]
%\centering
%\psfrag{estimador------}{$\mathcal{E}_{ocp,\T}$}
%\psfrag{Ndof(-1)}{\small $\mathsf{Ndof}^{-1}$}
%\begin{minipage}[b]{0.3\textwidth}
%\centering
%\includegraphics[trim={0 0 0 0},clip,width=3cm,height=3cm,scale=0.66]{fig/estimator-L.eps} \\
%\tiny{(C.1)}
%\end{minipage}
%\begin{minipage}[b]{0.3\textwidth}
%\centering
%\includegraphics[trim={0 0 0 0},clip,width=4cm,height=3cm,scale=0.66]{fig/velocity_state_L_3.eps} \\
%\tiny{(C.2)}
%\end{minipage}
%\begin{minipage}[b]{0.3\textwidth}
%\centering
%\includegraphics[trim={0 0 0 0},clip,width=4cm,height=3cm,scale=0.66]{fig/pressure_state_L_3.eps} \\
%\tiny{(C.3)}
%\end{minipage}
%\\
%\begin{minipage}[b]{0.4\textwidth}
%\centering
%\includegraphics[trim={0 0 0 0},clip,width=4cm,height=3cm,scale=0.66]{fig/control_L_3.eps} \\
%\tiny{(C.4)}
%\end{minipage}
%\begin{minipage}[b]{0.4\textwidth}
%\centering
%\includegraphics[trim={0 0 0 0},clip,width=4cm,height=3cm,scale=0.66]{fig/mesh_L_final.eps} \\
%\tiny{(C.5)}
%\end{minipage}
%\caption{Example 2: for the fully discrete scheme, we present the experimental convergence of the error estimator $\mathcal{E}_{ocp,\T}$ (C.1), the discrete optimal state velocity $|\bar{\mathbf{y}}_{h}|$ (C.2), the discrete optimal state pressure $\bar{\mathbf{p}}_{h}$ (C.3), the discrete optimal control $\bar{\mathbf{u}}_{h}$ (C.4), over a mesh that contains 8.326 elements and 24.978 vertices which was obtained after 40 adaptive iterations (C.5).}
%\label{fig:test_02}
%\end{figure}


%\section{Numerical experiments}
%\label{sec:numericalexperiments}
%
%In this section, we present a numerical experiment that illustrates the performance of the fully and semidiscrete methods presented in Sections \ref{fully_discrete_framework} and \ref{semi_discrete_framework}. To accomplish this task, we implemented a \texttt{C++} code to carry out the numerical examples presented. In doing so, all matrices were assembled exactly, and the global linear systems were solved using the multifrontal massively parallel sparse direct solver (MUMPS) \cite{MUMPS1,MUMPS2}. A quadrature formula was used to compute all integrals, guaranteeing accuracy by using polynomials of degree 19. %The ParaView application \cite{Ahrens2005ParaViewAE,Ayachit2015ThePG} was used to visualize finite element approximations.
%
%We now present the setting to the fully discrete scheme. For a given partition $\T_h$, we seek $(\bar{\mathbf{y}}_h,\bar{\mathsf{p}}_h,\bar{\mathbf{z}}_h,\bar{\mathsf{r}}_h,\bar{u}_h)\in \mathbf{X}_h \times M_h\times \mathbf{X}_h \times M_h\times \mathbb{U}_{ad,h}$ that solves the discrete problem \eqref{eq:brinkamn_problem_state_fd}--\eqref{eq:brinkman_problem_adjoint_fully}. To solve this optimality system, we use \textbf{Algoritms 1} and \textbf{2}. We observe that the \textbf{Algorithm 1} describes a variant of the \emph{primal-dual active set strategy}; see, for instance, \cite[section 2.12.4]{MR2583281}. In addition, \textbf{Algorithm 2} presents the also well--known \emph{Newton method}; see \cite[section 4.4.1]{MR1817388}. To introduce the \emph{Newton method}, we define $\mathcal{X}_h= \mathbf{X}_h \times M_h\times \mathbf{X}_h \times M_h\times \mathbb{U}_{ad,h}$ and present, for $\Phi=(\mathbf{y}_h,\mathsf{p}_h,\mathbf{z}_h,\mathsf{r}_h,u_h)$ and $\Theta=(\mathbf{v}_h,\mathsf{q}_h,\mathbf{v}_h,\mathsf{q}_h,t_h)$, the operator $\mathcal{F}_h: \mathcal{X}_h\to \mathcal{X}_h'$ as follows
%\begin{multline}\label{eq:newton_function}
%\langle F_h(\Phi),\Theta\rangle:=[(\nabla \mathbf{y}_{h},\nabla \mathbf{v}_{h})_{\mathbf{L}^2(\Omega)}+( u_{h}\mathbf{y}_{h}, \mathbf{v}_{h})_{\mathbf{L}^2(\Omega)}-(\mathsf{p}_{h},\text{div }\mathbf{v}_{h})_{L^2(\Omega)} 
%\\ -(\mathbf{f},\mathbf{v}_{h})_{\mathbf{L}^2(\Omega)}
%,(\mathsf{q}_{h},\text{div }\mathbf{y}_{h})_{L^2(\Omega)},
%(\nabla \mathbf{z}_{h},\nabla \mathbf{v}_{h})_{\mathbf{L}^2(\Omega)}+(u_{h}\mathbf{z}_{h}, \mathbf{v}_{h})_{\mathbf{L}^2(\Omega)}
%\\+(\mathsf{r}_{h},\text{div }\mathbf{v}_{h})_{L^2(\Omega)}
% - (\mathbf{y}_{h}-\mathbf{y}_{\Omega},\mathbf{v}_{h})_{\mathbf{L}^2(\Omega)},(\mathsf{s}_{h},\text{div }\bar{\mathbf{z}}_{h})_{L^2(\Omega)} , 
% \\(u_h\textbf{1}-\mathsf{a}\boldsymbol{\mathcal{X}}_{\mathsf{a}}-\mathsf{b}\boldsymbol{\mathcal{X}}_{\mathsf{b}}- \alpha^{-1}\Pi_{h}(\mathbf{y}_h\cdot\mathbf{z}_h)(\textbf{1}-\boldsymbol{\mathcal{X}}_{\mathsf{a}}-\boldsymbol{\mathcal{X}}_{\mathsf{b}}),t_h)_{L^2(\Omega)}
%]^{\mathsf{T}}
%\end{multline}
%Here, $\Pi_h$ represents $L^2$-projection operator onto piecewise constant functions over $\T_h$, and $\langle\cdot,\cdot\rangle$ denotes the duality pairing between $\mathcal{X}_h'$ and $\mathcal{X}_h$. Moreover, we consider $\boldsymbol{\mathcal{X}}_\mathsf{a},\boldsymbol{\mathcal{X}}_\mathsf{b}\in\mathbb{R}^{\#\T_h}$, $\textbf{1}=(1,\ldots,1)\in\mathbb{R}^{\#\T_h}$. We define for $k\in\mathbb{N}_0$, $\Phi_k:=(\mathbf{u}_h^k,\mathsf{p}_h^k,\mathbf{z}_h^k,\mathsf{r}_h^k,u_h^k)$. We now set an initial point $\Phi_0$, and we consider the following Newton iteration: $\Phi_{k+1}=\Phi_k + \eta$, where $\eta:= (\delta\mathbf{y}_h,\delta\mathsf{p}_h,\delta\mathbf{z}_h,\delta\mathsf{r}_h,\delta u_h) \in \mathcal{X}_h$ solves the iterative problem
%\begin{align}\label{eq:solve_incremental}
%\langle F_h'(\Phi_k)\eta,\Theta\rangle = -\langle F_h(\Phi_k),\Theta\rangle\quad \forall \Theta \in  \mathcal{X}_h.
%\end{align}
%Where, $F_h'(\Phi_k)\eta$ represents the G\^ateaux derivate of the operator $F_h$ in $\Phi_k$ evaluated at the direction $\eta$.
%
%%%%%
%\begin{algorithm}[ht]
%\caption{\textbf{Active set algorithm}}
%\label{Algorithm1}
%\textbf{Input:} Mesh $\mathscr{T}_h$, constraints $\mathsf{a}$ and $\mathsf{b}$, regularization parameter $\alpha$, and initial guess $\Phi_0=(\mathbf{y}_h^0,\mathsf{p}_h^0,\mathbf{z}_h^0,\mathsf{r}_h^0,u_h^0)\in\mathbf{X}_h \times M_h\times \mathbf{X}_h \times M_h\times \mathbb{U}_{ad,h}$;
%\\
%\textbf{Step} $\boldsymbol{1}$: Define $\boldsymbol{\mathcal{X}}_{\mathsf{a}}^{old}=(\mathcal{X}_{\mathsf{a},K}^{old})_{K\in\T}$, $\boldsymbol{\mathcal{X}}_{\mathsf{b}}^{old}=(\mathcal{X}_{\mathsf{b},K}^{old})_{K\in\T}\in\mathbb{R}^{\# \mathscr{T}_h}$ where $\mathcal{X}_{\mathsf{a},K}^{old},\mathcal{X}_{\mathsf{b},K}^{old}\in\{0,1\}$.
%
%\emph{Set:} $k=0.$
%\\
%\textbf{Step} $\boldsymbol{2}$: Applying \textbf{Algorithm 2} compute $[\mathbf{u}_h^{k+1},\mathsf{p}_h^{k+1},\mathbf{z}_h^{k+1},\mathsf{r}_h^{k+1},u_h^{k+1}]=\textbf{Newton}[\T,\mathsf{a},\mathsf{b},\alpha,\boldsymbol{\mathcal{X}}_{\mathsf{a}}^{old},\boldsymbol{\mathcal{X}}_{\mathsf{b}}^{old},\mathbf{u}_h^k,\mathsf{p}_h^k,\mathbf{z}_h^k,\mathsf{r}_h^k,u_h^k]$.
%\\
%\textbf{Step} $\boldsymbol{3}$: For each $K\in\mathscr{T}_h$ determine
%$$\mathcal{X}_{\mathsf{a},K}^{new}=\left\{ 
%\begin{array}{cl}
%1 & \text{if}\quad  \Pi_{K}(\alpha^{-1}\mathbf{y}_h^{k+1}\cdot\mathbf{z}_h^{k+1})<\mathsf{a},\\ 
%0 & \text{otherwise}.
%\end{array}%
%\right.$$ $$\mathcal{X}_{\mathsf{b},K}^{new}=\left\{ 
%\begin{array}{cl}
%1 & \text{if}\quad  \Pi_{K}(\alpha^{-1}\mathbf{y}_h^{k+1}\cdot\mathbf{z}_h^{k+1})>\mathsf{b},\\ 
%0 & \text{otherwise}.
%\end{array}%
%\right.$$
%where $\Pi_{K}$ denotes the $L^2$--projection onto piecewise constant functions over $K.$
%\\
%\textbf{Step} $\boldsymbol{4}$: If $\sum_{K\in\T}(|\mathcal{X}_{\mathsf{a},K}^{new}-\mathcal{X}_{\mathsf{a},K}^{old}|+|\mathcal{X}_{\mathsf{b},K}^{new}-\mathcal{X}_{\mathsf{b},K}^{old}|)=0$, set $(\bar{\mathbf{y}}_h,\bar{\mathsf{p}}_h,\bar{\mathbf{z}}_h,\bar{\mathsf{r}}_h,\bar{u}_h)=(\mathbf{u}_h^{k+1},\mathsf{p}_h^{k+1},\mathbf{z}_h^{k+1},\mathsf{r}_h^{k+1},u_h^{k+1})$. Otherwise, set $\boldsymbol{\mathcal{X}}_{\mathsf{a}}^{old}:=\boldsymbol{\mathcal{X}}_{\mathsf{a}}^{new}$, $\boldsymbol{\mathcal{X}}_{\mathsf{b}}^{old}:=\boldsymbol{\mathcal{X}}_{\mathsf{b}}^{new}$, and $k\leftarrow k+1$ and go to \textbf{Step 2}.
%\end{algorithm}
%%%%%
%
%%%TABLA 2
%%%%
%\begin{algorithm}[ht]
%\caption{\textbf{Newton method}}
%\label{Algorithm2}
%\textbf{Input:} Mesh $\mathscr{T}_h$, constraints $\mathsf{a}$ and $\mathsf{b}$, regularization parameter $\alpha$, initial guess $(\mathbf{y}_h^0,\mathsf{p}_h^0,\mathbf{z}_h^0,\mathsf{r}_h^0,u_h^0)\in\mathbf{X}_h \times M_h\times \mathbf{X}_h \times M_h\times \mathbb{U}_{ad,h}$, and $\boldsymbol{\mathcal{X}}_{\mathsf{a}},\boldsymbol{\mathcal{X}}_{\mathsf{b}}\in\mathbb{R}^{\# \mathscr{T}_h}$;
%
%\emph{Set:} $k=0.$
%\\
%\textbf{Step} $\boldsymbol{1}$: Given $(\mathbf{u}_h^k,\mathsf{p}_h^k,\mathbf{z}_h^k,\mathsf{r}_h^k,u_h^k)\in \mathbf{X}_h \times M_h\times \mathbf{X}_h \times M_h\times \mathbb{U}_{ad,h}$, compute the incremental $\eta=(\delta\mathbf{y}_h,\delta\mathsf{p}_h,\delta\mathbf{z}_h,\delta\mathsf{r}_h,\delta u_h) \in \mathbf{X}_h \times M_h\times \mathbf{X}_h \times M_h\times \mathbb{U}_{ad,h}$ as the solution to \eqref{eq:solve_incremental}.
%\\
%\textbf{Step} $\boldsymbol{2}$: Compute the variables $(\mathbf{u}_h^{k+1},\mathsf{p}_h^{k+1},\mathbf{z}_h^{k+1},\mathsf{r}_h^{k+1},u_h^{k+1}):=(\mathbf{u}_h^k,\mathsf{p}_h^k,\mathbf{z}_h^k,\mathsf{r}_h^k,u_h^k)+(\delta\mathbf{y}_h,\delta\mathsf{p}_h,\delta\mathbf{z}_h,\delta\mathsf{r}_h,\delta u_h)$.
%\\
%\textbf{Step} $\boldsymbol{3}$: If $\max\{\|\delta\mathbf{y}_h\|_{\mathbf{L}^\infty(\Omega)},\|\delta\mathsf{p}_h\|_{L^\infty(\Omega)},\|\delta\mathbf{z}_h\|_{\mathbf{L}^\infty(\Omega)},\|\delta\mathsf{r}_h\|_{L^\infty(\Omega)},\|\delta u_h\|_{L^\infty(\Omega)}\}<10^{-8}$, set $(\bar{\mathbf{y}}_h,\bar{\mathsf{p}}_h,\bar{\mathbf{z}}_h,\bar{\mathsf{r}}_h,\bar{u}_h)=(\mathbf{u}_h^{k+1},\mathsf{p}_h^{k+1},\mathbf{z}_h^{k+1},\mathsf{r}_h^{k+1},u_h^{k+1})$. Otherwise $k\leftarrow k+1$ and go to \textbf{Step 1}.
%\end{algorithm}
%%%%
%
%The total number of degrees of freedom is $\mathsf{Ndof}=2\text{dim} (\mathbf{X}_h)+2\text{dim}(M_h)+\dim(\mathbb{U}_{ad,h})$. We also present the framework of the semidiscrete problem. We observe that ....
%
%We present the setting of the numerical experiment for the bilinear optimal control of the Stokes--Brinkman equations: we set $d=2$, $\Omega=(0,1)\times (0,1)$, $\mathsf{a}= ... $, $\mathsf{b}= ... $, and $\alpha=...$. We observe that the desired state $\mathbf{y}_{\Omega}$ and the forcing term $\mathbf{f}$ are such that the exact optimal state and the optimal adjoint state are
%\begin{align}
%\bar{\mathbf{y}}=... \qquad \qquad \bar{\mathbf{z}}=...
%\end{align}
%To conclude, we consider that the optimal control $\bar{u}$ can be obtained in view of a projection formula: $\bar{u}(x):=\Pi_{[\mathsf{a},\mathsf{b}]}(\alpha^{-1}\bar{\mathbf{y}}\cdot\bar{\mathbf{z}})$ a.e. $x \in \Omega$,
%
%\textbf{Example 1:} ...
%
%\textbf{Example 2:} ... 

\bibliographystyle{siamplain}
\bibliography{biblio}
\end{document}


\begin{thebibliography}{10}

\bibitem{MR2424078}
{\sc R.~A. Adams and J.~J.~F. Fournier}, {\em Sobolev spaces}, vol.~140 of Pure
  and Applied Mathematics (Amsterdam), Elsevier/Academic Press, Amsterdam,
  second~ed., 2003.

\bibitem{MR1445736}
{\sc M.~Ainsworth and J.~T. Oden}, {\em A posteriori error estimators for the
  {S}tokes and {O}seen equations}, SIAM J. Numer. Anal., 34 (1997),
  pp.~228--245, \url{http://dx.doi.org/10.1137/S0036142994264092}.

\bibitem{MR3097958}
{\sc D.~Boffi, F.~Brezzi, and M.~Fortin}, {\em Mixed finite element methods and
  applications}, vol.~44 of Springer Series in Computational Mathematics,
  Springer, Heidelberg, 2013,
  \url{http://dx.doi.org/10.1007/978-3-642-36519-5}.

\bibitem{MR2373954}
{\sc S.~C. Brenner and L.~R. Scott}, {\em The mathematical theory of finite
  element methods}, vol.~15 of Texts in Applied Mathematics, Springer, New
  York, third~ed., 2008, \url{http://dx.doi.org/10.1007/978-0-387-75934-0}.

\bibitem{MR4679130}
{\sc E.~Casas, K.~Chrysafinos, and M.~Mateos}, {\em Semismooth {N}ewton method
  for boundary bilinear control}, IEEE Control Syst. Lett., 7 (2023),
  pp.~3549--3554, \url{http://dx.doi.org/10.1109/lcsys.2023.3337747}.

\bibitem{casas2024error}
{\sc E.~Casas, K.~Chrysafinos, and M.~Mateos}, {\em Error estimates for the
  discretization of bilinear control problems governed by semilinear elliptic
  pdes}, 2024, \href{http://arxiv.org/abs/2404.05658}{arXiv:2404.05658}.

\bibitem{MR3586845}
{\sc E.~Casas and M.~Mateos}, {\em Optimal control of partial differential
  equations}, in Computational mathematics, numerical analysis and
  applications, vol.~13 of SEMA SIMAI Springer Ser., Springer, Cham, 2017,
  pp.~3--59.

\bibitem{MR2272157}
{\sc E.~Casas and J.-P. Raymond}, {\em Error estimates for the numerical
  approximation of {D}irichlet boundary control for semilinear elliptic
  equations}, SIAM J. Control Optim., 45 (2006), pp.~1586--1611,
  \url{http://dx.doi.org/10.1137/050626600}.

\bibitem{MR2902693}
{\sc E.~Casas and F.~Tr\"{o}ltzsch}, {\em Second order analysis for optimal
  control problems: improving results expected from abstract theory}, SIAM J.
  Optim., 22 (2012), pp.~261--279, \url{http://dx.doi.org/10.1137/110840406}.

\bibitem{MR3103238}
{\sc Y.~Chen, Z.~Lu, and Y.~Huang}, {\em Superconvergence of triangular
  {R}aviart-{T}homas mixed finite element methods for a bilinear constrained
  optimal control problem}, Comput. Math. Appl., 66 (2013), pp.~1498--1513,
  \url{http://dx.doi.org/10.1016/j.camwa.2013.08.019}.

\bibitem{CiarletBook}
{\sc P.~G. Ciarlet}, {\em The finite element method for elliptic problems},
  SIAM, Philadelphia, PA, 2002,
  \url{http://dx.doi.org/10.1137/1.9780898719208}.

\bibitem{MR0977489}
{\sc M.~Dauge}, {\em Stationary {S}tokes and {N}avier-{S}tokes systems on two-
  or three-dimensional domains with corners. {I}. {L}inearized equations}, SIAM
  J. Math. Anal., 20 (1989), pp.~74--97,
  \url{http://dx.doi.org/10.1137/0520006}.

\bibitem{Guermond-Ern}
{\sc A.~Ern and J.-L. Guermond}, {\em Theory and practice of finite elements},
  vol.~159 of Applied Mathematical Sciences, Springer-Verlag, New York, 2004,
  \url{http://dx.doi.org/10.1007/978-1-4757-4355-5}.

\bibitem{MR3693332}
{\sc H.~Fu, H.~Guo, J.~Hou, and J.~Zhang}, {\em A stabilized mixed finite
  element approximation of bilinear state optimal control problems}, Comput.
  Math. Appl., 74 (2017), pp.~1246--1261,
  \url{http://dx.doi.org/10.1016/j.camwa.2017.06.010}.

\bibitem{MR4450052}
{\sc F.~Fuica and E.~Ot\'{a}rola}, {\em A posteriori error estimates for an
  optimal control problem with a bilinear state equation}, J. Optim. Theory
  Appl., 194 (2022), pp.~543--569,
  \url{http://dx.doi.org/10.1007/s10957-022-02039-6}.

\bibitem{MR3422453}
{\sc V.~Girault, R.~H. Nochetto, and L.~R. Scott}, {\em Max-norm estimates for
  {S}tokes and {N}avier-{S}tokes approximations in convex polyhedra}, Numer.
  Math., 131 (2015), pp.~771--822,
  \url{http://dx.doi.org/10.1007/s00211-015-0707-8}.

\bibitem{MR2121575}
{\sc V.~Girault, R.~H. Nochetto, and R.~Scott}, {\em Maximum-norm stability of
  the finite element {S}tokes projection}, J. Math. Pures Appl. (9), 84 (2005),
  pp.~279--330, \url{http://dx.doi.org/10.1016/j.matpur.2004.09.017}.

\bibitem{MR851383}
{\sc V.~Girault and P.-A. Raviart}, {\em Finite element methods for
  {N}avier-{S}tokes equations}, vol.~5 of Springer Series in Computational
  Mathematics, Springer-Verlag, Berlin, 1986,
  \url{http://dx.doi.org/10.1007/978-3-642-61623-5}.
\newblock Theory and algorithms.

\bibitem{MR0404849}
{\sc R.~B. Kellogg and J.~E. Osborn}, {\em A regularity result for the {S}tokes
  problem in a convex polygon}, J. Functional Analysis, 21 (1976),
  pp.~397--431, \url{http://dx.doi.org/10.1016/0022-1236(76)90035-5}.

\bibitem{MR1301452}
{\sc V.~A. Kozlov, V.~G. Maz'ya, and C.~Schwab}, {\em On singularities of
  solutions to the {D}irichlet problem of hydrodynamics near the vertex of a
  cone}, J. Reine Angew. Math., 456 (1994), pp.~65--97.

\bibitem{MR2536007}
{\sc A.~Kr\"{o}ner and B.~Vexler}, {\em A priori error estimates for elliptic
  optimal control problems with a bilinear state equation}, J. Comput. Appl.
  Math., 230 (2009), pp.~781--802,
  \url{http://dx.doi.org/10.1016/j.cam.2009.01.023}.

\bibitem{MR2680928}
{\sc K.~Kunisch, W.~Liu, Y.~Chang, N.~Yan, and R.~Li}, {\em Adaptive finite
  element approximation for a class of parameter estimation problems}, J.
  Comput. Math., 28 (2010), pp.~645--675,
  \url{http://dx.doi.org/10.4208/jcm.2009.10-m1003}.

\bibitem{MR2987056}
{\sc M.~Mitrea and M.~Wright}, {\em Boundary value problems for the {S}tokes
  system in arbitrary {L}ipschitz domains}, Ast\'{e}risque,  (2012),
  pp.~viii+241.

\bibitem{MR1259620}
{\sc J.~T. Oden, W.~Wu, and M.~Ainsworth}, {\em An a posteriori error estimate
  for finite element approximations of the {N}avier-{S}tokes equations},
  Comput. Methods Appl. Mech. Engrg., 111 (1994), pp.~185--202,
  \url{http://dx.doi.org/10.1016/0045-7825(94)90045-0}.

\bibitem{MR2583281}
{\sc F.~Tr\"{o}ltzsch}, {\em Optimal control of partial differential
  equations}, vol.~112 of Graduate Studies in Mathematics, American
  Mathematical Society, Providence, RI, 2010,
  \url{http://dx.doi.org/10.1090/gsm/112}.

\bibitem{MR4081184}
{\sc M.~Winkler}, {\em Error estimates for the finite element approximation of
  bilinear boundary control problems}, Comput. Optim. Appl., 76 (2020),
  pp.~155--199, \url{http://dx.doi.org/10.1007/s10589-020-00171-5}.

\end{thebibliography}
